\documentclass[11pt]{article}

% ============================================================
% Packages
% ============================================================
\usepackage[a4paper,margin=1in]{geometry}
\usepackage[T1]{fontenc}
\usepackage{lmodern}
\usepackage{microtype}
\usepackage{amsmath,amssymb,amsthm,mathtools}
\usepackage{bm}
\usepackage{graphicx}
\usepackage{booktabs}
\usepackage{enumitem}
\usepackage{tabularx}
\usepackage{array}
\usepackage{hyperref}
\usepackage[nameinlink,noabbrev]{cleveref}

\hypersetup{
  colorlinks=true,
  linkcolor=blue,
  citecolor=blue,
  urlcolor=blue
}

% ============================================================
% Macros / notation
% ============================================================
\newcommand{\dd}{\mathrm d}
\newcommand{\ii}{\mathrm i}
\newcommand{\ee}{\mathrm e}

\newcommand{\R}{\mathbb R}
\newcommand{\X}{\mathbf X}
\newcommand{\x}{\mathbf x}
\newcommand{\vct}{\mathbf v}
\newcommand{\nvec}{\mathbf n}
\newcommand{\grad}{\nabla}
\newcommand{\gradThree}{\nabla_3}
\newcommand{\gradFour}{\nabla_4}
\newcommand{\lapThree}{\nabla_3^2}
\newcommand{\lapFour}{\nabla_4^2}

\newcommand{\Wproj}{W}
\newcommand{\Zloc}{Z}
\newcommand{\Proj}[1]{\overline{#1}}

\newcommand{\kadd}{\kappa_{\mathrm{add}}}
\newcommand{\kPV}{\kappa_{\mathrm{PV}}}
\newcommand{\krho}{\kappa_{\rho}}
\newcommand{\betaPN}{\beta_{\mathrm{1PN}}}
\newcommand{\Kvec}{K_{\mathrm{vec}}}
\newcommand{\KEOS}{K_{\mathrm{EOS}}}
\newcommand{\alphaSq}{\alpha^2}

\newcommand{\PhiN}{\Phi}
\newcommand{\phibr}{\varphi}
\newcommand{\Sleak}{S_{\mathrm{leak}}}
\newcommand{\Cself}{C_{\mathrm{self}}}
\newcommand{\Cpara}{C_{\parallel}}
\newcommand{\CLong}{C_{L}}

\newcommand{\Exact}{\textbf{Exact}}
\newcommand{\Controlled}{\textbf{Controlled reduction}}
\newcommand{\Protocol}{\textbf{Protocol closure}}
\newcommand{\FullMatch}{\textbf{Full assembly}}
\newcommand{\Ansatz}{\textbf{Ansatz-level consistency}}

\newcommand{\claimclass}[1]{\paragraph{Claim class.}\emph{#1.}}
\newcommand{\TODO}[1]{\textbf{[TODO: #1]}}

\newcolumntype{L}{>{\raggedright\arraybackslash}X}

\newtheorem{claim}{Claim}
\newtheorem{remark}{Remark}

% ============================================================
% Title
% ============================================================
\title{A Controlled Derivation of the Full First Post-Newtonian Sector\\
from the Unified 4D Toy Model}
\author{Trevor Norris}
\date{\today}

\begin{document}
\maketitle

\begin{abstract}
% TODO: tighten after the main text is written.
We present a paper skeleton for deriving the full first post-Newtonian (1PN)
point-particle sector from the unified 4D toy model under explicit claim classes.
The narrative is organized so that each step is labeled as an exact identity,
controlled reduction, protocol-dependent closure, or full coefficient match.
The intended derivation chain is:
(i) the action-based 4D bulk/branes-observer framework and its exact projection identities,
(ii) the quasi-static Poisson hook on the brane,
(iii) a coherent-defect / small-body worldtube reduction to Newtonian particle motion,
(iv) scalar and optical dressing that fixes the self-sector coefficient,
(v) topology and internal-response closures that fix added-mass and breathing-response terms,
(vi) a constructive wake derivation of the EIH cross-velocity tensor,
and finally
(vii) assembly of the complete two-body 1PN Lagrangian and the standard perihelion result.
The paper makes explicit what is fixed within the stated closures
(\(q=1\), \(n=5\), \(\kadd=1/2\), \(\alphaSq=3/4\), \(a_H=0\), \(\Kvec=2/\pi^2\),
\(\kPV=3/2\), \(\betaPN=3\)) and what remains symbolic pending a more complete
moving-throat derivation.
\end{abstract}

% ============================================================
\section{Introduction and scope}
\label{sec:intro}

\subsection{Motivation: from coefficient fixing to full 1PN assembly}

The earlier papers in this series developed an effective, brane-facing description of defect dynamics at Newtonian and post-Newtonian order, while Ref.~\cite{norris2026_4d_model} supplied the unified \(4\)-spatial-dimensional generating model from which those effective sectors are meant to descend. The subsequent bridge paper then removed several phenomenological ``knobs'' by showing that key coefficients are fixed by topology, weak-field optics, wake matching, and wave-supported inertia. What has not yet appeared in one place is the full derivation narrative that starts from the 4D dictionary and ends with the standard conservative two-body first post-Newtonian (1PN) Einstein--Infeld--Hoffmann (EIH) Lagrangian. The purpose of the present paper is to provide that missing step.

By ``full 1PN'' we mean the complete conservative two-body ledger through order \(c^{-2}\): the quartic kinetic terms, the scalar self-sector terms proportional to \(\Phi v^2/c^2\), the velocity-dependent cross tensor proportional to \(\vct_A\!\cdot\!\vct_B\) and \((\vct_A\!\cdot\!\nvec)(\vct_B\!\cdot\!\nvec)\), and the static nonlinear term proportional to \(G^2/r^2\). Our goal is not to introduce another effective ansatz for those terms. Rather, it is to show that each coefficient in that ledger can be traced to a named ingredient of the 4D model: exact action-based identities, a controlled small-body reduction, one declared internal-response closure, and a constructive wake calculation.

This paper is therefore narrower than the unified 4D action paper and stronger than the earlier coefficient-by-coefficient bridge. It is narrower because we do not attempt to solve the fully dynamical moving-throat problem in the bulk. It is stronger because, within explicitly declared approximations, we now assemble the entire conservative two-body 1PN sector in one place and compare it directly to the standard EIH target. The claim is that the derivation chain is visible end to end:
\begin{enumerate}[leftmargin=1.5em]
  \item exact bulk dynamics and projection identities,
  \item a quasi-static near-zone Poisson hook on the brane,
  \item a coherent-defect / small-body reduction to Newtonian point-particle motion,
  \item scalar, optical, topological, and response corrections to the self sector,
  \item a constructive wake derivation of the cross-velocity tensor, and
  \item a final exact match to the standard two-body 1PN/EIH Lagrangian.
\end{enumerate}

A second purpose of the paper is bookkeeping clarity. In the 4D framework, brane observables are defined by a projection/measurement map rather than by stitched boundary conditions, so leakage, unresolved stress, and response terms are not optional decorations; they are the natural price of projection. For the 1PN problem this matters because it forces a clean separation between what is exact, what is reduced under an explicit regime assumption, and what depends on a declared response protocol. That separation is part of the result.

\subsection{What this paper claims}
\label{sec:claims}

The main claims are the following.
\begin{enumerate}[leftmargin=1.5em]
  \item The unified 4D model supplies the exact field-theoretic and projection framework needed to define a brane Newtonian potential: exact bulk balance laws, exact projected balances with explicit exchange terms, and a quasi-static near-zone Poisson hook.
  \item Under a coherent-defect / small-body closure---compact support, smooth external fields across the defect, negligible boundary fluxes at the retained order, and vanishing defect-centered dipole moment---the worldtube dynamics reduce to a Newtonian particle Lagrangian on the brane.
  \item Wave-supported rest energy reproduces the universal special-relativistic quartic kinetic coefficient, so the free-particle \(v^4\) term entering the 1PN bookkeeping is not inserted by hand.
  \item Scalar mass dressing together with the weak-field optical consistency condition fixes \(q=1\) and \(n=5\), and therefore fixes the self-sector \(\Phi v^2/c^2\) coefficient appearing in the pair Lagrangian.
  \item Local mass scaling and pair counting produce the static nonlinear \(G^2/r^2\) term, topology fixes \(\kadd=1/2\) for a \(w\)-uniform throat, and one explicit adiabatic breathing-response closure yields \(\kPV=3/2\), so that the familiar precession ledger closes to \(\betaPN=\krho+\kadd+\kPV=3\).
  \item A constructive isotropic Fourier-space wake basis yields the EIH cross coefficients \(\Cpara=-7/2\) and \(\CLong=-1/2\), with the unique minimal real solution \(\alphaSq=3/4\), \(a_H=0\), and \(\Kvec=2/\pi^2\).
  \item Combining these ingredients reproduces the standard conservative two-body 1PN/EIH Lagrangian exactly and yields the canonical perihelion coefficient in the test-mass orbit reduction.
\end{enumerate}

Stated differently, the point of the paper is not merely that the 4D model can be tuned to look 1PN-like. The point is that, once the derivation is organized sector by sector, the full conservative two-body 1PN structure can be recovered with a short list of named and auditable ingredients.

\subsection{What this paper does not claim}
\label{sec:nonclaims}

The scope should be kept equally explicit on the negative side.

First, we do \emph{not} claim an assumption-free theorem that begins with a fully solved moving-throat background in the 4D bulk and ends automatically with the point-particle EIH equations. The particle limit used here is a controlled reduction, not a numerically completed collective-coordinate derivation of the full defect profile.

Second, we do \emph{not} claim that every coefficient appearing in brane-effective dynamics is universally fixed once and for all. The coefficient \(\kPV\) is computed here only for one declared adiabatic one-degree-of-freedom closure. Outside that protocol, internal response coefficients should remain symbolic until the response problem is solved more completely in the same 4D theory.

Third, this paper is restricted to the conservative near-zone 1PN sector. It does not address radiation reaction, dissipative leakage, spin couplings, strong-field effects, or the full many-body problem beyond the conservative EIH structure. Nor does the present derivation by itself prove that all preferred-frame effects vanish in the completed brane-effective theory.

These non-claims are not weaknesses to be hidden. They are part of the correct interpretation of the result: the paper establishes a full 1PN derivation \emph{within a stated closure hierarchy}.

\subsection{Claim taxonomy}
\label{sec:taxonomy}

To keep the argument referee-safe, we will label derivation steps by claim class throughout.
\begin{description}[leftmargin=1.8em,style=nextline]
  \item[\Exact] identities that follow directly from the declared action, definitions, symmetry bookkeeping, or algebraic pair counting.
  \item[\Controlled reduction] steps that require an explicit regime assumption, such as quasi-staticity, small-body scaling, smooth external fields across the defect, harmonic far-field response, or low-frequency truncation.
  \item[\Protocol closure] coefficients that depend on a declared internal-response model rather than on kinematics alone. In the present paper the central example is the adiabatic one-degree-of-freedom closure used to compute \(\kPV\).
  \item[\Full assembly] the final exact algebraic equality between the derived reduced Lagrangian and the standard 1PN/EIH target once all previously derived ingredients are inserted.
\end{description}

This taxonomy is not cosmetic. It is the mechanism by which we separate exact 4D statements from the additional assumptions needed to turn them into a brane-facing point-particle theory.

\subsection{Roadmap}
\label{sec:roadmap}

Section~\ref{sec:exec-summary} gives the executive ledger of the coefficients and matches to be established. Section~\ref{sec:dictionary} restates the minimal 4D dictionary carried forward from the unified action paper: the action-level fields, projection map, exact balance laws, and the near-zone Poisson hook. Section~\ref{sec:newtonian-limit} then performs the coherent-defect / small-body reduction that turns the worldtube force law into a Newtonian particle Lagrangian.

The remaining sections build the 1PN ledger one sector at a time. Sections~\ref{sec:self-sector} and \ref{sec:static-term} derive the free-particle quartic term, the scalar/optical self-sector contribution, and the static nonlinear term. Section~\ref{sec:ledger} collects the topology and internal-response pieces, including \(\kadd\), \(\kPV\), and \(\betaPN\). Section~\ref{sec:wake} presents the constructive wake calculation for the cross-velocity tensor. Section~\ref{sec:full-eih} assembles all ingredients into the full two-body 1PN Lagrangian and compares it directly to the standard EIH form. Section~\ref{sec:orbit} reduces the result to the test-mass orbit problem and recovers the canonical perihelion coefficient. Section~\ref{sec:ledger-fixed-vs-symbolic} records which quantities are now fixed and which remain symbolic. Section~\ref{sec:verification} closes with the verification and reproducibility ledger tied to the symbolic harnesses.

% ============================================================
\section{Executive summary of results}
\label{sec:exec-summary}

This paper is the point at which the \(4\)D program becomes a full first post-Newtonian derivation rather than a coefficient-fixing story. Ref.~\cite{norris2026_4d_model} supplies the exact bulk action, projection dictionary, exact balance laws, and the near-zone Poisson hook needed to define a brane-facing Newtonian sector. The present paper then carries those ingredients through a controlled coherent-defect / small-body reduction, the scalar and optical self-sector bookkeeping, topology and internal-response corrections, and a constructive wake calculation for the velocity-dependent pair terms. The final output is an explicit conservative two-body \(1\)PN Lagrangian whose coefficients are traced to named ingredients of the \(4\)D model rather than inserted phenomenologically.

\paragraph{Notation warning (two different \(K\)'s).}
Throughout this paper we distinguish the equation-of-state constant \(\KEOS\) in
\[
P(\rho)=\KEOS \rho^n
\]
from the dimensionless vector-sector normalization \(\Kvec\) that appears in the wake/overlap tensor. Only \(n\) and \(\Kvec\) are fixed by the present derivation chain; \(\KEOS\) remains a dimensional model parameter set by units and normalization.

\paragraph{Headline result.}
The main claim of the paper is that, within the declared closure hierarchy, the toy model reproduces the full conservative two-body Einstein--Infeld--Hoffmann ledger through order \(c^{-2}\). The required coefficients are not fixed by a single argument but by a triangulated chain:
Newtonian matching fixes \(q=1\), weak-field optics fixes \(n=5\), wave-supported rest energy gives the universal quartic kinetic coefficient, topology fixes \(\kadd\), one explicit adiabatic response closure fixes \(\kPV\), and the wake overlap tensor fixes the velocity-dependent cross terms. The resulting reduced Lagrangian then agrees identically with the standard two-body EIH form.

Table~\ref{tab:headline-results} lists the quantities proved in the paper and the logic by which each is obtained.

\begin{table}[t]
  \centering
  \caption{Headline results established in this paper. The quartic kinetic coefficient is listed in Lagrangian normalization; the equivalent energy-expansion coefficient is \(3/8\).}
  \label{tab:headline-results}
  \small
  \setlength{\tabcolsep}{4pt}
  \renewcommand{\arraystretch}{1.18}
  \begin{tabularx}{\linewidth}{@{}L c L L@{}}
    \toprule
    Quantity & Value & Source of derivation & Claim class \\
    \midrule
    Newtonian scalar coupling & \(q=1\) & Newtonian matching of the scalar-dressed particle action & \Controlled \\
    EOS exponent & \(n=5\) & weak-field optical/refraction coefficient & \Controlled \\
    Free-particle quartic coefficient & \(1/8\) in \(L\) & wave-supported rest energy / SR expansion & \Controlled \\
    Self-sector pair coefficient & \(+3/2\) in the EIH pair term & scalar + optical dressing with \(q=1\), \(n=5\) & \Controlled \\
    Static 1PN coefficient & \(-1/2\) & local mass scaling + pair counting & \Exact / \Controlled \\
    Added-mass coefficient & \(\kadd=1/2\) & topology of a \(w\)-uniform throat slice & \Controlled \\
    Counterfactual bubble check & \(1/3\) & 4D bubble added-mass comparison & \Controlled \\
    Wake longitudinal fraction & \(\alphaSq=3/4\) & thermodynamic closure + wake matching & \Controlled \\
    Helical parity-even amplitude & \(a_H=0\) & minimal real wake solution & \Controlled \\
    Vector-sector normalization & \(\Kvec=2/\pi^2\) & EIH cross-term match & \Controlled \\
    Internal-response coefficient & \(\kPV=3/2\) & explicit adiabatic 1-DOF closure & \Protocol \\
    Precession ledger & \(\betaPN=3\) & \(\krho+\kadd+\kPV\) & \Protocol \\
    Cross coefficients & \(\Cpara=-7/2\), \(\CLong=-1/2\) & constructive wake overlap tensor & \Controlled \\
    Two-body 1PN Lagrangian & exact EIH match & full assembly of all derived ingredients & \FullMatch \\
    Perihelion coefficient & GR value & test-mass orbit reduction from \(\betaPN=3\) & \FullMatch \\
    \bottomrule
  \end{tabularx}
\end{table}

\paragraph{How the ledger closes.}
The scalar/self sector and the vector/cross sector are fixed by different physical inputs and only meet at the very end. On the scalar side, Newtonian reduction forces the mass-dressing exponent \(q=1\), weak-field optics forces the stiff polytropic exponent \(n=5\), and those two facts together fix the pair self coefficient multiplying \(v_A^2+v_B^2\). The universal free-particle quartic term follows from treating rest energy as wave-supported energy, so the \(v^4\) coefficient is not inserted by hand. The remaining static nonlinear term proportional to \(G^2/r^2\) then follows from local mass scaling together with the pair bookkeeping of the reduced two-body action.

On the response side, the \(w\)-uniform throat geometry fixes the added-mass contribution to the inertial ledger at \(\kadd=1/2\), while a counterfactual \(4\)D bubble would instead give \(1/3\). This is useful because it turns added mass into a topology discriminator rather than a tunable inertial correction. The pressure-volume / breathing contribution is more subtle: in general \(\kPV\) is a response coefficient and should remain symbolic, but for the explicit adiabatic one-degree-of-freedom closure adopted here it evaluates to
\[
\kPV=\frac32,
\qquad
\betaPN=\krho+\kadd+\kPV=3.
\]
That value is then the quantity that controls the final orbit-level perihelion ledger.

The velocity-dependent cross sector is fixed independently by a constructive isotropic Fourier-space wake basis. Matching the EIH cross tensor first yields a one-parameter family; the thermodynamic identity implied by the already-fixed \(n=5\) then collapses that family to the unique minimal real point
\[
\alphaSq=\frac34,
\qquad
a_H=0,
\qquad
\Kvec=\frac{2}{\pi^2},
\]
which in turn reproduces the standard EIH cross coefficients
\[
\Cpara=-\frac72,
\qquad
\CLong=-\frac12.
\]

\paragraph{Final assembled result.}
Once all sectors are inserted, the derived reduced two-body Lagrangian takes the form
\begin{align}
L_{\rm 1PN}^{\rm derived}
={}&
\frac12 m_A v_A^2 + \frac12 m_B v_B^2
+ \frac{m_A v_A^4 + m_B v_B^4}{8c^2}
+ \frac{Gm_A m_B}{r_{AB}}
\nonumber\\
&
+ \frac{Gm_A m_B}{c^2 r_{AB}}
\left[
\frac32\bigl(v_A^2+v_B^2\bigr)
-\frac72\,\vct_A\!\cdot\!\vct_B
-\frac12\,(\vct_A\!\cdot\!\nvec)(\vct_B\!\cdot\!\nvec)
\right]
- \frac{G^2m_A m_B(m_A+m_B)}{2c^2 r_{AB}^2}.
\label{eq:exec-eih-summary}
\end{align}
The claim of the paper is that this is not merely EIH-like: after the previous sectors are derived, it is algebraically identical to the standard conservative two-body EIH target. In the test-mass orbit reduction, the same ledger yields the standard perihelion coefficient because \(\betaPN=3\).

\paragraph{How to read the claim classes.}
Most rows in Table~\ref{tab:headline-results} are \Controlled statements: they require explicit regime assumptions such as quasi-staticity, coherent-defect / small-body reduction, or linear far-field wake closure, but are then derived without further tuning. The single genuinely \Protocol result is \(\kPV\), because it depends on the declared internal response model. The final equality to the two-body EIH Lagrangian and the perihelion result are \FullMatch statements: once the earlier ingredients are imported, the remaining algebra closes exactly.

\paragraph{What is fixed and what remains symbolic.}
Within the scope of this paper, the quantities
\[
q=1,\qquad n=5,\qquad \kadd=\frac12,\qquad
\alphaSq=\frac34,\qquad a_H=0,\qquad \Kvec=\frac{2}{\pi^2},
\qquad \kPV=\frac32\ \text{(for the adiabatic closure)},\qquad
\betaPN=3
\]
should be regarded as derived carry-forward values rather than knobs. What remains symbolic is equally important: the absolute EOS scale \(\KEOS\), localization-profile data, higher-order open-system response data, and any non-adiabatic generalization of \(\kPV\). The paper therefore claims a full \(1\)PN derivation within a stated closure hierarchy, not an assumption-free theorem of the fully solved moving-throat bulk dynamics.

% ============================================================
\section{Minimal 4D dictionary carried into the 1PN derivation}
\label{sec:dictionary}

This section is intentionally downstream of the unified 4D action paper. We do not
repeat the full parent-theory derivation here. Instead, we record only the objects and
identities that the present 1PN calculation actually uses. The purpose of the section is
to freeze notation, to keep the distinction between exact identities and controlled
reductions visible, and to make clear what data are imported from the parent 4D model
before any post-Newtonian bookkeeping begins.

The two structural principles carried forward are simple.

First, the fundamental dynamics live in four spatial dimensions,
\[
\mathbf{X}=(x,y,z,w)\in\mathbb R^4,
\]
with time added in the usual way. What a brane observer measures is \emph{not} imposed by
boundary matching; it is defined operationally by a projection map localized near
\(w=0\).

Second, the appearance of familiar \(3\)-dimensional operators on the brane is not taken
as an axiom. It arises either as an exact consequence of the projection map
(open-system continuity and the longitudinal identity) or as a controlled reduction
(quasi-static Poisson behavior, zero-mode electromagnetic localization, small-body
worldtube limits). That exact-versus-reduced distinction is the part of the dictionary
that matters most for the present paper.

\subsection{Bulk variables, projection, and reduction}
\label{sec:dictionary-bulk}

We work on a \((4+1)\)-dimensional spacetime with coordinates
\begin{equation}
X^M=(t,x,y,z,w)\equiv (t,\x,w),
\qquad
\X=(x,y,z,w)\in\R^4,
\qquad
\x=(x,y,z)\in\R^3.
\end{equation}
The parent theory contains a localized matter sector, an optional localized gauge sector,
and a dynamical throat geometry. For the present paper it is enough to keep the following
objects explicit:
\begin{equation}
\psi(\X,t),
\qquad
\rho(\X,t)=|\psi(\X,t)|^2,
\qquad
A_M(\X,t),
\qquad
\Zloc(w),
\qquad
a(t),\ L(t).
\end{equation}
Here \(\psi\) is the bulk matter field, \(\rho\) its density, \(A_M\) the optional gauge
potential, \(\Zloc(w)\) the localization profile multiplying the bulk Maxwell kinetic
term, and \((a,L)\) the effective throat radius and axial extent that parameterize the
geometry sector.

A key distinction inherited from the 4D paper is the difference between
\emph{projection} and \emph{reduction}.

\paragraph{Projection.}
A brane observer measures weighted averages across the extra coordinate using a fixed,
nonnegative kernel \(\Wproj(w)\) normalized by
\begin{equation}
\int_{-\infty}^{\infty}\Wproj(w)\,\dd w = 1.
\label{eq:dictionary-Wnorm}
\end{equation}
For any bulk scalar \(Q(t,\x,w)\), the projected brane observable is
\begin{equation}
\Proj{Q}(t,\x)
\equiv
\int_{-\infty}^{\infty}\Wproj(w)\,Q(t,\x,w)\,\dd w.
\label{eq:dictionary-projection}
\end{equation}
In particular, the projected density and projected brane flux are
\begin{equation}
\Proj{\rho}(t,\x)
=
\int_{-\infty}^{\infty}\Wproj(w)\,\rho(t,\x,w)\,\dd w,
\qquad
\mathbf{J}_{\rm br}(t,\x)
=
\int_{-\infty}^{\infty}\Wproj(w)\,\mathbf{j}_{xyz}(t,\x,w)\,\dd w,
\label{eq:dictionary-brane-observables}
\end{equation}
where \(\mathbf{j}_{xyz}\) denotes the \((x,y,z)\) components of the bulk number current.

\paragraph{Reduction.}
Separately, one may integrate over \(w\) under an explicit ansatz in order to obtain a
brane-effective theory with renormalized couplings. This is not the same operation as the
measurement map \eqref{eq:dictionary-projection}. Projection defines what the brane
observer sees; reduction produces effective equations only after additional assumptions
have been stated. Keeping those two operations distinct is essential throughout the 1PN
derivation.

Where \(\Proj{\rho}>0\), the projected brane velocity is defined kinematically by
\begin{equation}
\mathbf{v}_{\rm br}(t,\x)\equiv \frac{\mathbf{J}_{\rm br}(t,\x)}{\Proj{\rho}(t,\x)}.
\label{eq:dictionary-vbrane}
\end{equation}
Later sections will reduce this open-system brane kinematics to an effective Newtonian
particle sector under explicit near-zone and small-body assumptions.

\subsection{Minimal action-side ingredients}
\label{sec:dictionary-action}

The present paper does not need the full parent action written out term by term, but it
does need the exact identities obtained from it. The relevant action-side facts are the
following.

\paragraph{Matter sector.}
The bulk matter field defines an exact number density \(\rho=|\psi|^2\) and bulk number
current \(j^i\), with
\begin{equation}
j^i = \rho\,v^i,
\label{eq:dictionary-bulk-velocity}
\end{equation}
which introduces the bulk hydrodynamic velocity \(v^i\) wherever \(\rho>0\).

\paragraph{Equation of state.}
The barotropic matter sector is parameterized by
\begin{equation}
P(\rho)=\KEOS \rho^n.
\label{eq:dictionary-eos}
\end{equation}
For the purposes of this section, \(n\) is left symbolic. Later sections will fix its
value from the optical weak-field consistency condition.

\paragraph{Localized gauge sector.}
When retained, the bulk Maxwell operator is weighted by \(\Zloc(w)\), so that variation of
the parent action yields an exact localized Maxwell equation of the form
\begin{equation}
\partial_M\!\left(\Zloc(w)F^{MN}\right)
+\frac{1}{\xi}\,\partial^N(\partial\!\cdot\!A)
=
\mu_0\bigl(J_\psi^N+J_{\rm ext}^N\bigr).
\label{eq:dictionary-maxwell}
\end{equation}
This is included here not because we will solve the full bulk gauge problem in the
Newtonian sector, but because it makes clear that any brane-effective transverse/vector
sector is grounded in an exact action-based channel rather than in a purely kinematic
ansatz.

\paragraph{Geometry sector.}
The localized throat is represented by two collective geometric variables,
\begin{equation}
a(t)\quad\text{(effective radius)},
\qquad
L(t)\quad\text{(effective axial extent)}.
\label{eq:dictionary-geometry-dofs}
\end{equation}
At the level of the parent theory, these variables enter the confinement geometry and can
exchange energy with the fields. In the present paper we only need the structural fact
that such internal degrees of freedom exist and can support a controlled breathing
response. Their explicit adiabatic closure is introduced later, when \(\kPV\) is
computed.

\paragraph{Projection-induced stress.}
Because projection does not commute with nonlinear products, the projected momentum sector
is not generically closed. A convenient structural measure of that non-commutation is the
projected extra-stress tensor
\begin{equation}
R_{ab}
\equiv
\Proj{\frac{j_a j_b}{\rho}}
-
\Proj{\rho}\,v_{{\rm br},a}v_{{\rm br},b},
\qquad a,b\in\{x,y,z\}.
\label{eq:dictionary-extra-stress}
\end{equation}
The present 1PN derivation will not need the full projected momentum equation in the main
text, but it \emph{does} rely on the fact that leakage and projection-induced stresses are
real structural channels of the parent theory. The controlled reductions used below amount
to saying when those channels are subleading.

\subsection{Exact projected continuity and leakage}
\label{sec:dictionary-continuity}

From the parent matter dynamics one obtains the exact bulk continuity identity
\begin{equation}
\partial_t \rho + \partial_i j^i = 0,
\qquad i\in\{x,y,z,w\}.
\label{eq:dictionary-bulk-continuity}
\end{equation}
Projecting this equation with the kernel \(\Wproj(w)\) and separating the \(w\)-flux gives
an exact open-system brane continuity equation:
\begin{equation}
\partial_t \Proj{\rho} + \gradThree\!\cdot \mathbf{J}_{\rm br}
=
\Sleak,
\label{eq:dictionary-proj-continuity}
\end{equation}
with leakage source
\begin{equation}
\Sleak(t,\x)
\equiv
-\left[\Wproj(w)\,j^w(t,\x,w)\right]_{-\infty}^{+\infty}
+
\int_{-\infty}^{\infty}\Wproj'(w)\,j^w(t,\x,w)\,\dd w.
\label{eq:dictionary-leakage}
\end{equation}
Under the common fast-decay assumption
\begin{equation}
\Wproj(w)\,j^w(t,\x,w)\to 0
\qquad\text{as }|w|\to\infty,
\label{eq:dictionary-fast-decay}
\end{equation}
this simplifies to
\begin{equation}
\Sleak(t,\x)
=
\int_{-\infty}^{\infty}\Wproj'(w)\,j^w(t,\x,w)\,\dd w.
\label{eq:dictionary-leakage-simplified}
\end{equation}

The point of \eqref{eq:dictionary-proj-continuity} is conceptual as much as algebraic.
Even when the bulk matter sector is exactly conservative, the projected brane density is
generically an open-system observable. Leakage is therefore not an optional phenomenology
term added later; it is forced by the projection definition itself.

\subsection{The exact longitudinal identity and the Poisson hook}
\label{sec:poisson-hook}

The 1PN derivation does not begin from Newton's law as a postulate. It begins from the
exact projected continuity equation together with a Helmholtz decomposition of the brane
velocity. Write
\begin{equation}
\mathbf{v}_{\rm br}(t,\x)
=
\gradThree \phibr(t,\x) + \mathbf{v}_T(t,\x),
\qquad
\gradThree\!\cdot\mathbf{v}_T = 0,
\label{eq:dictionary-helmholtz}
\end{equation}
where \(\phibr\) is the longitudinal brane velocity potential and \(\mathbf{v}_T\) is the
divergence-free transverse part.

Using \(\mathbf{J}_{\rm br}=\Proj{\rho}\,\mathbf{v}_{\rm br}\) in
\eqref{eq:dictionary-proj-continuity} and expanding the brane divergence gives the exact
longitudinal identity
\begin{equation}
\boxed{
\Proj{\rho}\,\lapThree \phibr
=
\Sleak
-
\partial_t\Proj{\rho}
-
\bigl(\gradThree \Proj{\rho}\bigr)\!\cdot\!\bigl(\gradThree\phibr+\mathbf{v}_T\bigr).
}
\label{eq:dictionary-longitudinal-identity}
\end{equation}
This equation is exact once the projection map and Helmholtz decomposition have been
declared. It is not an additional constitutive assumption.

A Poisson-type equation on the brane appears only in a controlled near-zone regime. If
\begin{enumerate}[leftmargin=1.5em]
  \item \(\Proj{\rho}\approx \rho_0\) is slowly varying over the region of interest,
  \item \(\partial_t\Proj{\rho}\) is quasi-static at the order retained,
  \item the advective correction \((\gradThree \Proj{\rho})\!\cdot\!(\gradThree\phibr+\mathbf{v}_T)\) is subleading, and
  \item the transverse sector is negligible or perturbative,
\end{enumerate}
then \eqref{eq:dictionary-longitudinal-identity} reduces to the Poisson hook
\begin{equation}
\lapThree \phibr
\;\simeq\;
\frac{S_{\rm eff}}{\rho_0},
\qquad
S_{\rm eff}\equiv \Sleak\ \text{plus any retained slow-source terms}.
\label{eq:dictionary-poisson-hook}
\end{equation}
For a localized effective source,
\begin{equation}
S_{\rm eff}(\x)\sim \mathcal S\,\delta^{(3)}(\x),
\end{equation}
the Green's function of the \(3\)-dimensional Laplacian gives
\begin{equation}
\phibr(\x)\sim -\frac{\mathcal S}{4\pi\rho_0}\,\frac{1}{r},
\qquad
\gradThree\phibr \sim \frac{1}{r^2},
\qquad
r=|\x|.
\label{eq:dictionary-inverse-square}
\end{equation}
This is the route by which inverse-square longitudinal behavior enters the present
program. In the next section, a controlled worldtube reduction turns that kinematic
Poisson-sector structure into the effective Newtonian potential \(\PhiN\) used in the
particle Lagrangian.

It is important not to overstate the claim at this stage. Equation
\eqref{eq:dictionary-longitudinal-identity} is exact. Equation
\eqref{eq:dictionary-poisson-hook} is a regime statement. The identification of the
resulting near-zone scalar with the Newtonian potential used by a reduced point particle
is a further controlled constitutive step carried out later in the paper.

\subsection{Controlled EM reduction and standard brane operators}
\label{sec:dictionary-em}

Although the scalar Newtonian sector is the main focus of the next section, the parent 4D
paper also provides a second example of how familiar brane operators emerge as controlled
reductions rather than axioms. Under the zero-mode assumptions
\begin{equation}
A_w = 0,
\qquad
\partial_w A_\mu = 0,
\qquad
\mu\in\{0,1,2,3\},
\label{eq:dictionary-zero-mode}
\end{equation}
and with a brane-compatible gauge choice, integrating
\eqref{eq:dictionary-maxwell} over \(w\) yields the effective brane Maxwell equation
\begin{equation}
\partial_\mu F^{\mu\nu}
=
\mu_0^{\rm eff}\,J_{\rm eff}^\nu,
\qquad
\mu_0^{\rm eff}
=
\frac{\mu_0}{Z_{\rm int}},
\qquad
Z_{\rm int}\equiv \int_{-\infty}^{\infty}\Zloc(w)\,\dd w.
\label{eq:dictionary-brane-maxwell}
\end{equation}
We do not use \eqref{eq:dictionary-brane-maxwell} as an engine for the Newtonian scalar
derivation itself. Its role here is organizational: it reminds the reader that the same
exact-versus-controlled logic used for the Poisson hook also governs the transverse/vector
sector of the 4D model.

\subsection{What is actually carried forward into the 1PN paper}
\label{sec:dictionary-note}

Only a small subset of the parent theory is needed below.

What we carry forward explicitly is:
\begin{enumerate}[leftmargin=1.5em]
  \item the bulk fields \((\psi,A_M)\), the density \(\rho\), and the geometric variables \((a,L)\);
  \item the projection map \(\Proj{\cdot}\) and the distinction between projection and reduction;
  \item the exact open-system projected continuity equation with leakage;
  \item the exact longitudinal identity and its quasi-static Poisson hook;
  \item the structural existence of projection-induced stress and internal geometric response sectors.
\end{enumerate}

What we do \emph{not} repeat in full is the detailed parent-action variation, the full
projected momentum algebra, or the full higher-mode electromagnetic analysis. Those
results belong to the unified 4D paper and are only invoked here through the minimal
dictionary above.

\claimclass{\Exact for the projection map, the bulk continuity law, the leakage identity,
the extra-stress definition, and the longitudinal identity. \Controlled for the
quasi-static Poisson hook \eqref{eq:dictionary-poisson-hook} and the zero-mode
electromagnetic reduction \eqref{eq:dictionary-brane-maxwell}.}


% ============================================================
\section{From the Poisson hook to the Newtonian point-particle limit}
\label{sec:newtonian-limit}

This section is the first controlled bridge from the exact projection-language of the
parent \(4\)D model to the effective particle language used in the rest of the paper.
The point is not to claim an assumption-free theorem of the fully solved moving-throat
bulk dynamics. The point is narrower and more useful: once the exact longitudinal
identity has entered a quasi-static Poisson regime, a compact coherent defect admits a
controlled monopole limit in which its center-of-mass motion is governed by the usual
Newtonian point-particle Lagrangian.

The logic has two steps. First, the Poisson hook of
\cref{eq:dictionary-poisson-hook} tells us that the longitudinal brane scalar behaves,
in the near zone, like a \(3\)-dimensional potential sourced by a localized effective
monopole. Second, a small-body worldtube reduction shows that a compact defect in such a
smooth external field responds at leading order only through its total mass monopole,
with finite-size corrections suppressed by powers of \(a/r\), where \(a\) is the defect
size and \(r\) the external variation scale.

Operationally, we therefore define the measured Newtonian potential \(\PhiN\) as the
properly normalized near-zone scalar channel on the brane,
\begin{equation}
\PhiN(\x,t)\equiv \lambda_\Phi\,\phibr(\x,t),
\label{eq:newtonian-potential-definition}
\end{equation}
with \(\lambda_\Phi\) fixed by monopole calibration. For an isolated compact source \(B\),
the near-zone solution outside the source support is then
\begin{equation}
\Phi_B(\x,t)
=
-\frac{Gm_B}{|\x-\mathbf X_B(t)|}
+ O\!\left(\frac{Gm_B a_B^2}{|\x-\mathbf X_B|^3}\right),
\label{eq:isolated-source-newtonian-potential}
\end{equation}
where \(G\) is the measured brane normalization of the scalar channel and \(a_B\) is the
source size. The remaining task in this section is to show that a second compact defect
moving in this field reduces, at leading order, to a Newtonian point particle.

\subsection{Controlled near-zone assumptions}
\label{sec:near-zone-assumptions}

The reduction carried out below uses the following explicit assumptions.
\begin{enumerate}[leftmargin=1.5em]
  \item \textbf{Compact support / small body.}
  The defect occupies a region of characteristic size \(a\) that is small compared with
  the inter-body separation and with the scale over which the external field varies:
  \(a\ll r\).

  \item \textbf{Smooth external field across the support.}
  Over the defect worldtube, the external potential may be Taylor-expanded about the
  center-of-mass position.

  \item \textbf{Negligible worldtube boundary fluxes at the retained order.}
  Any residual leakage, entrainment, or stress flux through the defect boundary is beyond
  the Newtonian order kept here.

  \item \textbf{Defect-centered coordinates.}
  The spatial origin is chosen so that the defect dipole moment vanishes.

  \item \textbf{Quasi-static near zone.}
  The dropped terms in the exact longitudinal identity remain parametrically small, so
  the Poisson hook of \cref{eq:dictionary-poisson-hook} is valid.

  \item \textbf{Conservative sector only.}
  Radiation, dissipative exchange, and higher-response effects are deferred to later
  work and are not part of the Newtonian limit derived in this section.
\end{enumerate}

Under these assumptions, the Poisson-hook scalar can be treated as an ordinary Newtonian
external potential over the support of a compact defect, while the defect's internal
structure enters only through higher multipoles.

\subsection{Coherent-defect worldtube closure}
\label{sec:coherent-defect}

Let the defect center be \(\mathbf X_{\rm cm}(t)\), and introduce defect-centered
coordinates
\begin{equation}
\boldsymbol{\xi}=\x-\mathbf X_{\rm cm}(t).
\end{equation}
Write the defect density in the form
\begin{equation}
\rho_{\rm def}(\x,t)=\rho_*(\boldsymbol{\xi},t),
\end{equation}
where \(\rho_*\) is localized in a worldtube \(W_t\) around \(\boldsymbol{\xi}=0\).
Nothing in the leading monopole reduction depends on the profile being exactly Gaussian,
but it is convenient to use an isotropic Gaussian as an explicit representative closure,
\begin{equation}
\rho_{\rm def}(\boldsymbol{\xi})
=
\frac{M}{\pi^{3/2}a^3}
\exp\!\left(-\frac{\xi^2+\eta^2+\zeta^2}{a^2}\right),
\label{eq:defect-gaussian-profile}
\end{equation}
because its moments are elementary and match the symbolic harness exactly.

The coherent-defect closure is then encoded in the lowest moments:
\begin{align}
\int_{W_t}\rho_{\rm def}\,\dd^3\xi &= M,
\label{eq:defect-mass-moment}
\\[3pt]
\int_{W_t}\rho_{\rm def}\,\xi_i\,\dd^3\xi &= 0,
\label{eq:defect-dipole-zero}
\\[3pt]
\int_{W_t}\rho_{\rm def}\,\xi_i\xi_j\,\dd^3\xi &= Q_{ij}.
\label{eq:defect-quadrupole-def}
\end{align}
For the isotropic Gaussian \eqref{eq:defect-gaussian-profile},
\begin{equation}
Q_{ij}=\frac{Ma^2}{2}\,\delta_{ij}.
\label{eq:defect-isotropic-second-moment}
\end{equation}

Equation \eqref{eq:defect-dipole-zero} is the critical centering condition: once the
dipole vanishes, the first finite-size correction to a smooth external force is
quadrupolar, not dipolar. That is why the monopole force law is the correct leading
small-body limit.

\subsection{Momentum and kinetic-energy split}
\label{sec:cm-split}

We now decompose the defect velocity into coherent translation plus internal motion:
\begin{equation}
\vct(\boldsymbol{\xi},t)
=
\dot{\mathbf X}_{\rm cm}(t)
+
\mathbf u_{\rm int}(\boldsymbol{\xi},t).
\label{eq:defect-velocity-split}
\end{equation}
The coherent-translation closure is the statement that the internal flow has zero
mass-weighted mean,
\begin{equation}
\int_{W_t}\rho_{\rm def}\,\mathbf u_{\rm int}\,\dd^3\xi = 0.
\label{eq:coherent-translation-closure}
\end{equation}
Under this condition, the total momentum is
\begin{equation}
\mathbf P
=
\int_{W_t}\rho_{\rm def}\,\vct\,\dd^3\xi
=
M\dot{\mathbf X}_{\rm cm}.
\label{eq:defect-total-momentum}
\end{equation}
Likewise, the total kinetic energy separates into center-of-mass motion plus internal
energy:
\begin{align}
T
&=
\frac12\int_{W_t}\rho_{\rm def}\,|\vct|^2\,\dd^3\xi
\nonumber\\
&=
\frac12 M\dot{\mathbf X}_{\rm cm}^{\,2}
+
\frac12\int_{W_t}\rho_{\rm def}\,|\mathbf u_{\rm int}|^2\,\dd^3\xi.
\label{eq:defect-kinetic-split}
\end{align}
At Newtonian order, the second term is an internal energy contribution. It can renormalize
the defect rest energy, but it does not alter the leading external center-of-mass
kinematics so long as the coherent-translation closure is maintained.

This is the first place where the reduced description already displays an important
feature that later becomes the Newtonian matching condition of the scalar sector:
the same total monopole \(M\) appears in the kinetic term and in the leading coupling to
the external potential below. In the reduced Newtonian description, inertial and
gravitational mass therefore coincide automatically.

\subsection{Newtonian monopole force law}
\label{sec:monopole-force}

Let the defect move in a smooth external potential \(\PhiN(\x,t)\). The worldtube force is
\begin{equation}
\mathbf F_{\rm ext}(t)
=
-\int_{W_t}\rho_{\rm def}(\boldsymbol{\xi},t)\,
\gradThree \PhiN\!\bigl(\mathbf X_{\rm cm}(t)+\boldsymbol{\xi},t\bigr)\,\dd^3\xi
+
\mathbf F_{\partial W},
\label{eq:worldtube-force}
\end{equation}
where \(\mathbf F_{\partial W}\) denotes boundary-flux terms. By assumption,
\(\mathbf F_{\partial W}\) is beyond the retained Newtonian order.

Taylor-expand the external field about the center:
\begin{align}
\partial_i \PhiN(\mathbf X_{\rm cm}+\boldsymbol{\xi},t)
&=
\partial_i\PhiN(\mathbf X_{\rm cm},t)
+
\xi_j\,\partial_j\partial_i\PhiN(\mathbf X_{\rm cm},t)
\nonumber\\
&\quad
+
\frac12\xi_j\xi_k\,\partial_j\partial_k\partial_i\PhiN(\mathbf X_{\rm cm},t)
+\cdots.
\label{eq:potential-taylor-expansion}
\end{align}
Using \eqref{eq:defect-mass-moment}--\eqref{eq:defect-quadrupole-def}, the force becomes
\begin{equation}
F_i
=
- M\,\partial_i\PhiN(\mathbf X_{\rm cm},t)
-\frac12\,Q_{jk}\,\partial_j\partial_k\partial_i\PhiN(\mathbf X_{\rm cm},t)
+\cdots.
\label{eq:force-multipole-expansion}
\end{equation}
The dipole term drops out identically because of
\eqref{eq:defect-dipole-zero}. Therefore the leading force is monopolar:
\begin{equation}
\mathbf F_{\rm ext}
=
- M\,\gradThree \PhiN(\mathbf X_{\rm cm},t)
+
O\!\left(Q\,\nabla^3\PhiN\right).
\label{eq:monopole-force-leading}
\end{equation}

For an external field sourced by a distant compact body,
\(\PhiN\sim Gm/r\), the monopole force scales as \(MGm/r^2\), whereas the leading
finite-size correction scales as \(MGm\,a^2/r^4\). Their ratio is therefore
\begin{equation}
\frac{|\mathbf F_{\rm quad}|}{|\mathbf F_{\rm mono}|}
=
O\!\left(\frac{a^2}{r^2}\right).
\label{eq:quadrupole-suppression}
\end{equation}
This is the explicit small-body control parameter of the Newtonian reduction.

Combining \eqref{eq:defect-total-momentum} with \eqref{eq:monopole-force-leading}, we
obtain the reduced center-of-mass equation of motion
\begin{equation}
M\ddot{\mathbf X}_{\rm cm}
=
- M\,\gradThree\PhiN(\mathbf X_{\rm cm},t)
+
O\!\left(M\,\frac{a^2}{r^2}\,\frac{\PhiN}{r}\right).
\label{eq:newtonian-center-of-mass-eom}
\end{equation}
At leading order, this is equivalent to the Euler--Lagrange equation of the Newtonian
point-particle Lagrangian
\begin{equation}
L_{\rm N}
=
\frac12 M\dot{\mathbf X}_{\rm cm}^{\,2}
-
M\PhiN(\mathbf X_{\rm cm},t).
\label{eq:newtonian-particle-lagrangian}
\end{equation}
Once the reduction has been made, the passage from
\eqref{eq:newtonian-particle-lagrangian} to
\eqref{eq:newtonian-center-of-mass-eom} is exact.

\paragraph{Two-body pair counting.}
For two compact defects \(A\) and \(B\), with
\begin{equation}
r_{AB}=|\mathbf X_A-\mathbf X_B|,
\end{equation}
the near-zone field of \(B\) evaluated at \(A\) is
\begin{equation}
\Phi_B(\mathbf X_A)=-\frac{Gm_B}{r_{AB}},
\qquad
\Phi_A(\mathbf X_B)=-\frac{Gm_A}{r_{AB}}.
\end{equation}
Summing the one-body kinetic pieces and pair-counting the interaction once gives the
standard Newtonian two-body Lagrangian
\begin{equation}
L_{\rm N}^{(2)}
=
\frac12 m_A v_A^2
+
\frac12 m_B v_B^2
+
\frac{Gm_A m_B}{r_{AB}}.
\label{eq:newtonian-two-body-lagrangian}
\end{equation}
More generally, the \(N\)-body Newtonian sector reads
\begin{equation}
L_{\rm N}^{(N)}
=
\sum_A \frac12 m_A v_A^2
+
\sum_{A<B}\frac{Gm_A m_B}{r_{AB}}.
\label{eq:newtonian-nbody-lagrangian}
\end{equation}
Equation \eqref{eq:newtonian-two-body-lagrangian} is the baseline pair Lagrangian on top
of which the subsequent \(1\)PN self, cross, and static nonlinear sectors will be added.

\claimclass{\Controlled for the identification of the Poisson-hook scalar with the
effective Newtonian potential, the coherent-defect / small-body worldtube closure, and
the neglect of boundary-flux terms at the retained order. Exact for the Euler--Lagrange
equivalence between \cref{eq:newtonian-particle-lagrangian} and the leading monopole
equation of motion once the reduced Lagrangian has been stated.}


% ============================================================
\section{Scalar and optical dressing of the self sector}
\label{sec:self-sector}

This section isolates the part of the \(1\)PN pair Lagrangian proportional to
\((v_A^2+v_B^2)\). Conceptually this sector is distinct from the vector wake sector:
the wake overlap tensor is bilinear in the velocities of \emph{two} bodies and therefore
naturally produces terms such as \(\vct_A\!\cdot\!\vct_B\) and
\((\vct_A\!\cdot\!\nvec)(\vct_B\!\cdot\!\nvec)\), but not the one-body ``self''
combination \(v_A^2+v_B^2\). The self terms must therefore come from a different
channel. In the present framework that channel is the combined effect of
\begin{enumerate}[leftmargin=1.5em]
  \item scalar mass dressing in the Newtonian potential, and
  \item the optical/barotropic renormalization of the local limiting propagation speed.
\end{enumerate}

The logic of the section is straightforward. We first write the minimal reduced
scalar/optical worldline ansatz that captures both effects at the order needed for
\(1\)PN bookkeeping. Its Newtonian term fixes the scalar dressing coefficient \(q\).
Independently, the optical weak-field coefficient fixes the EOS exponent \(n\). Once
those two numbers are known, the coefficient of the mixed \(\PhiN v^2/c^2\) term follows
algebraically. We also record the universal free-particle \(v^4\) coefficient, since it
enters the same two-body \(1\)PN ledger.

A bookkeeping remark is important at the outset. In this section we keep the
\(\PhiN v^2/c^2\) channel and the free \(v^4/c^2\) channel, but we do \emph{not} use this
section to derive the static nonlinear \(\PhiN^2/c^2\) term. That term is conceptually
separate and is derived in \cref{sec:static-term} by local mass scaling plus pair
counting.

\subsection{Scalar-dressed particle action and Newtonian matching}
\label{sec:q-matching}

The reduced one-body ansatz we need is
\begin{equation}
L_{\rm sc}
=
-M_0\bigl(1+q\,\epsilon_\Phi\bigr)c^2
\sqrt{1-\frac{\epsilon_v^2}{1+(n-1)\epsilon_\Phi}},
\label{eq:scalar-optical-ansatz}
\end{equation}
with the dimensionless weak-field parameters
\begin{equation}
\epsilon_\Phi \equiv \frac{\PhiN}{c^2},
\qquad
\epsilon_v^2 \equiv \frac{v^2}{c^2}.
\label{eq:eps-defs}
\end{equation}
The prefactor \(1+q\,\epsilon_\Phi\) is the scalar mass dressing, while the denominator
\(1+(n-1)\epsilon_\Phi\) encodes the leading optical/barotropic shift in the local
propagation speed inherited from the same microphysics that controls the refractive
sector. Equation \eqref{eq:scalar-optical-ansatz} is a \Controlled reduction: it is the
minimal worldline form that keeps both channels visible while truncating to the order
relevant for conservative \(1\)PN motion.

Expanding \eqref{eq:scalar-optical-ansatz} to first order in \(\epsilon_\Phi\) and to
quartic order in \(v/c\) gives
\begin{align}
L_{\rm sc}
=
{}&-M_0 c^2
- q\,M_0\PhiN
+ \frac12 M_0 v^2
+ \frac18 M_0\frac{v^4}{c^2}
\nonumber\\
&\hspace{2em}
+ M_0\left(\frac q2-\frac{n-1}{2}\right)\frac{\PhiN v^2}{c^2}
+ O\!\left(\frac{v^6}{c^4},\frac{\PhiN^2}{c^2},\frac{\PhiN v^4}{c^4}\right).
\label{eq:scalar-optical-expanded}
\end{align}
The Newtonian worldline derived in \cref{sec:newtonian-limit} is
\begin{equation}
L_{\rm N}
=
\frac12 M_0 v^2 - M_0\PhiN.
\label{eq:newtonian-worldline-recalled}
\end{equation}
Comparing \eqref{eq:scalar-optical-expanded} with
\eqref{eq:newtonian-worldline-recalled}, the coefficient of the Newtonian potential fixes
\begin{equation}
q=1.
\label{eq:q-equals-one}
\end{equation}

This is the first key result of the section. The coefficient multiplying the scalar
potential is not left as a phenomenological parameter; once the one-body reduction is
required to reproduce the Newtonian point-particle limit, the scalar dressing exponent is
forced to unity.

\subsection{Optical consistency fixes the EOS exponent}
\label{sec:n-five}

The second ingredient is the optical weak-field coefficient. For the barotropic closure
\begin{equation}
P(\rho)=\KEOS \rho^n,
\label{eq:eos-self-section}
\end{equation}
the characteristic propagation speed obeys
\begin{equation}
c_s^2(\rho)=\frac{dP}{d\rho}=n\KEOS \rho^{n-1}.
\label{eq:sound-speed-self-section}
\end{equation}
Linearizing around a background density \(\rho_0\) with
\(\delta \equiv (\rho-\rho_0)/\rho_0\), one finds
\begin{equation}
\frac{\Delta c_s}{c}
\simeq
\frac{n-1}{2}\,\delta,
\qquad
N(\rho)\equiv \frac{c}{c_s(\rho)}
\simeq
1-\frac{n-1}{2}\,\delta.
\label{eq:index-linearized}
\end{equation}
In the weak static regime used throughout the series, the background density perturbation
satisfies
\begin{equation}
\delta \simeq \frac{\PhiN}{c^2},
\qquad
\PhiN(r)=-\frac{GM}{r},
\label{eq:delta-from-potential}
\end{equation}
so the far-zone refractive profile becomes
\begin{equation}
N(r)
\simeq
1+\alpha_n\frac{GM}{c^2 r},
\qquad
\alpha_n\equiv \frac{n-1}{2}.
\label{eq:alpha-n-def}
\end{equation}
The weak-field optical tests require the GR-matching coefficient
\begin{equation}
\alpha_n = 2,
\label{eq:alpha-n-two}
\end{equation}
and therefore
\begin{equation}
\boxed{\,n=5\,}.
\label{eq:n-equals-five}
\end{equation}

The role of this result should be read narrowly and correctly. We are not promoting the
optical refractive metric to the full effective metric for massive matter motion. What we
\emph{do} import from the optical sector is the stiffness exponent \(n\), because it is a
property of the same barotropic microphysics that enters the reduced scalar/optical
worldline ansatz \eqref{eq:scalar-optical-ansatz}. Once the optical sector has fixed
\(n\), that value can legitimately be carried into the massive-body \(1\)PN bookkeeping.

\subsection{1PN self coefficient from the scalar+optical expansion}
\label{sec:self-coeff}

The mixed coefficient of \(\PhiN v^2/c^2\) in
\eqref{eq:scalar-optical-expanded} is
\begin{equation}
\Cself
=
\frac q2-\frac{n-1}{2}.
\label{eq:self-coeff-general}
\end{equation}
Using \eqref{eq:q-equals-one} and \eqref{eq:n-equals-five},
\begin{equation}
\Cself
=
\frac12 - 2
=
-\frac32.
\label{eq:self-coeff-minus-three-halves}
\end{equation}

To translate this one-body coefficient into the pair Lagrangian, consider body \(A\) in
the field of body \(B\):
\begin{equation}
L_A
=
\frac12 m_A v_A^2
- m_A\Phi_B
+ \frac18 m_A\frac{v_A^4}{c^2}
+ m_A \Cself \frac{\Phi_B v_A^2}{c^2}
+\cdots,
\label{eq:body-A-self-expansion}
\end{equation}
with
\begin{equation}
\Phi_B(\mathbf X_A)=-\frac{Gm_B}{r_{AB}}.
\label{eq:phi-B-at-A}
\end{equation}
Therefore
\begin{equation}
m_A \Cself \frac{\Phi_B v_A^2}{c^2}
=
\frac{Gm_A m_B}{c^2 r_{AB}}
\bigl(-\Cself\bigr)v_A^2.
\label{eq:self-sign-translation-A}
\end{equation}
The same argument for body \(B\) gives
\begin{equation}
L_{\rm self}^{(AB)}
=
\frac{Gm_A m_B}{c^2 r_{AB}}
\bigl(-\Cself\bigr)\bigl(v_A^2+v_B^2\bigr).
\label{eq:self-pair-term-general}
\end{equation}
Since \(\Cself=-3/2\), the pair coefficient is
\begin{equation}
-\Cself=+\frac32.
\label{eq:self-pair-coeff-plus-three-halves}
\end{equation}
This is exactly the EIH self-sector coefficient multiplying
\((v_A^2+v_B^2)\) in the conservative two-body \(1\)PN Lagrangian.

Two points deserve emphasis. First, this coefficient is fixed without any use of the wake
overlap tensor; it belongs to a genuinely different sector of the theory. Second, the
derivation depends only on the already-fixed values \(q=1\) and \(n=5\). Once those are
established, the self coefficient is no longer an adjustable number.

\subsection{Free kinematic expansion}
\label{sec:kinematic-expansion}

The same reduced expansion also carries the universal free-particle quartic term. Setting
\(\PhiN=0\) in \eqref{eq:scalar-optical-ansatz} gives
\begin{equation}
L_{\rm free}
=
- M_0 c^2\sqrt{1-\frac{v^2}{c^2}}
=
- M_0 c^2
+ \frac12 M_0 v^2
+ \frac18 M_0\frac{v^4}{c^2}
+ O\!\left(\frac{v^6}{c^4}\right).
\label{eq:free-lagrangian-expansion}
\end{equation}
The coefficient \(1/8\) is the one needed in the two-body \(1\)PN \emph{Lagrangian}.
This should not be confused with the coefficient \(3/8\) that appears in the
low-velocity expansion of the relativistic \emph{energy},
\begin{equation}
E(v)=\gamma M_0 c^2
=
M_0 c^2
+ \frac12 M_0 v^2
+ \frac38 M_0\frac{v^4}{c^2}
+ \cdots.
\label{eq:energy-expansion-reminder}
\end{equation}
Both statements are correct; they refer to different functions related by Legendre
transform. In the present paper it is \eqref{eq:free-lagrangian-expansion} that enters
the EIH matching.

The physical interpretation is the same as in the bridge paper: if the defect's inertia
is wave-supported, then the free relativistic coefficients are not inserted by hand but
come from the trapped-mode kinematics. For the purposes of the present section, however,
all we need operationally is the Lagrangian coefficient \(1/8\).

\paragraph{Summary of the self-sector route.}
The direct scalar/optical route therefore fixes
\begin{equation}
q=1,
\qquad
n=5,
\qquad
\Cself=-\frac32,
\qquad
-\Cself=+\frac32
\ \text{in the pair Lagrangian}.
\label{eq:self-sector-summary}
\end{equation}
Later, in \cref{sec:ledger}, the same self-sector number will reappear from the
topology-plus-response ledger
\(\betaPN=\krho+\kadd+\kPV\). The equality of those two routes is an important internal
consistency check of the full \(1\)PN construction.

\claimclass{\Controlled. Equation \eqref{eq:scalar-optical-ansatz} is a reduced
scalar/optical worldline ansatz, Newtonian matching fixes \(q=1\), the optical weak-field
coefficient fixes \(n=5\), and the \(\PhiN v^2/c^2\) self coefficient then follows
algebraically. The free \(v^4\) coefficient in the Lagrangian is the universal
special-relativistic value \(1/8\).}


% ============================================================
\section{Static nonlinear term from local mass scaling}
\label{sec:static-term}

This section isolates the velocity-independent \(1\)PN interaction term. It is
conceptually separate from both of the sectors treated so far. The scalar/optical
self-sector of \cref{sec:self-sector} produces terms proportional to
\(\PhiN v^2/c^2\), so it vanishes when the bodies are instantaneously at rest. The wake
overlap sector, to be derived later, is bilinear in the body velocities and also vanishes
in the static limit. The standard EIH ledger nevertheless contains a nonzero
order-\(c^{-2}\) interaction even for static bodies. In the present framework that term
comes from a different mechanism: the effective mass entering the Newtonian interaction is
itself dressed by the local scalar environment.

This is the ``gravity gravitates'' sector of the reduced theory. In the full \(N\)-body
problem it is most naturally written as a local-potential-squared term, which can then be
reorganized into the familiar symmetric static EIH form. In the special case of two
bodies it collapses to the standard
\[
-\frac{G^2 m_A m_B (m_A+m_B)}{2 c^2 r_{AB}^2}
\]
interaction. The purpose of the section is to show that this term follows directly from
local mass scaling plus exact pair counting once the scalar dressing coefficient has
already been fixed by Newtonian matching.

\subsection{Local mass scaling as the source of the static sector}
\label{sec:mass-scaling-ansatz}

Let
\begin{equation}
\Phi_A^{\rm loc}(\mathbf X_A)
\equiv
-\sum_{C\neq A}\frac{Gm_C}{r_{AC}}
\label{eq:phi-local-def}
\end{equation}
denote the Newtonian potential at body \(A\) due to all \emph{other} bodies. The guiding
assumption is that the same scalar dressing that fixed the Newtonian one-body coupling
also dresses the local mass entering the pair-counted interaction energy:
\begin{equation}
m_A^{\rm eff}
=
m_A\left(1+\krho\,\frac{\Phi_A^{\rm loc}}{c^2}\right)
+ O(c^{-4}).
\label{eq:local-mass-scaling}
\end{equation}
Here \(\krho\) is simply the static-sector notation for the local mass-scaling
coefficient. In the present theory it is not a new free parameter: it is the same scalar
dressing coefficient that appeared as \(q\) in \cref{sec:self-sector}, so
\begin{equation}
\krho=q=1.
\label{eq:kappa-rho-equals-one}
\end{equation}

The pair-counted interaction part of the reduced Lagrangian is then written as
\begin{equation}
L_{\rm pot}
=
-\frac12\sum_A m_A^{\rm eff}\,\Phi_A^{\rm loc}.
\label{eq:pair-counted-potential}
\end{equation}
Expanding \eqref{eq:pair-counted-potential} with
\eqref{eq:local-mass-scaling} gives
\begin{equation}
L_{\rm pot}
=
-\frac12\sum_A m_A \Phi_A^{\rm loc}
-\frac{\krho}{2c^2}\sum_A m_A\bigl(\Phi_A^{\rm loc}\bigr)^2
+ O(c^{-4}).
\label{eq:pair-counted-expanded}
\end{equation}
The first term is the ordinary Newtonian interaction, while the second is the full static
nonlinear sector generated by the same local mass dressing. Using
\eqref{eq:phi-local-def}, the static term may be written compactly as
\begin{equation}
L_{\rm stat}
=
-\frac{\krho G^2}{2c^2}
\sum_A
m_A
\left(
\sum_{C\neq A}\frac{m_C}{r_{AC}}
\right)^2.
\label{eq:static-compact}
\end{equation}
Equation \eqref{eq:static-compact} is the key formula of the section. Once the reduced
ansatz \eqref{eq:local-mass-scaling} has been declared, the rest is exact algebra and
pair counting.

Before turning to the \(1\)PN piece, note that the Newtonian part of
\eqref{eq:pair-counted-expanded} reproduces the standard interaction exactly:
\begin{align}
-\frac12\sum_A m_A \Phi_A^{\rm loc}
&=
\frac12\sum_A\sum_{C\neq A}\frac{Gm_A m_C}{r_{AC}}
\nonumber\\
&=
\sum_{A<C}\frac{Gm_A m_C}{r_{AC}}.
\label{eq:newtonian-pair-counting}
\end{align}
Thus the same pair-counting rule yields both the ordinary Newtonian potential and its
static \(1\)PN correction.

\subsection{Two-body pair counting}
\label{sec:pair-counting}

For two bodies \(A\) and \(B\),
\begin{equation}
\Phi_A^{\rm loc}=-\frac{Gm_B}{r_{AB}},
\qquad
\Phi_B^{\rm loc}=-\frac{Gm_A}{r_{AB}}.
\label{eq:two-body-local-potentials}
\end{equation}
Substituting these into \eqref{eq:pair-counted-expanded} gives
\begin{align}
L_{\rm pot}^{(AB)}
&=
-\frac12\left[
m_A\left(1+\krho\frac{\Phi_A^{\rm loc}}{c^2}\right)\Phi_A^{\rm loc}
+
m_B\left(1+\krho\frac{\Phi_B^{\rm loc}}{c^2}\right)\Phi_B^{\rm loc}
\right]
+ O(c^{-4})
\nonumber\\[3pt]
&=
\frac{Gm_A m_B}{r_{AB}}
-\frac{\krho G^2 m_A m_B(m_A+m_B)}{2c^2 r_{AB}^2}
+ O(c^{-4}).
\label{eq:two-body-static-general}
\end{align}
Therefore the two-body static \(1\)PN contribution is
\begin{equation}
L_{\rm stat}^{(AB)}
=
-\frac{\krho G^2 m_A m_B(m_A+m_B)}{2c^2 r_{AB}^2}.
\label{eq:static-1pn-term-general}
\end{equation}
Using \eqref{eq:kappa-rho-equals-one}, this becomes
\begin{equation}
\boxed{
L_{\rm stat}^{(AB)}
=
-\frac{G^2 m_A m_B(m_A+m_B)}{2c^2 r_{AB}^2}.
}
\label{eq:static-1pn-term}
\end{equation}
This is exactly the static coefficient required in the conservative two-body EIH
Lagrangian.

Two comments are worth recording explicitly.

First, the coefficient \(-1/2\) is not imported from the EIH target. It is the direct
consequence of expanding the locally dressed mass only to first order in
\(\Phi_A^{\rm loc}/c^2\) and then pair-counting exactly once. Second, this sector carries
no new free parameter once Newtonian matching has fixed \(q=1\): the only coefficient
that appears here is \(\krho\), and in the present framework \(\krho=q=1\).

\subsection{Three-body consistency pattern}
\label{sec:three-body-pattern}

Although the present paper ultimately specializes to the two-body \(1\)PN ledger, the
mechanism above automatically generates the correct static pattern in the full many-body
problem. For three bodies \(A,B,C\),
\begin{align}
\Phi_A^{\rm loc} &= -G\left(\frac{m_B}{r_{AB}}+\frac{m_C}{r_{AC}}\right), \nonumber\\
\Phi_B^{\rm loc} &= -G\left(\frac{m_A}{r_{AB}}+\frac{m_C}{r_{BC}}\right), \nonumber\\
\Phi_C^{\rm loc} &= -G\left(\frac{m_A}{r_{AC}}+\frac{m_B}{r_{BC}}\right).
\label{eq:three-body-local-potentials}
\end{align}
Substituting into \eqref{eq:static-compact} and expanding gives both the pair-collapsed
\(r^{-2}\) pieces and the genuine triplet static terms:
\begin{align}
L_{\rm stat}^{(ABC)}
={}&
-\frac{\krho G^2}{2c^2}
\left[
\frac{m_A m_B(m_A+m_B)}{r_{AB}^2}
+
\frac{m_A m_C(m_A+m_C)}{r_{AC}^2}
+
\frac{m_B m_C(m_B+m_C)}{r_{BC}^2}
\right]
\nonumber\\
&
-\frac{\krho G^2 m_A m_B m_C}{c^2}
\left[
\frac{1}{r_{AB}r_{AC}}
+
\frac{1}{r_{AB}r_{BC}}
+
\frac{1}{r_{AC}r_{BC}}
\right].
\label{eq:three-body-static-expanded}
\end{align}
The important point is the pattern of the genuine three-body terms: each distinct product
of inverse separations appears with unit weight. This is exactly the coefficient pattern
checked in the symbolic harness. Depending on symmetrization convention, the same result
may be rewritten into the standard prime-sum EIH static sector.

For the present two-body paper we do not need the full many-body rewriting, but it is
useful to record \eqref{eq:three-body-static-expanded} for two reasons. First, it shows
that the mechanism responsible for the two-body \(r^{-2}\) term is not ad hoc; it extends
consistently to the full static nonlinear hierarchy. Second, it clarifies why this sector
is often described as the ``gravity gravitates'' term: once the local interaction mass is
itself sensitive to the ambient potential, the pair-counted Newtonian energy necessarily
feeds back into a quadratic potential contribution.

\paragraph{Summary of the static route.}
The static nonlinear term is therefore fixed by a short chain:
\begin{equation}
m_A^{\rm eff}(\Phi_A^{\rm loc})
=
m_A\left(1+\krho\frac{\Phi_A^{\rm loc}}{c^2}\right),
\qquad
\krho=q=1,
\qquad
L_{\rm stat}
=
-\frac{\krho}{2c^2}\sum_A m_A(\Phi_A^{\rm loc})^2.
\label{eq:static-route-summary}
\end{equation}
Specializing to two bodies then yields the exact coefficient
\begin{equation}
-\frac12
\quad\text{for the}\quad
\frac{G^2m_A m_B(m_A+m_B)}{c^2 r_{AB}^2}
\quad\text{term.}
\end{equation}
This is the static contribution that will enter the full two-body \(1\)PN assembly in
\cref{sec:full-eih}.

\claimclass{\Controlled for the local mass-scaling ansatz
\eqref{eq:local-mass-scaling}, which imports the already-fixed scalar dressing
coefficient into the static sector. Exact for the pair counting, the expansion of
\eqref{eq:pair-counted-potential}, the compact many-body form
\eqref{eq:static-compact}, and the resulting two-body coefficient
\eqref{eq:static-1pn-term} once the reduced ansatz has been declared.}


% ============================================================
\section{Topology and internal response: \texorpdfstring{$\kadd$}{kappaadd}, \texorpdfstring{$\kPV$}{kappaPV}, and the 1PN ledger}
\label{sec:ledger}

This section closes the scalar/self ledger from a second direction. In
\cref{sec:self-sector}, the coefficient of the mixed \(\PhiN v^2/c^2\) term was obtained
directly from the reduced scalar/optical worldline ansatz and the already-fixed values
\(q=1\) and \(n=5\). Here we reorganize the same self-sector bookkeeping into three
physically distinct contributions:
\begin{equation}
\betaPN = \krho + \kadd + \kPV.
\label{eq:beta-ledger-def}
\end{equation}
The first term, \(\krho\), is the scalar density/mass-dressing contribution already fixed
by Newtonian matching:
\begin{equation}
\krho = q = 1.
\label{eq:krho-one-ledger}
\end{equation}
The second term, \(\kadd\), comes from exterior added mass and is controlled by the
topology of the defect as seen by the surrounding flow. The third term, \(\kPV\), is an
internal breathing/pressure-volume response coefficient. Unlike \(\kadd\), it is not a
purely kinematic topology result; it becomes derived only after a response protocol has
been declared.

The logic of the section is therefore intentionally split. First, we compute
\(\kadd\) from the exterior potential-flow problem and show that the
\(w\)-uniform throat gives \(\kadd=1/2\), while a compact 4D bubble would give
\(1/3\). Second, we present one explicit adiabatic one-degree-of-freedom closure for the
internal breathing mode and show that it yields \(\kPV=3/2\). The sum then gives
\begin{equation}
\betaPN = 3,
\end{equation}
which reproduces the self-sector coefficient found earlier.

\subsection{Added mass as a topology discriminator}
\label{sec:added-mass}

The added-mass contribution is computed in the usual quasi-static, low-Mach potential-flow
regime. The assumptions are:
\begin{enumerate}[leftmargin=1.5em]
  \item the exterior rearrangement is slow enough that compressibility corrections are
  negligible at the retained order,
  \item the exterior flow is irrotational,
  \item the defect translates rigidly through the ambient medium at velocity
  \(\mathbf U\), and
  \item the surrounding flow decays sufficiently rapidly at infinity.
\end{enumerate}
Under these assumptions, the exterior velocity may be written as
\begin{equation}
\mathbf v_{\rm ext}=\grad \phi,
\qquad
\nabla^2 \phi = 0
\quad
\text{outside the excluded region.}
\label{eq:potential-flow-ledger}
\end{equation}
The kinetic energy stored in this exterior rearrangement is
\begin{equation}
T_{\rm ext}
=
\frac12 \rho_0 \int_{\rm ext} |\grad \phi|^2\, dV
=
\frac12 M_{\rm add} U^2,
\label{eq:added-mass-def}
\end{equation}
which defines the added mass \(M_{\rm add}\). Writing the displaced mass as
\begin{equation}
M_{\rm disp}=\rho_0 V_{\rm exc},
\label{eq:displaced-mass}
\end{equation}
we introduce the dimensionless added-mass coefficient
\begin{equation}
M_{\rm add}=\kadd M_{\rm disp}.
\label{eq:kadd-def}
\end{equation}

For a translating \(d_{\rm eff}\)-dimensional ball, the \(\ell=1\) harmonic of the
exterior potential-flow problem gives the standard result
\begin{equation}
\kadd(d_{\rm eff})=\frac{1}{d_{\rm eff}-1}.
\label{eq:kadd-general}
\end{equation}
In particular,
\begin{equation}
\kadd(3)=\frac12,
\qquad
\kadd(4)=\frac13.
\label{eq:kadd-3d-4d}
\end{equation}

The reason this matters for the present toy model is that the defect is not treated as a
compact \(4\)-ball. The object carried forward from the earlier papers is a
\(w\)-extended throat. In the regime relevant for the \(1\)PN reduction, the exterior
rearrangement is uniform along \(w\), so one may take
\begin{equation}
\phi(x,y,z,w)=\phi_3(x,y,z),
\qquad
\nabla_4^2 \phi = \nabla_3^2 \phi_3.
\label{eq:throat-potential-factorization}
\end{equation}
Each \(w\)-slice therefore solves the ordinary \(3\)-dimensional sphere problem, and the
total exterior kinetic energy is simply extensive in the throat length \(L\). The
effective exterior-flow dimension is thus
\begin{equation}
d_{\rm eff}=3,
\end{equation}
so the throat result is
\begin{equation}
\boxed{\kadd=\frac12.}
\label{eq:kadd-half}
\end{equation}

It is useful to record the counterfactual comparison explicitly. Had the defect instead
been a compact \(4\)D bubble \(B^4\), the exterior flow would have been genuinely
four-dimensional and one would obtain
\begin{equation}
\kadd^{(B^4)}=\frac13.
\label{eq:kadd-bubble}
\end{equation}
This contrast is conceptually important: \(\kadd\) is not a tunable inertial correction.
It is a topology discriminator. Once the object is specified to be a \(w\)-uniform
throat rather than a compact \(4\)D bubble, the added-mass contribution is fixed.

\subsection{Adiabatic one-degree-of-freedom closure for \texorpdfstring{$\kPV$}{kappaPV}}
\label{sec:kappaPV-closure}

The breathing/pressure-volume contribution is different in character. In the full parent
theory, the defect can exchange energy between trapped waves, exterior rearrangement, and
geometric breathing. Without a specific dynamical protocol, the coefficient \(\kPV\)
should remain symbolic. In this paper we make one explicit choice: an \emph{adiabatic}
one-degree-of-freedom closure in which the internal geometry relaxes quasistatically to a
local minimum of an effective rest-energy functional.

The protocol is:
\begin{enumerate}[leftmargin=1.5em]
  \item the external control parameter is the ambient density \(\rho\) at the defect
  location,
  \item the weak-field relation \(\delta \rho/\rho_0 \simeq \PhiN/c^2\) is used to link
  ambient density to the local Newtonian environment,
  \item the geometry is reduced to one breathing variable \(a\), with fixed aspect ratio
  \begin{equation}
  L=\Lambda a,
  \label{eq:fixed-aspect-ratio}
  \end{equation}
  and
  \item for each \(\rho\), the equilibrium radius \(a_*(\rho)\) minimizes the effective
  energy.
\end{enumerate}

Within this closure, we write the reduced equilibrium functional as
\begin{equation}
F(a,\rho)=E_w+E_f+E_{\rm PV},
\label{eq:F-split}
\end{equation}
with the sector scalings
\begin{equation}
E_w=\frac{\mathcal C_w\,c_s(\rho)}{a},
\qquad
E_f=\frac{\mathcal C_f}{\rho\,a^2},
\qquad
E_{\rm PV}=\mathcal C_{\rm PV}\,P(\rho)\,V(a).
\label{eq:sector-scalings}
\end{equation}
Using the same barotropic closure as before,
\begin{equation}
P(\rho)=\KEOS \rho^n,
\qquad
c_s(\rho)\propto \rho^{(n-1)/2},
\label{eq:closure-eos-ledger}
\end{equation}
and the fixed-aspect-ratio volume
\begin{equation}
V(a)=\pi \Lambda a^3,
\label{eq:volume-a}
\end{equation}
the \(a\)-dependence of the three sectors is \((-1,-2,+3)\). Equilibrium
\(\partial_a F=0\) therefore implies the virial relation
\begin{equation}
E_w+2E_f=3E_{\rm PV}.
\label{eq:virial-ledger}
\end{equation}

Define the dimensionless internal partition parameter
\begin{equation}
x\equiv \frac{E_f}{E_w}.
\label{eq:x-def}
\end{equation}
Then \eqref{eq:virial-ledger} gives
\begin{equation}
\frac{E_{\rm PV}}{E_w}=\frac{1+2x}{3},
\label{eq:EPV-over-Ew}
\end{equation}
and the total equilibrium energy becomes
\begin{equation}
F
=
E_w\left(1+x+\frac{1+2x}{3}\right)
=
E_w\,\frac{4+5x}{3}.
\label{eq:F-over-Ew}
\end{equation}

Because \(a_*(\rho)\) minimizes \(F\), the equilibrium logarithmic density slope obeys an
envelope-theorem form:
\begin{equation}
\frac{d\ln F_{\rm eq}}{d\ln \rho}
=
\frac{
\left(\frac{n-1}{2}\right)E_w
+
(-1)E_f
+
nE_{\rm PV}
}{F}.
\label{eq:envelope-slope-general}
\end{equation}
Substituting \eqref{eq:EPV-over-Ew} and \eqref{eq:F-over-Ew} yields
\begin{equation}
\frac{d\ln F_{\rm eq}}{d\ln \rho}
=
\frac{\frac{5n-3}{2} + (2n-3)x}{4+5x}.
\label{eq:envelope-slope-x}
\end{equation}
Using the already-fixed optical value \(n=5\), this reduces to
\begin{equation}
\frac{d\ln F_{\rm eq}}{d\ln \rho}
=
\frac{11+7x}{4+5x}.
\label{eq:dlnF-n5}
\end{equation}

We now connect this slope to the \(1\)PN bookkeeping. In the adiabatic closure, the
breathing-response contribution is estimated by subtracting the already-accounted scalar
mass-dressing piece:
\begin{equation}
\kPV
\equiv
\frac{d\ln F_{\rm eq}}{d\ln \rho}
-
\krho.
\label{eq:kPV-def}
\end{equation}
Since \(\krho=1\), the target value \(\kPV=3/2\) is equivalent to
\begin{equation}
\frac{d\ln F_{\rm eq}}{d\ln \rho}
=
\krho+\kPV
=
\frac52.
\label{eq:dlnF-target}
\end{equation}
Applying this to \eqref{eq:dlnF-n5} gives
\begin{equation}
\frac{11+7x}{4+5x}=\frac52
\qquad\Longrightarrow\qquad
x=\frac{2}{11}.
\label{eq:x-two-elevenths}
\end{equation}
Therefore the adiabatic one-degree-of-freedom closure predicts
\begin{equation}
\boxed{\kPV=\frac32.}
\label{eq:kPV-three-halves}
\end{equation}

The same solution fixes the internal energy partition:
\begin{equation}
E_w:E_f:E_{\rm PV}=11:2:5,
\label{eq:partition-11-2-5}
\end{equation}
or, equivalently,
\begin{equation}
(f_w,f_f,f_{\rm PV})
=
\left(\frac{11}{18},\frac{2}{18},\frac{5}{18}\right).
\label{eq:partition-fractions}
\end{equation}
This is a useful point to emphasize. Even though \(\kPV\) is a protocol-level quantity,
once the protocol is declared its value is not merely fitted; it comes with a definite
predicted partition of internal energy among wave, flow, and pressure-volume sectors.

It is also worth recording one secondary prediction of the same closure. Differentiating
the equilibrium condition gives the breathing response slope
\begin{equation}
\frac{d\ln a_*}{d\ln \rho}
=
-\frac{57}{64}
\approx -0.890625
\qquad
(n=5,\ x=2/11).
\label{eq:breathing-slope}
\end{equation}
This number is not needed directly in the conservative two-body \(1\)PN assembly, but it
is a useful internal diagnostic of the adiabatic closure and a concrete target for future
numerical tests.

\subsection{The 1PN self-sector ledger}
\label{sec:self-ledger}

We now combine the three contributions. The self-sector/precession ledger is
\begin{equation}
\betaPN = \krho + \kadd + \kPV.
\label{eq:beta-ledger}
\end{equation}
Using
\begin{equation}
\krho=1,
\qquad
\kadd=\frac12,
\qquad
\kPV=\frac32,
\end{equation}
we obtain
\begin{equation}
\boxed{\betaPN=3.}
\label{eq:beta-three}
\end{equation}

This is the value that later controls the orbit-level perihelion coefficient, but it
already has an immediate algebraic consequence for the present paper. In one-body form,
the self-sector coefficient multiplying \(\PhiN v^2/c^2\) may be written as
\begin{equation}
-\frac{\betaPN}{2}
=
-\frac32.
\label{eq:beta-to-self}
\end{equation}
This is exactly the coefficient found directly in \cref{sec:self-sector}:
\begin{equation}
\Cself=-\frac32.
\label{eq:self-coeff-repeat}
\end{equation}
Equivalently, in the two-body pair Lagrangian one has
\begin{align}
L_{\rm self}^{(AB)}
&=
-\frac{\betaPN}{2c^2}
\left[
m_A\Phi_B(\mathbf X_A)\,v_A^2
+
m_B\Phi_A(\mathbf X_B)\,v_B^2
\right]
\nonumber\\[3pt]
&=
\frac{\betaPN}{2}\,
\frac{Gm_A m_B}{c^2 r_{AB}}
\bigl(v_A^2+v_B^2\bigr),
\label{eq:pair-self-beta}
\end{align}
so \(\betaPN=3\) implies the pair coefficient
\begin{equation}
+\frac{\betaPN}{2}=+\frac32,
\label{eq:pair-self-plus-three-halves}
\end{equation}
again in exact agreement with \cref{sec:self-sector}.

This agreement matters conceptually. The coefficient \(+3/2\) in the two-body pair
Lagrangian is not resting on a single derivation route. It is reproduced both by the
direct scalar/optical expansion and by the independent topology-plus-response ledger. The
former fixes the coefficient from the reduced worldline structure; the latter explains how
that same number decomposes into scalar mass dressing, exterior added mass, and internal
breathing response.

\paragraph{What is fixed and what remains conditional.}
The distinction between \(\kadd\) and \(\kPV\) should be kept explicit.

\begin{itemize}[leftmargin=1.5em]
  \item \(\kadd=1/2\) is a geometry/topology result for a \(w\)-uniform throat in the
  controlled low-Mach potential-flow regime.
  \item \(\kPV=3/2\) is \emph{not} claimed as an exact theorem of the full parent theory.
  It is the value selected by the particular adiabatic one-degree-of-freedom closure used
  in this paper.
\end{itemize}

Within the scope of that declared closure hierarchy, however, the ledger is fully closed:
\begin{equation}
\betaPN=3.
\end{equation}

\claimclass{\Controlled for the exterior added-mass result
\(\kadd=1/2\), which depends on the quasi-static potential-flow regime and on the throat
topology. \Protocol for the adiabatic one-degree-of-freedom closure yielding
\(\kPV=3/2\). Exact for the algebraic ledger relation
\(\betaPN=\krho+\kadd+\kPV\) and for its equivalence to the pair self coefficient once the
individual ingredients have been supplied.}


% ============================================================
\section{Constructive wake derivation of the EIH cross terms}
\label{sec:wake}

This section derives the velocity-dependent \emph{cross} sector of the conservative
two-body \(1\)PN interaction. It is crucial to separate this sector cleanly from the
self-sector derived in \cref{sec:self-sector}. A bilinear overlap of the wake produced by
body \(A\) with the wake produced by body \(B\) is, by construction, bilinear in
\(\vct_A\) and \(\vct_B\). It can therefore generate terms proportional to
\(\vct_A\!\cdot\!\vct_B\) and
\((\vct_A\!\cdot\!\hat{\mathbf n}_{AB})(\vct_B\!\cdot\!\hat{\mathbf n}_{AB})\), but it
cannot generate the one-body combination \(v_A^2+v_B^2\). The latter belongs to the
scalar/optical and topology/response routes already treated. The purpose of the present
section is narrower: to show that the full EIH cross tensor follows constructively from an
isotropic linear-response wake basis with no additional phenomenological offsets.

The result is that the parity-even conservative cross sector is fixed by the unique point
\begin{equation}
\alpha^2=\frac34,
\qquad
a_H=0,
\qquad
K_{\rm vec}=\frac{2}{\pi^2},
\label{eq:wake-unique-point}
\end{equation}
which reproduces the standard EIH coefficients
\begin{equation}
C_\parallel=-\frac72,
\qquad
C_L=-\frac12.
\label{eq:wake-eih-targets}
\end{equation}
The logic proceeds in three stages:
\begin{enumerate}[leftmargin=1.5em]
  \item declare an isotropic Fourier-space wake basis and evaluate the overlap tensor;
  \item match the resulting tensor to the EIH cross structures, yielding a one-parameter
  family of solutions; and
  \item use the already-fixed EOS exponent \(n=5\) to collapse that family to the unique
  parity-even point \eqref{eq:wake-unique-point}.
\end{enumerate}

\subsection{Target cross sector}
\label{sec:wake-target}

For two bodies \(A,B\), define
\begin{equation}
\mathbf r_{AB}=\mathbf X_A-\mathbf X_B,
\qquad
r_{AB}=|\mathbf r_{AB}|,
\qquad
\hat{\mathbf n}_{AB}=\frac{\mathbf r_{AB}}{r_{AB}}.
\label{eq:wake-separation-defs}
\end{equation}
The most general rotationally invariant parity-even bilinear cross sector at this order is
\begin{equation}
L_{\rm vec}^{(AB)}
=
\frac{Gm_A m_B}{c^2 r_{AB}}
\left[
C_\parallel\,\vct_A\!\cdot\!\vct_B
+
C_L\,(\vct_A\!\cdot\!\hat{\mathbf n}_{AB})
     (\vct_B\!\cdot\!\hat{\mathbf n}_{AB})
\right].
\label{eq:wake-general-cross}
\end{equation}
The EIH target values are
\begin{equation}
C_\parallel^{\rm(EIH)}=-\frac72,
\qquad
C_L^{\rm(EIH)}=-\frac12.
\label{eq:wake-eih-target}
\end{equation}
The task is to derive these numbers from the toy model rather than insert them.

\subsection{Isotropic Fourier-space wake basis}
\label{sec:wake-basis}

We assume a linear far-field response in which the translational disturbance created by a
moving defect is represented in Fourier space as a sum of three isotropic basis
contributions,
\begin{equation}
\mathbf u(\mathbf k;\mathbf v)
=
\mathbf u_T(\mathbf k;\mathbf v)
+
\mathbf u_H(\mathbf k;\mathbf v)
+
\mathbf u_L(\mathbf k;\mathbf v).
\label{eq:wake-total-basis}
\end{equation}
These are taken to be
\begin{align}
\mathbf u_T(\mathbf k;\mathbf v)
&=
i\,a_T\,
\frac{\mathbf P_T(\hat{\mathbf k})\,\mathbf v}{k},
\label{eq:wake-uT}
\\[3pt]
\mathbf u_H(\mathbf k;\mathbf v)
&=
i\,a_H\,
\frac{\mathbf k\times \mathbf v}{k^2},
\label{eq:wake-uH}
\\[3pt]
\mathbf u_L(\mathbf k;\mathbf v)
&=
i\,a_L\,
\frac{\mathbf k(\mathbf k\!\cdot\!\mathbf v)}{k^3},
\label{eq:wake-uL}
\end{align}
with
\begin{equation}
\hat{\mathbf k}=\frac{\mathbf k}{k},
\qquad
\mathbf P_T(\hat{\mathbf k})=\mathbf I-\hat{\mathbf k}\hat{\mathbf k},
\qquad
a_T\equiv 1.
\label{eq:wake-projectors}
\end{equation}
The three pieces have the intended interpretations:
\begin{itemize}[leftmargin=1.5em]
  \item \(\mathbf u_T\): transverse projector / translation-shear response,
  \item \(\mathbf u_H\): helical transverse response,
  \item \(\mathbf u_L\): longitudinal (compressible) response.
\end{itemize}

Each mode is chosen to scale as \(v/k\). This is the minimal scaling that makes the
bilinear overlap behave as \(u_A\!\cdot\!u_B\sim 1/k^2\), so that the Fourier integral
naturally produces a \(1/r\) pair interaction. The assumptions of this section are thus:
\begin{enumerate}[leftmargin=1.5em]
  \item a linear far-field response regime,
  \item isotropy in Fourier space,
  \item harmonic \(1/r\)-type near-zone pair scaling after the transform, and
  \item restriction to the parity-even conservative sector.
\end{enumerate}
No claim is made here about dissipative or odd-parity response channels.

\subsection{Bilinear overlap and isotropic tensor decomposition}
\label{sec:wake-overlap}

The pair overlap kernel is taken in the form
\begin{equation}
I_{AB}(r_{AB})
\propto
\int d^3k\;
\mathbf u(\mathbf k;\vct_A)\!\cdot\!\mathbf u(-\mathbf k;\vct_B)\,
e^{i\mathbf k\cdot \mathbf r_{AB}},
\label{eq:wake-overlap-kernel}
\end{equation}
with the overall normalization absorbed later into a dimensionless coupling
\(K_{\rm vec}\). Because the integrand is bilinear in \(\vct_A\) and \(\vct_B\), isotropy
implies that after integration the result must take the form
\begin{equation}
I_{AB}(r_{AB})
=
\frac{1}{r_{AB}}
\left[
\mathcal S_\parallel\,
\vct_A\!\cdot\!\vct_B
+
\mathcal S_L\,
(\vct_A\!\cdot\!\hat{\mathbf n}_{AB})
(\vct_B\!\cdot\!\hat{\mathbf n}_{AB})
\right].
\label{eq:wake-isotropic-decomposition}
\end{equation}

A convenient way to evaluate \eqref{eq:wake-overlap-kernel} is to align the separation
vector with the \(z\)-axis while doing the integral, so that
\begin{equation}
\hat{\mathbf n}_{AB}=\hat{\mathbf z},
\qquad
(\vct_A\!\cdot\!\hat{\mathbf n}_{AB})
(\vct_B\!\cdot\!\hat{\mathbf n}_{AB})
=
v_{Az}v_{Bz}.
\end{equation}
Carrying out the angular and radial integrations then yields the exact shape coefficients
within the declared wake basis:
\begin{equation}
\mathcal S_\parallel
=
\pi^2\bigl(-1+a_H^2-a_L^2\bigr),
\qquad
\mathcal S_L
=
\pi^2\bigl(-1+a_H^2+a_L^2\bigr).
\label{eq:wake-shapes}
\end{equation}
Only the squares \(a_H^2\) and \(a_L^2\) appear in the parity-even sector, as expected.

It is useful at this point to define
\begin{equation}
\alpha^2 \equiv a_L^2,
\label{eq:alpha-from-aL}
\end{equation}
where \(\alpha^2\) is the longitudinal/compressible wake fraction carried forward in the
rest of the paper. Introducing the dimensionless overall normalization \(K_{\rm vec}\),
the cross coefficients in \eqref{eq:wake-general-cross} become
\begin{equation}
C_\parallel
=
K_{\rm vec}\pi^2\bigl(-1+a_H^2-\alpha^2\bigr),
\qquad
C_L
=
K_{\rm vec}\pi^2\bigl(-1+a_H^2+\alpha^2\bigr).
\label{eq:wake-C-parallel-long}
\end{equation}

This equation is the central output of the constructive wake calculation. It shows that
the entire parity-even cross tensor is controlled by only three quantities:
\(K_{\rm vec}\), \(\alpha^2\), and \(a_H^2\).

\subsection{EIH matching gives a one-parameter family}
\label{sec:wake-family}

Imposing the EIH targets \eqref{eq:wake-eih-target} on
\eqref{eq:wake-C-parallel-long} gives
\begin{align}
-\frac72
&=
K_{\rm vec}\pi^2\bigl(-1+a_H^2-\alpha^2\bigr),
\label{eq:wake-match-para}
\\[3pt]
-\frac12
&=
K_{\rm vec}\pi^2\bigl(-1+a_H^2+\alpha^2\bigr).
\label{eq:wake-match-long}
\end{align}
Taking the difference and the sum isolates the two invariants:
\begin{align}
C_L-C_\parallel
&=
2K_{\rm vec}\pi^2\alpha^2
=
3,
\label{eq:wake-difference}
\\[3pt]
C_L+C_\parallel
&=
2K_{\rm vec}\pi^2(-1+a_H^2)
=
-4.
\label{eq:wake-sum}
\end{align}
Equivalently,
\begin{equation}
K_{\rm vec}\alpha^2=\frac{3}{2\pi^2},
\qquad
K_{\rm vec}(1-a_H^2)=\frac{2}{\pi^2}.
\label{eq:wake-invariants}
\end{equation}
Solving for \(\alpha^2\) and \(K_{\rm vec}\) in terms of \(a_H^2\) gives
\begin{equation}
\alpha^2(a_H^2)=\frac34(1-a_H^2),
\qquad
K_{\rm vec}(a_H^2)=\frac{2}{\pi^2(1-a_H^2)},
\qquad
0\le a_H^2<1.
\label{eq:wake-family}
\end{equation}

Thus EIH matching alone does \emph{not} yet pick a unique point. It fixes two invariants
and leaves one apparent freedom, which may be parameterized by \(a_H^2\). Put
differently, the cross-tensor structure is already highly constrained by the wake basis,
but one more physical input is needed to collapse the family.

\subsection{Thermodynamic closure from the already-fixed EOS exponent}
\label{sec:wake-thermo}

The additional input comes from the same barotropic microphysics that already fixed
\(n=5\) in the optical sector. The action-level barotropic identity is
\begin{equation}
P=\rho h-U,
\label{eq:wake-thermo-identity}
\end{equation}
and for a polytrope
\begin{equation}
P(\rho)=K_{\rm EOS}\rho^n,
\label{eq:wake-thermo-polytrope}
\end{equation}
the compatible internal energy density is
\begin{equation}
U(\rho)=\frac{K_{\rm EOS}}{n-1}\rho^n,
\qquad
\frac{U}{P}=\frac{1}{n-1}.
\label{eq:wake-UP-ratio}
\end{equation}
The minimal thermodynamic closure used here identifies the longitudinal wake fraction with
the non-internal part of the barotropic budget:
\begin{equation}
\alpha^2_{\rm thermo}
\equiv
1-\frac{U}{P}
=
1-\frac{1}{n-1}.
\label{eq:wake-thermo-alpha}
\end{equation}
Since the optical sector already fixed
\begin{equation}
n=5,
\label{eq:wake-n-five}
\end{equation}
we obtain
\begin{equation}
\boxed{
\alpha^2=\alpha^2_{\rm thermo}=1-\frac14=\frac34.
}
\label{eq:wake-alpha-three-quarters}
\end{equation}

This is the key closure step. The cross sector is no longer independent of the scalar and
optical sectors: once the EOS exponent has been fixed by weak-field optics, the same
microphysics fixes the longitudinal wake fraction.

\subsection{Collapse to the unique parity-even point}
\label{sec:wake-unique}

Substituting \eqref{eq:wake-alpha-three-quarters} into the EIH family
\eqref{eq:wake-family} gives
\begin{equation}
\frac34=\frac34(1-a_H^2)
\qquad\Longrightarrow\qquad
a_H^2=0.
\label{eq:wake-aH-zero}
\end{equation}
Hence
\begin{equation}
a_H=0
\label{eq:wake-aH-zero-sign}
\end{equation}
for the minimal real parity-even solution, and the normalization becomes
\begin{equation}
K_{\rm vec}=\frac{2}{\pi^2}.
\label{eq:wake-K-two-over-pi2}
\end{equation}
Equivalently, one may write the selected amplitude set as
\begin{equation}
a_T=1,
\qquad
a_H=0,
\qquad
a_L=\pm\frac{\sqrt3}{2},
\label{eq:wake-amplitudes}
\end{equation}
where only \(a_L^2=\alpha^2=3/4\) enters the conservative parity-even sector.

Back-substituting
\eqref{eq:wake-alpha-three-quarters}, \eqref{eq:wake-aH-zero-sign}, and
\eqref{eq:wake-K-two-over-pi2} into \eqref{eq:wake-C-parallel-long} yields
\begin{align}
C_\parallel
&=
\frac{2}{\pi^2}\pi^2\left(-1-\frac34\right)
=
-\frac72,
\label{eq:wake-backsub-para}
\\[3pt]
C_L
&=
\frac{2}{\pi^2}\pi^2\left(-1+\frac34\right)
=
-\frac12.
\label{eq:wake-backsub-long}
\end{align}
Thus the full parity-even cross tensor is reproduced exactly:
\begin{equation}
\boxed{
C_\parallel=-\frac72,
\qquad
C_L=-\frac12.
}
\label{eq:wake-final-cross}
\end{equation}

\paragraph{Why no additional offset belongs here.}
It is worth stating explicitly what this section does \emph{not} do. No scalar or optical
offset is added to \(C_L\) or \(C_\parallel\). Such a step would double-count physics
already assigned to the self-sector. The wake overlap is bilinear in the two body
velocities and therefore already exhausts the parity-even cross tensor at this order. The
self terms belong elsewhere; the cross terms belong here.

\paragraph{Summary of the wake route.}
The cross sector therefore closes through the chain
\begin{equation}
\text{isotropic wake basis}
\;\longrightarrow\;
\eqref{eq:wake-C-parallel-long}
\;\longrightarrow\;
\text{EIH family } \eqref{eq:wake-family}
\;\longrightarrow\;
\alpha^2=1-\frac{1}{n-1}
\;\longrightarrow\;
n=5
\;\longrightarrow\;
\left(\alpha^2,a_H,K_{\rm vec}\right)
=
\left(\frac34,0,\frac{2}{\pi^2}\right),
\label{eq:wake-route-summary}
\end{equation}
and finally
\begin{equation}
\left(C_\parallel,C_L\right)
=
\left(-\frac72,-\frac12\right).
\end{equation}
This is the cross-velocity input used in the full two-body \(1\)PN assembly of
\cref{sec:full-eih}.

\claimclass{\Controlled for the isotropic linear-response wake basis, the harmonic
far-field overlap ansatz, the minimal-real-solution criterion, and the thermodynamic
closure \(\alpha^2=1-1/(n-1)\). Exact for the algebraic tensor decomposition,
the EIH-family relations \eqref{eq:wake-family}, and the back-substitution that reproduces
\((C_\parallel,C_L)=(-7/2,-1/2)\) once the closure inputs have been supplied.}

% ============================================================
\section{Full two-body 1PN Lagrangian and exact EIH match}
\label{sec:full-eih}

This is the assembly section. No new dynamical ingredient is introduced here, and no new
coefficient is fitted. Once the Newtonian reduction, the scalar/optical self sector, the
static nonlinear sector, and the wake cross tensor have all been fixed, the remaining task
is purely algebraic: collect the previously derived terms into a single reduced two-body
Lagrangian and compare the result with the standard conservative Einstein--Infeld--Hoffmann
(EIH) target.

As throughout, we suppress the additive rest-energy constants
\(-m_A c^2-m_B c^2\), since they do not affect the equations of motion. We also use
\begin{equation}
\mathbf r_{AB}=\mathbf X_A-\mathbf X_B,
\qquad
r_{AB}=|\mathbf r_{AB}|,
\qquad
\nvec\equiv \hat{\mathbf n}_{AB}=\frac{\mathbf r_{AB}}{r_{AB}}.
\label{eq:full-eih-separation}
\end{equation}

\subsection{Derived two-body Lagrangian}
\label{sec:derived-2body}

The final two-body \(1\)PN Lagrangian is the sum of five pieces:
\begin{equation}
L_{\rm 1PN}^{\rm derived}
=
L_{\rm N}^{(AB)}
+
L_{v^4}^{(AB)}
+
L_{\rm self}^{(AB)}
+
L_{\rm cross}^{(AB)}
+
L_{\rm stat}^{(AB)}.
\label{eq:assembly-split}
\end{equation}
Each piece has already been derived in an earlier section.

\paragraph{Newtonian sector.}
From \cref{sec:newtonian-limit},
\begin{equation}
L_{\rm N}^{(AB)}
=
\frac12 m_A v_A^2
+
\frac12 m_B v_B^2
+
\frac{Gm_A m_B}{r_{AB}}.
\label{eq:assembly-newtonian}
\end{equation}

\paragraph{Universal quartic kinematics.}
From \cref{sec:kinematic-expansion},
\begin{equation}
L_{v^4}^{(AB)}
=
\frac{m_A v_A^4+m_B v_B^4}{8c^2}.
\label{eq:assembly-v4}
\end{equation}

\paragraph{Self sector.}
From \cref{sec:self-sector,sec:ledger}, the scalar/optical route and the
topology/response ledger agree on the same pair coefficient:
\begin{equation}
L_{\rm self}^{(AB)}
=
\frac{Gm_A m_B}{c^2 r_{AB}}\,
\frac32\bigl(v_A^2+v_B^2\bigr).
\label{eq:assembly-self}
\end{equation}

\paragraph{Cross sector.}
From \cref{sec:wake}, the constructive wake calculation yields
\(\Cpara=-7/2\) and \(\CLong=-1/2\), so
\begin{equation}
L_{\rm cross}^{(AB)}
=
\frac{Gm_A m_B}{c^2 r_{AB}}
\left[
-\frac72\,\vct_A\!\cdot\!\vct_B
-\frac12\,(\vct_A\!\cdot\!\nvec)(\vct_B\!\cdot\!\nvec)
\right].
\label{eq:assembly-cross}
\end{equation}

\paragraph{Static nonlinear sector.}
From \cref{sec:static-term},
\begin{equation}
L_{\rm stat}^{(AB)}
=
-\frac{G^2m_A m_B(m_A+m_B)}{2c^2 r_{AB}^2}.
\label{eq:assembly-static}
\end{equation}

Substituting \cref{eq:assembly-newtonian,eq:assembly-v4,eq:assembly-self,eq:assembly-cross,eq:assembly-static}
into \cref{eq:assembly-split} gives the final assembled expression
\begin{align}
L_{\rm 1PN}^{\rm derived}
={}&
\frac12 m_A v_A^2 + \frac12 m_B v_B^2
+ \frac{m_A v_A^4 + m_B v_B^4}{8c^2}
+ \frac{Gm_A m_B}{r_{AB}}
\nonumber\\
&
+ \frac{Gm_A m_B}{c^2 r_{AB}}
\left[
\frac32\bigl(v_A^2+v_B^2\bigr)
-\frac72\,\vct_A\!\cdot\!\vct_B
-\frac12\,(\vct_A\!\cdot\!\nvec)(\vct_B\!\cdot\!\nvec)
\right]
- \frac{G^2m_A m_B(m_A+m_B)}{2c^2 r_{AB}^2}.
\label{eq:derived-eih}
\end{align}

Equation \eqref{eq:derived-eih} is the main conservative \(1\)PN output of the paper.
Every coefficient in it has now been traced back to a named ingredient of the \(4\)D
model:
the point-particle Newtonian limit, wave-supported free kinematics, scalar/optical
dressing, topology and internal response, and the constructive wake overlap tensor.
Nothing remains to be chosen at this stage.

It is useful to emphasize the division of labor among sectors. The coefficient \(+3/2\)
multiplying \(v_A^2+v_B^2\) is \emph{not} part of the wake calculation and is reproduced
independently by the scalar/optical route and the
\(\betaPN=\krho+\kadd+\kPV\) ledger. Conversely, the coefficients \(-7/2\) and \(-1/2\)
are genuinely wake-sector outputs and are not dressed by additional scalar offsets. The
static \(G^2/r^2\) term is a third independent sector, fixed by local mass scaling and
pair counting. Equation \eqref{eq:derived-eih} is therefore a true assembly of distinct
previous derivations, not a single reparameterized ansatz.

\subsection{Standard EIH target and exact equality}
\label{sec:eih-target}

In the same Lagrangian convention used throughout this paper, the standard conservative
two-body Einstein--Infeld--Hoffmann target is
\begin{align}
L_{\rm 1PN}^{\rm EIH}
={}&
\frac12 m_A v_A^2 + \frac12 m_B v_B^2
+ \frac{m_A v_A^4 + m_B v_B^4}{8c^2}
+ \frac{Gm_A m_B}{r_{AB}}
\nonumber\\
&
+ \frac{Gm_A m_B}{c^2 r_{AB}}
\left[
\frac32\bigl(v_A^2+v_B^2\bigr)
-\frac72\,\vct_A\!\cdot\!\vct_B
-\frac12\,(\vct_A\!\cdot\!\nvec)(\vct_B\!\cdot\!\nvec)
\right]
- \frac{G^2m_A m_B(m_A+m_B)}{2c^2 r_{AB}^2}.
\label{eq:eih-target}
\end{align}
Comparing \cref{eq:derived-eih,eq:eih-target}, the residual vanishes identically:
\begin{equation}
\Delta L
\equiv
L_{\rm 1PN}^{\rm derived}-L_{\rm 1PN}^{\rm EIH}
=0.
\label{eq:eih-residual-zero}
\end{equation}

This is the sense in which the paper achieves a \emph{full} conservative two-body \(1\)PN
derivation within the stated closure hierarchy. The equality in
\cref{eq:eih-residual-zero} is not a fit, and it is not a qualitative resemblance. Once
the previously derived coefficients are imported, the final answer is exactly the standard
EIH Lagrangian.

A useful way to read the result is coefficient by coefficient. The entire provenance of
\cref{eq:derived-eih} is summarized in \cref{tab:coefficient-provenance}.

\begin{table}[t]
  \centering
  \caption{Coefficient provenance in the final two-body \(1\)PN Lagrangian.}
  \label{tab:coefficient-provenance}
  \small
  \renewcommand{\arraystretch}{1.15}
  \begin{tabularx}{\linewidth}{@{}L c L@{}}
    \toprule
    Coefficient in \cref{eq:derived-eih} & Value & Derived in \\
    \midrule
    \(v^4\) coefficient & \(1/8\) & \cref{sec:kinematic-expansion} \\
    \(v_A^2+v_B^2\) coefficient & \(+3/2\) & \cref{sec:self-sector,sec:ledger} \\
    \(\vct_A\!\cdot\!\vct_B\) coefficient & \(-7/2\) & \cref{sec:wake} \\
    \((\vct_A\!\cdot\!\nvec)(\vct_B\!\cdot\!\nvec)\) coefficient & \(-1/2\) & \cref{sec:wake} \\
    Static \(G^2/r^2\) coefficient & \(-1/2\) & \cref{sec:static-term} \\
    \bottomrule
  \end{tabularx}
\end{table}

Two final remarks help fix the scope of the match.

First, the equality established here is a statement about the \emph{conservative}
near-zone two-body sector through order \(c^{-2}\). Nothing in
\cref{eq:eih-residual-zero} addresses dissipative radiation reaction, spin couplings, or
higher-order post-Newtonian structure.

Second, the present section contains no remaining protocol dependence of its own. The only
protocol-level input entering the final assembly is \(\kPV=3/2\), and that quantity was
already isolated and labeled as a protocol closure in \cref{sec:ledger}. Once that input
is accepted, the rest of the final match is exact algebra.

\claimclass{\FullMatch. Once the previously derived ingredients are imported, the equality
to the standard conservative two-body EIH target is exact algebra, with no additional
fitting or coefficient freedom in this section.}

% ============================================================
\section{Test-mass orbit reduction and perihelion shift}
\label{sec:orbit}

This section extracts the most familiar observable consequence of the \(1\)PN ledger:
the perihelion shift in the test-mass limit. No new coefficient is fixed here. The point
is instead to show that the full two-body result of \cref{sec:full-eih}, when reduced to a
light body orbiting a heavy central source, collapses to the same orbital-sector
bookkeeping used earlier in the series. In that reduction, the entire precession-relevant
\(1\)PN correction is controlled by the single combination
\begin{equation}
\betaPN=\krho+\kadd+\kPV=3.
\label{eq:orbit-beta-three}
\end{equation}
The final perihelion shift is then
\begin{equation}
\Delta\phi_{\rm model}
=
\frac{2\pi\betaPN\mu}{c^2 a(1-e^2)},
\label{eq:orbit-deltaphi-general-intro}
\end{equation}
which becomes the standard GR coefficient once \(\betaPN=3\).

The section is organized in two steps. First, we take the test-mass limit of the full
assembled \(1\)PN Lagrangian and show how it reduces to a compact orbit Lagrangian with a
position-dependent kinetic prefactor. Second, we derive the corresponding Binet equation
and the perihelion advance.

\subsection{Reduced orbit Lagrangian}
\label{sec:orbit-lagrangian}

Let body \(A\) be a light test mass \(m\), and body \(B\) a heavy central mass \(M\), with
\begin{equation}
m\ll M,
\qquad
\mu\equiv GM.
\label{eq:test-mass-limit-def}
\end{equation}
Work in the central-body rest frame,
\begin{equation}
\mathbf X_B=\mathbf 0,
\qquad
\vct_B=\mathbf 0,
\qquad
r\equiv r_{AB}=|\mathbf X_A|.
\label{eq:central-rest-frame}
\end{equation}
Dividing the full two-body Lagrangian \eqref{eq:derived-eih} by the test mass \(m\) and
dropping terms suppressed by further powers of \(m/M\), the specific test-particle
Lagrangian becomes
\begin{equation}
\mathcal L_{\rm test}
=
\frac12 v^2
+\frac{\mu}{r}
+\frac{1}{c^2}
\left[
\frac18 v^4
+\frac32\,\frac{\mu}{r}\,v^2
-\frac12\,\frac{\mu^2}{r^2}
\right],
\label{eq:test-particle-lagrangian-full}
\end{equation}
where
\begin{equation}
v^2=\dot r^2+r^2\dot\phi^2.
\label{eq:polar-vsq}
\end{equation}

Equation \eqref{eq:test-particle-lagrangian-full} is already the correct \(1\)PN
test-mass limit of the derived EIH sector. To isolate the part relevant for perihelion,
introduce the Newtonian specific energy
\begin{equation}
\mathcal E_{\rm N}
\equiv
\frac12 v^2-\frac{\mu}{r},
\label{eq:newtonian-specific-energy}
\end{equation}
which is conserved at leading order. Then
\begin{equation}
\frac18 v^4-\frac12\frac{\mu^2}{r^2}
=
\frac12 \mathcal E_{\rm N}^2
+
\mathcal E_{\rm N}\frac{\mu}{r}.
\label{eq:v4-static-identity}
\end{equation}
The first term in \eqref{eq:v4-static-identity} is a constant, while the second is still
of pure \(1/r\) form. At the order retained, such terms renormalize only the Kepler energy
and effective semimajor-axis bookkeeping; they do not change the closed-orbit character of
a pure Newtonian \(1/r\) problem and therefore do not by themselves generate perihelion
precession.

The precession-relevant \(1\)PN part is thus the velocity-coupled sector, which may be
packaged as a position-dependent inertia:
\begin{equation}
\mathcal L_{\rm orb}
=
\frac12\bigl(1+\sigma(r)\bigr)\bigl(\dot r^2+r^2\dot\phi^2\bigr)
+\frac{\mu}{r},
\label{eq:orbit-lagrangian}
\end{equation}
with
\begin{equation}
\sigma(r)=\frac{\betaPN\mu}{c^2 r}.
\label{eq:sigma-def}
\end{equation}
For the present model,
\begin{equation}
\betaPN=3,
\label{eq:beta-three-repeat}
\end{equation}
so
\begin{equation}
\sigma(r)=\frac{3\mu}{c^2 r}.
\label{eq:sigma-three}
\end{equation}

This makes precise the sense in which the earlier orbital-sector bookkeeping is recovered
from the full \(1\)PN derivation. The self-sector coefficient \(+3/2\) in the pair
Lagrangian becomes, in the test-mass central-field reduction, the effective kinetic
prefactor \(\betaPN/2\), and the already-derived value \(\betaPN=3\) carries the full
precession information.

Because \(\phi\) is cyclic in \eqref{eq:orbit-lagrangian}, the canonical specific angular
momentum is conserved:
\begin{equation}
\ell
\equiv
\frac{\partial \mathcal L_{\rm orb}}{\partial \dot\phi}
=
\bigl(1+\sigma(r)\bigr) r^2 \dot\phi.
\label{eq:orbit-angular-momentum}
\end{equation}
Likewise, the corresponding specific energy is
\begin{equation}
\mathcal E
=
\frac12\bigl(1+\sigma(r)\bigr)\bigl(\dot r^2+r^2\dot\phi^2\bigr)
-\frac{\mu}{r}.
\label{eq:orbit-energy}
\end{equation}

\subsection{Perihelion coefficient}
\label{sec:perihelion}

To derive the orbit equation, set
\begin{equation}
u(\phi)\equiv \frac{1}{r(\phi)},
\qquad
\epsilon \equiv \frac{\betaPN\mu}{c^2}.
\label{eq:u-and-epsilon}
\end{equation}
Then
\begin{equation}
1+\sigma(r)=1+\epsilon u,
\qquad
\dot\phi=\frac{\ell u^2}{1+\epsilon u},
\qquad
\dot r=-\frac{\ell u'}{1+\epsilon u},
\label{eq:orbit-kinematic-identities}
\end{equation}
where a prime denotes \(d/d\phi\).

Substituting \eqref{eq:orbit-kinematic-identities} into the energy integral
\eqref{eq:orbit-energy} yields
\begin{equation}
(u')^2+u^2
=
\frac{2(1+\epsilon u)}{\ell^2}\bigl(\mathcal E+\mu u\bigr).
\label{eq:first-integral-u}
\end{equation}
Expanding to first order in \(c^{-2}\) gives
\begin{equation}
(u')^2+u^2
=
\frac{2\mathcal E}{\ell^2}
+\frac{2\mu}{\ell^2}u
+\frac{2\epsilon\mathcal E}{\ell^2}u
+\frac{2\epsilon\mu}{\ell^2}u^2
+O(c^{-4}).
\label{eq:first-integral-u-expanded}
\end{equation}
Differentiating with respect to \(\phi\) and dividing by \(2u'\) gives the \(1\)PN Binet
equation
\begin{equation}
u''+u
=
\frac{\mu}{\ell^2}
+\frac{\epsilon\mathcal E}{\ell^2}
+\frac{2\epsilon\mu}{\ell^2}u
+O(c^{-4}).
\label{eq:binet-general}
\end{equation}

The constant term
\begin{equation}
\frac{\mu}{\ell^2}+\frac{\epsilon\mathcal E}{\ell^2}
\label{eq:constant-shift}
\end{equation}
only renormalizes the effective semilatus rectum. Define
\begin{equation}
\frac{1}{p_{\rm eff}}
\equiv
\frac{\mu+\epsilon\mathcal E}{\ell^2},
\label{eq:p-eff-def}
\end{equation}
so that \eqref{eq:binet-general} becomes
\begin{equation}
u''+
\left(
1-\frac{2\betaPN\mu^2}{c^2\ell^2}
\right)u
=
\frac{1}{p_{\rm eff}}
+O(c^{-4}).
\label{eq:binet-frequency-shift}
\end{equation}
Thus the solution has the standard slightly detuned Kepler form
\begin{equation}
u(\phi)
=
\frac{1}{p_{\rm eff}}
\left[
1+e\cos(\omega\phi)
\right],
\label{eq:detuned-kepler}
\end{equation}
with
\begin{equation}
\omega^2
=
1-\frac{2\betaPN\mu^2}{c^2\ell^2}.
\label{eq:omega-squared}
\end{equation}
To first post-Newtonian order,
\begin{equation}
\omega
=
1-\frac{\betaPN\mu^2}{c^2\ell^2}
+O(c^{-4}).
\label{eq:omega-expanded}
\end{equation}

A radial cycle therefore corresponds to an azimuthal advance
\begin{equation}
\Phi_r=\frac{2\pi}{\omega},
\end{equation}
so the perihelion shift per orbit is
\begin{equation}
\Delta\phi_{\rm model}
=
\Phi_r-2\pi
=
2\pi\left(\frac{1}{\omega}-1\right)
=
\frac{2\pi\betaPN\mu^2}{c^2\ell^2}
+O(c^{-4}).
\label{eq:deltaphi-in-ell}
\end{equation}
Using the Newtonian relation
\begin{equation}
\ell^2=\mu p=\mu a(1-e^2),
\label{eq:newtonian-ell-p}
\end{equation}
with semimajor axis \(a\) and eccentricity \(e\), this becomes
\begin{equation}
\boxed{
\Delta\phi_{\rm model}
=
\frac{2\pi\betaPN\mu}{c^2 a(1-e^2)}.
}
\label{eq:deltaphi-model}
\end{equation}

Finally, substituting the derived ledger value \(\betaPN=3\) gives
\begin{equation}
\boxed{
\Delta\phi_{\rm model}
=
\frac{6\pi\mu}{c^2 a(1-e^2)}.
}
\label{eq:deltaphi-model-beta3}
\end{equation}
This is exactly the standard Schwarzschild/GR \(1\)PN perihelion target:
\begin{equation}
\Delta\phi_{\rm GR}
=
\frac{6\pi\mu}{c^2 a(1-e^2)}.
\label{eq:deltaphi-gr}
\end{equation}
Hence
\begin{equation}
\Delta\phi_{\rm model}=\Delta\phi_{\rm GR}.
\label{eq:deltaphi-equality}
\end{equation}

\paragraph{Interpretation.}
At the orbit level, the result says that once the full \(1\)PN derivation has fixed
\(\betaPN=3\), the model reproduces the canonical perihelion coefficient without any
further tuning. The observable precession is carried by the same self-sector ledger that
was decomposed earlier into scalar mass dressing, added mass, and internal breathing
response:
\begin{equation}
\betaPN
=
\krho+\kadd+\kPV
=
1+\frac12+\frac32
=
3.
\end{equation}
Thus the perihelion result is not an independent postulate. It is the orbit-level
expression of the full coefficient ledger already derived in the previous sections.

\claimclass{\FullMatch. This section introduces no new fit parameters. Once the full
two-body \(1\)PN result has been reduced to the test-mass orbit sector, the perihelion
shift follows algebraically and matches the standard GR coefficient exactly when
\(\betaPN=3\).}

% ============================================================
\section{Fixed versus symbolic quantities}
\label{sec:ledger-fixed-vs-symbolic}

One of the main purposes of this paper is to make the coefficient status of the toy model
more transparent. Earlier stages of the program necessarily mixed three different kinds of
objects:
\begin{enumerate}[leftmargin=1.5em]
  \item quantities already fixed by exact identities or controlled reductions,
  \item quantities fixed only after a declared response protocol is chosen, and
  \item quantities that should still remain symbolic because they belong to sectors not
  yet solved.
\end{enumerate}
Now that the conservative two-body \(1\)PN assembly has been completed, it is useful to
record that distinction explicitly.

The point of this section is therefore not to weaken the main result. On the contrary, it
sharpens it. Within the closure hierarchy adopted in this paper, the conservative two-body
\(1\)PN sector no longer contains free \emph{dimensionless} coefficients. What remains
symbolic belongs to absolute normalization, higher-order response, finite-size structure,
or sectors beyond the scope of the present \(1\)PN derivation.

\subsection{Quantities fixed by the present derivation chain}
\label{sec:fixed-quantities}

The following quantities should now be regarded as derived carry-forward values rather than
tunable parameters.

\paragraph{Newtonian/scalar sector.}
Newtonian matching of the reduced one-body action fixes the scalar mass-dressing exponent:
\begin{equation}
q=1.
\label{eq:fixed-q}
\end{equation}
Equivalently, in the static/local-mass notation,
\begin{equation}
\krho = 1.
\label{eq:fixed-krho}
\end{equation}
This means that the Newtonian potential couples to the defect monopole with unit strength
in the reduced worldline description.

\paragraph{Optical/barotropic sector.}
Weak-field optical consistency fixes the stiffness exponent of the barotropic closure
\begin{equation}
P(\rho)=\KEOS \rho^n
\end{equation}
to be
\begin{equation}
n=5.
\label{eq:fixed-n}
\end{equation}
This is one of the key ``knob-removal'' results of the program: the exponent \(n\) is no
longer left phenomenological once the optical sector is required to reproduce the weak-field
coefficient.

\paragraph{Free relativistic kinematics.}
The reduced wave-supported particle sector reproduces the universal special-relativistic
quartic term,
\begin{equation}
L_{\rm free}
=
- mc^2+\frac12 mv^2+\frac18 m\frac{v^4}{c^2}+O(c^{-4}),
\label{eq:fixed-free-quartic}
\end{equation}
so the coefficient
\begin{equation}
\frac18
\end{equation}
in the \(1\)PN \emph{Lagrangian} is fixed. Equivalently, the corresponding low-velocity
energy expansion carries the standard
\begin{equation}
\frac38
\end{equation}
coefficient.

\paragraph{Static nonlinear sector.}
The local-mass-scaling route fixes the two-body static coefficient to
\begin{equation}
L_{\rm stat}^{(AB)}
=
-\frac{G^2m_A m_B(m_A+m_B)}{2c^2 r_{AB}^2}.
\label{eq:fixed-static-term}
\end{equation}
Thus the static \(G^2/r^2\) coefficient is fixed at
\begin{equation}
-\frac12.
\end{equation}

\paragraph{Topology/added-mass sector.}
For the actual \(w\)-uniform throat geometry used in the present model, the exterior
potential-flow calculation fixes
\begin{equation}
\kadd=\frac12.
\label{eq:fixed-kadd}
\end{equation}
This should be read as a geometry/topology result, not as a fit parameter. A compact
\(4\)-ball would instead give the counterfactual value
\begin{equation}
\kadd^{(B^4)}=\frac13,
\label{eq:counterfactual-bubble-kadd}
\end{equation}
which is useful as a diagnostic but is \emph{not} the value adopted by the actual throat
model.

\paragraph{Wake/cross sector.}
The constructive isotropic wake calculation, together with the thermodynamic closure tied
to the already-fixed \(n=5\), fixes the conservative parity-even wake parameters to
\begin{equation}
\alpha^2=\frac34,
\qquad
a_H=0,
\qquad
K_{\rm vec}=\frac{2}{\pi^2}.
\label{eq:fixed-wake-parameters}
\end{equation}
These in turn fix the EIH cross coefficients to
\begin{equation}
C_\parallel=-\frac72,
\qquad
C_L=-\frac12.
\label{eq:fixed-cross-coeffs}
\end{equation}
Hence the entire parity-even velocity-dependent cross tensor is fixed.

\paragraph{Full conservative two-body \(1\)PN sector.}
Once the previous ingredients are assembled, the reduced two-body Lagrangian is fixed to
be
\begin{align}
L_{\rm 1PN}^{\rm derived}
={}&
\frac12 m_A v_A^2 + \frac12 m_B v_B^2
+ \frac{m_A v_A^4 + m_B v_B^4}{8c^2}
+ \frac{Gm_A m_B}{r_{AB}}
\nonumber\\
&
+ \frac{Gm_A m_B}{c^2 r_{AB}}
\left[
\frac32\bigl(v_A^2+v_B^2\bigr)
-\frac72\,\vct_A\!\cdot\!\vct_B
-\frac12\,(\vct_A\!\cdot\!\nvec)(\vct_B\!\cdot\!\nvec)
\right]
- \frac{G^2m_A m_B(m_A+m_B)}{2c^2 r_{AB}^2},
\label{eq:fixed-full-eih}
\end{align}
and this equals the standard conservative EIH target exactly:
\begin{equation}
L_{\rm 1PN}^{\rm derived}=L_{\rm 1PN}^{\rm EIH}.
\label{eq:fixed-full-equality}
\end{equation}

Thus, once the adopted closure hierarchy is accepted, the conservative two-body \(1\)PN
sector contains no remaining free dimensionless coefficients.

\subsection{Quantities fixed only within the declared response protocol}
\label{sec:protocol-fixed}

A smaller set of quantities is also fixed in this paper, but only after one explicit
internal-response protocol is declared. These should be treated as \emph{protocol-fixed},
not universally exact.

\paragraph{Breathing/pressure-volume response.}
Under the adiabatic one-degree-of-freedom closure of \cref{sec:kappaPV-closure}, the
internal breathing coefficient is fixed to
\begin{equation}
\kPV=\frac32.
\label{eq:protocol-fixed-kPV}
\end{equation}
This is a genuine derivation \emph{within that closure}, but it should not be interpreted
as an assumption-free theorem of the full parent \(4\)D dynamics.

\paragraph{Self-sector/precession ledger.}
With
\begin{equation}
\krho=1,
\qquad
\kadd=\frac12,
\qquad
\kPV=\frac32,
\end{equation}
the self-sector ledger becomes
\begin{equation}
\betaPN=\krho+\kadd+\kPV=3.
\label{eq:protocol-fixed-beta}
\end{equation}
This is the value that reproduces both the pair self coefficient \(+3/2\) and the
standard perihelion coefficient in the test-mass orbit reduction.

\paragraph{Internal energy partition and breathing slope.}
The same adiabatic closure also fixes auxiliary diagnostic quantities:
\begin{equation}
E_w:E_f:E_{\rm PV}=11:2:5,
\label{eq:protocol-fixed-partition}
\end{equation}
and
\begin{equation}
\frac{d\ln a_*}{d\ln \rho}=-\frac{57}{64}.
\label{eq:protocol-fixed-breathing-slope}
\end{equation}
These are not needed to state the final two-body \(1\)PN Lagrangian, but they are useful
predictions of the chosen closure and therefore belong in the protocol-fixed ledger.

\subsection{Quantities that remain symbolic}
\label{sec:still-symbolic}

What remains symbolic is just as important as what has been fixed. The present paper does
\emph{not} claim that the entire brane-effective theory has been solved. The following
data should still be treated as symbolic or unsolved.

\paragraph{Absolute dimensional normalization.}
The present paper fixes several dimensionless coefficients, but it does not by itself fix
the overall dimensional scale of the barotropic sector. In particular, in
\begin{equation}
P(\rho)=\KEOS \rho^n,
\qquad n=5,
\label{eq:symbolic-KEOS}
\end{equation}
the exponent \(n\) is fixed, but the dimensional constant \(\KEOS\) remains a model
parameter set by normalization and units. Similarly, \(G\) and \(c\) enter the reduced
brane description as measured effective constants; the paper fixes how they appear in the
dimensionless coefficient ledger, not their absolute numerical values from first
principles.

\paragraph{Projection and localization profiles.}
The detailed shapes of the projection kernel and gauge-localization profile,
\begin{equation}
\Wproj(w),
\qquad
\Zloc(w),
\label{eq:symbolic-profiles}
\end{equation}
remain symbolic. The \(1\)PN derivation uses only structural properties of these objects
(normalization, localization, and the existence of the corresponding exact projected
identities), not a unique closed-form choice for their profiles.

\paragraph{Finite-size and higher-multipole data.}
The Newtonian point-particle limit uses the coherent-defect / small-body regime in which
higher multipoles are suppressed. The explicit values of finite-size moments such as
\begin{equation}
Q_{ij}=\int \rho_{\rm def}\,\xi_i\xi_j\,d^3\xi
\label{eq:symbolic-finite-size}
\end{equation}
remain profile-dependent and are not fixed universally here. They would enter at higher
order in \(a/r\) and in finite-size corrections beyond the retained monopole limit.

\paragraph{General internal response beyond the adiabatic closure.}
Outside the specific adiabatic one-degree-of-freedom protocol adopted in
\cref{sec:kappaPV-closure}, the breathing/pressure-volume response should remain symbolic.
In particular, one should expect more general response data such as
\begin{equation}
\kPV \;\to\; \kPV(\omega,\text{protocol}),
\label{eq:symbolic-kPV-general}
\end{equation}
together with possible damping, phase lag, compliance, or multi-mode coupling
coefficients. None of these more general response functions is fixed by the present paper.

\paragraph{Open-system transport beyond the quasi-static reduction.}
The parent \(4\)D theory contains exact leakage and projection-stress channels through
\begin{equation}
\Sleak,
\qquad
R_{ab},
\qquad
\mathbf F_{\partial W},
\label{eq:symbolic-open-system}
\end{equation}
but the conservative \(1\)PN derivation uses a regime in which these are either encoded in
the Poisson hook or neglected at the retained order. Their fully dynamical, non-quasi-static
forms remain symbolic and belong to a more general open-system treatment.

\paragraph{Beyond-\(1\)PN and dissipative sectors.}
The present paper is restricted to the conservative near-zone \(1\)PN problem. The
following sectors are therefore left unresolved:
\begin{equation}
\text{radiation reaction},\quad
\text{dissipative leakage},\quad
\text{spin couplings},\quad
\text{higher-PN terms},\quad
\text{strong-field corrections}.
\label{eq:symbolic-beyond-1pn}
\end{equation}
Likewise, the wake derivation fixes only the minimal real \emph{parity-even conservative}
point; odd-parity or dissipative wake sectors are not claimed to vanish in full generality.

\paragraph{Fully solved moving-throat bulk dynamics.}
Most importantly, the paper does not claim a completed first-principles solution of the
fully dynamical moving defect profile in the parent \(4\)D bulk. The worldtube/particle
limit used here is a controlled reduction, not a theorem obtained by explicitly solving
the full moving-throat PDE problem for arbitrary motion.

\subsection{Compact status ledger}
\label{sec:status-ledger-table}

For convenience, \cref{tab:fixed-symbolic-ledger} summarizes the status of the main
quantities entering the present paper.

\begin{table}[t]
  \centering
  \caption{Status ledger for the main quantities entering the \(1\)PN derivation.}
  \label{tab:fixed-symbolic-ledger}
  \small
  \renewcommand{\arraystretch}{1.15}
  \begin{tabularx}{\linewidth}{@{}L L L@{}}
    \toprule
    Quantity & Status & Comments \\
    \midrule
    \(q\), \(\krho\) & fixed & Newtonian matching gives \(q=\krho=1\) \\
    \(n\) & fixed & optical consistency gives \(n=5\) \\
    \(1/8\) quartic Lagrangian coefficient & fixed & universal relativistic/wave-supported kinematics \\
    Static \(G^2/r^2\) coefficient & fixed & local mass scaling + pair counting gives \(-1/2\) \\
    \(\kadd\) & fixed & throat topology gives \(\kadd=1/2\) \\
    \(\alpha^2\), \(a_H\), \(K_{\rm vec}\) & fixed & wake + thermodynamic closure give \(\alpha^2=3/4\), \(a_H=0\), \(K_{\rm vec}=2/\pi^2\) \\
    \(C_\parallel\), \(C_L\) & fixed & \((-7/2,-1/2)\) from the wake construction \\
    \(L_{\rm 1PN}^{\rm derived}\) & fixed & exact equality to the conservative EIH target \\
    \(\kPV\) & protocol-fixed & equals \(3/2\) in the adiabatic 1-DOF closure \\
    \(\betaPN\) & protocol-fixed & equals \(3\) once \(\kPV=3/2\) is adopted \\
    \(\KEOS\) & symbolic & dimensional EOS normalization remains free \\
    \(\Wproj(w)\), \(\Zloc(w)\) & symbolic & only structural properties are used here \\
    finite-size moments / higher multipoles & symbolic & suppressed in the monopole small-body limit \\
    general response functions & symbolic & \(\kPV(\omega,\text{protocol})\), damping, compliance, etc. \\
    beyond-\(1\)PN and dissipative sectors & symbolic & outside the scope of the present paper \\
    \bottomrule
  \end{tabularx}
\end{table}

\subsection{Interpretation}
\label{sec:interpretation-fixed-symbolic}

The cleanest way to summarize the paper is therefore the following.

\begin{quote}
Within the declared closure hierarchy, the conservative two-body \(1\)PN sector of the
brane-effective theory is fully fixed and matches the standard EIH result exactly. What
remains symbolic belongs to absolute normalization, generalized internal response,
finite-size structure, open-system dynamics beyond the quasi-static reduction, or sectors
beyond conservative \(1\)PN order.
\end{quote}

That is the appropriate strength of the claim. It is stronger than a phenomenological fit,
because the conservative two-body \(1\)PN coefficients are no longer free. It is weaker
than a complete theorem of the fully solved parent \(4\)D dynamics, because some response
and higher-order sectors remain unresolved. The value of the present paper is that it
separates those two statements cleanly.

\claimclass{\Exact for the status classification itself once the earlier derivations are in
place. The fixed entries are the outputs of earlier \Exact, \Controlled, or \FullMatch
steps; the protocol-fixed entries are the outputs of the declared adiabatic closure; the
symbolic entries are outside the solved conservative two-body \(1\)PN sector.}

% ============================================================
\section{Verification and reproducibility ledger}
\label{sec:verification}

The derivation in this paper is analytic, not computational. The Wolfram Language (WL)
scripts accompanying the project are therefore not presented as black-box generators of the
final result. Their role is narrower and more auditable: they act as a referee suite that
checks the nontrivial algebraic identities, coefficient matches, and end-to-end assembly
steps used in the main text. In particular, the scripts are organized so that each
load-bearing statement is tested as an explicit symbolic residual, a named coefficient
comparison, or a manufactured-data self-test.

This distinction matters for interpretation. The scripts verify that the coefficient ledger
derived in the previous sections is internally consistent \emph{within the stated closure
hierarchy}. They do not replace the analytic argument, and they do not upgrade the paper
into an assumption-free theorem of the fully solved moving-throat bulk dynamics. Their
purpose is reproducibility and auditability: to show that every nontrivial step appearing
in the final \(1\)PN assembly can be checked independently.

\subsection{Structure of the referee suite}
\label{sec:verification-scripts}

The project uses three main WL files, each with a distinct role. Their functions are
summarized in \cref{tab:verification}.

\begin{table}[t]
  \centering
  \caption{WL referee suite used to verify the algebra and bookkeeping of the \(1\)PN derivation.}
  \label{tab:verification}
  \small
  \renewcommand{\arraystretch}{1.15}
  \begin{tabularx}{\linewidth}{@{}L L L@{}}
    \toprule
    Script & Role in the derivation & Reported summary \\
    \midrule
    \texttt{4d\_gravity\_1pn\_master\_harness.wl}
    &
    Consolidated gravity-side referee harness: projection normalization, leakage split,
    exact longitudinal identity, Poisson-hook reduction, Newtonian coupling \(q=1\),
    optical coefficient and \(n=5\), added mass, vector-family invariants, adiabatic
    \(\kPV\) closure, \(\betaPN=3\), and orbit/perihelion bookkeeping. Includes
    manufactured-data checks for extraction machinery.
    &
    Reference transcript reports \texttt{PASS count = 60}, \texttt{FAIL count = 0},
    and overall pass.
    \\[4pt]

    \texttt{vector\_wake\_rebuild.wl}
    &
    Self-contained constructive wake module for the parity-even velocity-dependent cross
    tensor. Derives the overlap shapes, solves the EIH ratio constraint, selects a real
    wake, and fixes the overall normalization.
    &
    Reports a real solution with \(a_H=0\), \(a_L=\pm \sqrt{3}/2\),
    \(K_{\rm vec}=2/\pi^2\), reproducing
    \(C_\parallel=-7/2\) and \(C_L=-1/2\).
    \\[4pt]

    \texttt{4d\_full\_1pn\_derivation\_with\_vectorwake.wl}
    &
    End-to-end \(1\)PN assembly script. Imports the wake module, implements the coherent
    defect / small-body worldtube reduction, checks the scalar/self sector, static
    nonlinear term, topology and response ledger, full two-body EIH equality, and
    perihelion reduction.
    &
    Reference transcript reports \texttt{Passes: 37}, \texttt{Fails: 0},
    \texttt{Skips: 0}, with all checks passed.
    \\
    \bottomrule
  \end{tabularx}
\end{table}

The division of labor among these scripts mirrors the logic of the paper itself.
\texttt{vector\_wake\_rebuild.wl} isolates the cross-velocity tensor so that its
\(-7/2\) and \(-1/2\) coefficients are not contaminated by self-sector bookkeeping.
\texttt{4d\_gravity\_1pn\_master\_harness.wl} acts as a broad regression suite for the
gravity-side identities and reduced ledgers. Finally,
\texttt{4d\_full\_1pn\_derivation\_with\_vectorwake.wl} performs the paper-level assembly:
it imports the constructive wake result and verifies that the full reduced two-body
Lagrangian matches the standard conservative EIH target exactly.

\subsection{What is actually being checked}
\label{sec:verification-checked}

The referee suite is designed around the same claim taxonomy used in the main text, and
the checks naturally fall into four groups.

\paragraph{(i) Exact symbolic identities.}
These are checks that do not depend on a further dynamical closure once the definitions are
declared. Representative examples include:
\begin{enumerate}[leftmargin=1.5em]
  \item projection-kernel normalization and the integration-by-parts leakage split,
  \item kinematic identities for projected velocity and current,
  \item the exact longitudinal identity obtained from projected continuity plus Helmholtz
  decomposition,
  \item pair-counted Newtonian interaction bookkeeping,
  \item local-mass-scaling expansion of the static nonlinear sector, and
  \item virial/algebraic identities inside the adiabatic response closure.
\end{enumerate}
These checks are implemented as symbolic residuals reduced to zero under explicitly stated
assumptions.

\paragraph{(ii) Controlled reductions.}
These are checks whose correctness is conditional on the stated regime assumptions rather
than on pure algebra alone. Representative examples include:
\begin{enumerate}[leftmargin=1.5em]
  \item the quasi-static Poisson hook and \(1/r^2\) longitudinal field scaling,
  \item the coherent-defect / small-body worldtube reduction to Newtonian monopole motion,
  \item finite-size suppression at order \((a/r)^2\),
  \item the exterior potential-flow calculation giving \(\kadd=1/2\) for the
  \(w\)-uniform throat and \(1/3\) for the counterfactual \(4\)D bubble,
  \item the optical weak-field coefficient \(\alpha_n=(n-1)/2\) and the resulting
  selection \(n=5\), and
  \item the isotropic Fourier-space wake basis leading to the EIH cross-tensor family.
\end{enumerate}
These checks certify that the reductions used in the paper are algebraically self-consistent
once their regime assumptions are accepted.

\paragraph{(iii) Protocol-level closures.}
The main protocol-dependent ingredient in the present paper is the adiabatic
one-degree-of-freedom breathing closure. The referee suite checks:
\begin{equation}
x=\frac{E_f}{E_w}=\frac{2}{11},
\qquad
\kPV=\frac32,
\qquad
\betaPN=\krho+\kadd+\kPV=3,
\qquad
\frac{d\ln a_*}{d\ln \rho}=-\frac{57}{64}.
\label{eq:verification-protocol-checks}
\end{equation}
These are not presented as exact theorems of the full parent theory; they are documented as
the outputs of the declared adiabatic closure used in this paper.

\paragraph{(iv) Full assembly checks.}
The end-to-end checks are the most important for the present paper. They verify that
\begin{enumerate}[leftmargin=1.5em]
  \item the imported wake module reproduces the EIH cross coefficients,
  \item the scalar/optical route and the \(\betaPN\)-ledger route agree on the same
  self-sector coefficient,
  \item the assembled reduced two-body \(1\)PN Lagrangian matches the standard EIH target
  exactly, and
  \item the test-mass orbit reduction reproduces the canonical perihelion coefficient once
  \(\betaPN=3\).
\end{enumerate}
These are the script-level counterparts of the main analytic claims of
\cref{sec:self-sector,sec:ledger,sec:wake,sec:full-eih,sec:orbit}.

\subsection{Reference run summaries}
\label{sec:verification-summaries}

The project files are written in explicit PASS/FAIL style, so the verification status is
easy to inspect without reading the code line by line. In the reference transcripts
included with the scripts, the broad master harness reports
\begin{equation}
\texttt{PASS count = 60},
\qquad
\texttt{FAIL count = 0},
\label{eq:verification-master-count}
\end{equation}
and the end-to-end \(1\)PN assembly script reports
\begin{equation}
\texttt{Passes: 37},
\qquad
\texttt{Fails: 0},
\qquad
\texttt{Skips: 0}.
\label{eq:verification-assembly-count}
\end{equation}
The wake-rebuild module reports a successful real-parameter solution and prints the selected
parameters explicitly:
\begin{equation}
a_H=0,
\qquad
a_L=\pm\frac{\sqrt3}{2},
\qquad
K_{\rm vec}=\frac{2}{\pi^2},
\label{eq:verification-wake-output}
\end{equation}
followed by the recovered targets
\begin{equation}
C_\parallel=-\frac72,
\qquad
C_L=-\frac12.
\label{eq:verification-wake-targets}
\end{equation}

These reported summaries should be read as reproducibility metadata, not as substitutes for
the derivation itself. Their value is that they provide a compact audit trail: the reader
can see at a glance that the symbolic checks close cleanly, and can then inspect the named
checks in the scripts if more detail is desired.

\subsection{What the referee suite does \emph{not} verify}
\label{sec:verification-scope}

It is equally important to state the scope limits of the verification suite.

First, the scripts do \emph{not} solve the fully dynamical moving-throat PDE problem in the
parent \(4\)D bulk. The particle-limit step is checked in the form used in this paper:
as a coherent-defect / small-body reduction with explicit assumptions on smooth external
fields, vanishing dipole moment, and negligible boundary fluxes at the retained order.

Second, the suite verifies the conservative near-zone \(1\)PN ledger only. It does not
claim to cover dissipative leakage, radiation reaction, spin couplings, higher-order
post-Newtonian terms, or non-adiabatic internal response.

Third, some checks in the master harness are explicitly labeled as \emph{synthetic
self-tests}. These do not derive new physics; rather, they verify the reliability of the
symbolic extraction machinery by applying it to manufactured or exactly solvable data.
Their presence strengthens the reproducibility story, but they should not be confused with
new dynamical claims.

In short, the scripts verify the algebra and bookkeeping of the paper \emph{as stated},
not a stronger claim than the paper itself makes.

\subsection{Reproducibility conventions}
\label{sec:verification-conventions}

Three reproducibility conventions are worth recording explicitly.

\paragraph{Named checks.}
Every nontrivial statement checked by the scripts is given a human-readable name in the
PASS/FAIL output. This makes it possible to compare the paper's prose directly with the
verification transcript.

\paragraph{Local assumptions.}
The scripts do not hide regime assumptions globally. Quasi-staticity, small-body scaling,
adiabatic closure, or parity-even wake restriction are declared in the relevant check
blocks, so the reader can see which equalities are exact and which are conditional.

\paragraph{Modular assembly.}
The constructive wake sector is verified in its own file and then imported into the full
assembly script. This modular structure matters because it separates the genuinely bilinear
cross tensor from the scalar/self bookkeeping and makes double counting easier to rule out.

For readability, the main text reports only the verification ledger and the all-pass
summaries. A fuller PASS/FAIL transcript can be placed in
Appendix~\ref{app:verification} or in a separate repository note accompanying the paper.

\paragraph{Bottom line.}
The purpose of the referee suite is not to ask the reader to trust a computation instead of
a derivation. It is to provide an auditable backstop showing that the coefficient ledger of
the paper closes exactly under the stated assumptions:
\begin{equation}
L_{\rm 1PN}^{\rm derived}=L_{\rm 1PN}^{\rm EIH},
\qquad
\Delta\phi_{\rm model}=\Delta\phi_{\rm GR},
\label{eq:verification-bottom-line}
\end{equation}
with the intermediate coefficients independently checked along the way.

\claimclass{Documentary / reproducibility note. This section does not add a new dynamical
derivation. It records how the accompanying WL referee suite checks the exact identities,
controlled reductions, protocol closures, and full-match assembly steps used elsewhere in
the paper.}

% ============================================================
\section{Conclusions}
\label{sec:conclusions}

The main result of this paper is that the \(4\)D toy model now supports a full
\emph{conservative two-body \(1\)PN derivation} within a clearly stated closure hierarchy.
The point of the paper is not merely that the model can be tuned to resemble the
Einstein--Infeld--Hoffmann (EIH) Lagrangian. Rather, once the derivation is organized by
sector, every coefficient in the conservative two-body \(1\)PN ledger can be traced to a
named ingredient of the parent \(4\)D framework.

The derivation chain established here is:
\begin{enumerate}[leftmargin=1.5em]
  \item exact projection identities in the parent \(4\)D theory, including open-system
  continuity with leakage and the exact longitudinal identity;
  \item a controlled quasi-static near-zone Poisson hook, which supplies inverse-square
  longitudinal behavior on the brane;
  \item a coherent-defect / small-body worldtube reduction, which yields the Newtonian
  point-particle limit;
  \item scalar and optical dressing, which fix \(q=1\) and \(n=5\) and therefore the
  self-sector coefficient;
  \item local mass scaling plus exact pair counting, which generate the static nonlinear
  \(G^2/r^2\) term;
  \item topology and internal response, which give \(\kadd=1/2\),
  \(\kPV=3/2\) under the explicit adiabatic closure, and therefore
  \(\betaPN=3\); and
  \item a constructive isotropic wake calculation, which fixes
  \(\alpha^2=3/4\), \(a_H=0\), \(K_{\rm vec}=2/\pi^2\), and hence the EIH cross
  coefficients \(-7/2\) and \(-1/2\).
\end{enumerate}

When these ingredients are assembled, the reduced two-body Lagrangian is
\begin{align}
L_{\rm 1PN}^{\rm derived}
={}&
\frac12 m_A v_A^2 + \frac12 m_B v_B^2
+ \frac{m_A v_A^4 + m_B v_B^4}{8c^2}
+ \frac{Gm_A m_B}{r_{AB}}
\nonumber\\
&
+ \frac{Gm_A m_B}{c^2 r_{AB}}
\left[
\frac32\bigl(v_A^2+v_B^2\bigr)
-\frac72\,\vct_A\!\cdot\!\vct_B
-\frac12\,(\vct_A\!\cdot\!\nvec)(\vct_B\!\cdot\!\nvec)
\right]
- \frac{G^2m_A m_B(m_A+m_B)}{2c^2 r_{AB}^2},
\label{eq:conclusion-final-eih}
\end{align}
which matches the standard conservative EIH target exactly. In the test-mass orbit
reduction, the same ledger yields
\begin{equation}
\Delta\phi_{\rm model}
=
\frac{2\pi\betaPN\mu}{c^2 a(1-e^2)}
=
\frac{6\pi\mu}{c^2 a(1-e^2)},
\label{eq:conclusion-perihelion}
\end{equation}
so the canonical perihelion coefficient is recovered once \(\betaPN=3\).

Conceptually, the paper does two things at once. First, it upgrades the earlier
coefficient-fixing program into a full conservative \(1\)PN assembly. Second, it makes the
status of each ingredient explicit. Some steps are \Exact statements of the parent theory
or of pair counting. Some are \Controlled reductions, such as the quasi-static Poisson
hook, the coherent-defect particle limit, and the isotropic wake basis. One step is a
\Protocol closure: the adiabatic one-degree-of-freedom computation of \(\kPV\). The final
equality to the standard two-body EIH Lagrangian is then a \FullMatch statement once the
previous ingredients are inserted.

That distinction is the correct way to read the strength of the result. This paper does
\emph{not} claim an assumption-free theorem derived from a fully solved moving-throat bulk
background for arbitrary motion. The coherent-defect particle limit is a controlled
worldtube reduction, and the value \(\kPV=3/2\) is derived only within the declared
adiabatic closure. The paper therefore establishes a full conservative \(1\)PN derivation
\emph{within a stated closure hierarchy}, not a stronger theorem than that.

Within that hierarchy, however, the bookkeeping is now substantially cleaner than in the
earlier stages of the program. The conservative two-body \(1\)PN sector contains no
remaining free dimensionless coefficients. The quantities
\begin{equation}
q=1,
\qquad
n=5,
\qquad
\kadd=\frac12,
\qquad
\alpha^2=\frac34,
\qquad
a_H=0,
\qquad
K_{\rm vec}=\frac{2}{\pi^2},
\qquad
\betaPN=3
\label{eq:conclusion-fixed-ledger}
\end{equation}
are now derived carry-forward values, with \(\kPV=3/2\) fixed inside the declared
adiabatic response protocol. What remains symbolic belongs to a different category:
absolute dimensional normalization, generalized internal response,
finite-size/higher-multipole structure, open-system dynamics beyond the quasi-static
reduction, and sectors beyond conservative \(1\)PN order.

The accompanying WL referee suite is important in this context because it makes the
derivation auditable. The scripts do not replace the analytic argument, but they do verify
that the exact identities, controlled reductions, protocol-level closures, and end-to-end
assembly steps used in the paper close consistently, including the final equality to the
standard EIH target and the perihelion coefficient.

The natural next steps are now clear. The first is to replace the present adiabatic
one-degree-of-freedom breathing closure by a more general internal-response calculation,
so that \(\kPV\) can be derived as a frequency- and protocol-dependent quantity inside the
same parent theory. The second is to push the defect reduction beyond the leading monopole
limit, making finite-size and higher-multipole corrections explicit. The third is to move
beyond the conservative near-zone \(1\)PN sector toward radiation, dissipative leakage,
and higher-order post-Newtonian structure. More ambitiously, one would like to derive the
particle-limit dynamics from a fully solved moving-throat bulk background rather than from
the controlled collective reduction used here.

Those open problems do not undo the main result of the paper. What has been shown is that
the \(4\)D toy model is now capable of reproducing the full conservative two-body \(1\)PN
ledger in a way that is both sector-by-sector transparent and algebraically exact at the
final assembly stage. That is the sense in which the model has progressed from a
suggestive Newtonian/optical framework to a genuine \(1\)PN derivation.

\claimclass{\FullMatch as a summary statement of the paper's central result, with the
important scope condition that the derivation is complete within the stated closure
hierarchy rather than as an assumption-free theorem of the fully solved parent \(4\)D
dynamics.}

% ============================================================
\appendix

\section{Notation and conventions}
\label{app:notation}

This appendix collects the notation and bookkeeping conventions used throughout the paper.
The main text is intentionally organized by derivation sector, so several symbols reappear
in different physical contexts. The purpose of this appendix is to keep those usages
unambiguous and to make clear which operations are exact definitions, which are reductions,
and which are merely naming conventions.

\begin{table}[t]
  \centering
  \caption{Core notation used throughout the paper.}
  \label{tab:notation-core}
  \small
  \renewcommand{\arraystretch}{1.15}
  \begin{tabularx}{\linewidth}{@{}L L@{}}
    \toprule
    Symbol & Meaning \\
    \midrule
    \(X^M=(t,x,y,z,w)\) & \((4+1)\)-dimensional coordinates of the parent theory \\
    \(\X=(x,y,z,w)\in\R^4\) & bulk spatial position \\
    \(\x=(x,y,z)\in\R^3\) & brane spatial position \\
    \(\Proj{Q}\) & brane projection (measurement map) of a bulk quantity \(Q\) \\
    \(\Wproj(w)\) & normalized projection kernel in the extra coordinate \\
    \(\rho=|\psi|^2\) & bulk matter density \\
    \(\PhiN\) & effective Newtonian potential on the brane \\
    \(\phibr\) & longitudinal brane scalar/velocity potential before Newtonian normalization \\
    \(\Sleak\) & leakage source in projected continuity \\
    \(\vct_A,\vct_B\) & ordinary 3D brane velocities of bodies \(A,B\) \\
    \(\mathbf r_{AB}=\mathbf X_A-\mathbf X_B\) & body separation vector \\
    \(r_{AB}=|\mathbf r_{AB}|\), \(\nvec=\mathbf r_{AB}/r_{AB}\) & separation magnitude and unit direction \\
    \(\KEOS\) & equation-of-state constant in \(P(\rho)=\KEOS\rho^n\) \\
    \(n\) & EOS exponent, fixed here to \(5\) \\
    \(\Kvec\) & vector-sector normalization in the wake overlap tensor \\
    \(\krho,\kadd,\kPV\) & scalar mass-dressing, added-mass, and breathing-response coefficients \\
    \(\betaPN\) & self-sector/precession ledger, \(\betaPN=\krho+\kadd+\kPV\) \\
    \(\alphaSq\) & longitudinal wake fraction \\
    \(\Cself,\Cpara,\CLong\) & self, transverse-cross, and longitudinal-cross coefficients \\
    \bottomrule
  \end{tabularx}
\end{table}

\subsection{Coordinates, indices, and derivative conventions}
\label{app:notation-indices}

We work in a parent \((4+1)\)-dimensional spacetime with coordinates
\begin{equation}
X^M=(t,x,y,z,w)\equiv (t,\x,w),
\qquad
\X=(x,y,z,w)\in\R^4,
\qquad
\x=(x,y,z)\in\R^3.
\label{eq:notation-coordinates}
\end{equation}
The index conventions are:
\begin{align}
M,N,\dots &\in \{0,1,2,3,4\}
&& \text{full spacetime indices}, \nonumber\\
\mu,\nu,\dots &\in \{0,1,2,3\}
&& \text{brane spacetime indices}, \nonumber\\
i,j,k,\dots &\in \{x,y,z,w\}
&& \text{bulk spatial indices}, \nonumber\\
a,b,c,\dots &\in \{x,y,z\}
&& \text{brane spatial indices}.
\label{eq:notation-index-ranges}
\end{align}
Repeated spatial indices are summed unless stated otherwise. By contrast, body labels
\(A,B,C,\dots\) are never summed implicitly; sums over bodies are always written out
explicitly.

An overdot denotes differentiation with respect to coordinate time \(t\),
\begin{equation}
\dot f \equiv \frac{df}{dt},
\end{equation}
while in orbit equations a prime denotes differentiation with respect to the azimuthal
angle \(\phi\),
\begin{equation}
u' \equiv \frac{du}{d\phi}.
\end{equation}

The differential operators used in the main text are
\begin{equation}
\gradThree=\nabla_3,
\qquad
\lapThree=\nabla_3^2,
\qquad
\gradFour=\nabla_4,
\qquad
\lapFour=\nabla_4^2.
\label{eq:notation-gradients}
\end{equation}
Unless explicitly indicated otherwise, all dot products such as
\(\vct_A\!\cdot\!\vct_B\) and \(\nvec\!\cdot\!\vct_A\) are ordinary Euclidean 3D brane
dot products.

\subsection{Projection versus reduction}
\label{app:notation-projection}

A central convention of the paper is the distinction between \emph{projection} and
\emph{reduction}.

\paragraph{Projection.}
Projection is the operational rule defining what a brane observer measures. For any bulk
scalar \(Q(t,\x,w)\),
\begin{equation}
\Proj{Q}(t,\x)
\equiv
\int_{-\infty}^{\infty}\Wproj(w)\,Q(t,\x,w)\,\dd w,
\qquad
\int_{-\infty}^{\infty}\Wproj(w)\,\dd w=1.
\label{eq:notation-projection}
\end{equation}
In particular,
\begin{equation}
\Proj{\rho}(t,\x)
=
\int \Wproj(w)\,\rho(t,\x,w)\,\dd w,
\qquad
\mathbf J_{\rm br}(t,\x)
=
\int \Wproj(w)\,\mathbf j_{xyz}(t,\x,w)\,\dd w,
\label{eq:notation-projected-density-current}
\end{equation}
and where \(\Proj{\rho}>0\) the projected brane velocity is defined by
\begin{equation}
\mathbf v_{\rm br}
\equiv
\frac{\mathbf J_{\rm br}}{\Proj{\rho}}.
\label{eq:notation-vbrane}
\end{equation}

\paragraph{Reduction.}
Reduction is a separate operation in which one integrates over \(w\) under an explicit
ansatz or scaling regime in order to obtain an effective lower-dimensional description.
Projection tells us what the brane observer measures; reduction tells us how one may
derive an effective brane theory under additional assumptions. The paper keeps these two
operations distinct throughout.

This distinction matters especially for the continuity law. The projected brane density is
generically an open-system observable:
\begin{equation}
\partial_t\Proj{\rho}
+
\gradThree\!\cdot\mathbf J_{\rm br}
=
\Sleak.
\label{eq:notation-open-continuity}
\end{equation}
Thus leakage is not added by hand at the effective level; it is already implied by the
projection map.

\subsection{Potential, pair, and orbit conventions}
\label{app:notation-potentials}

The Newtonian potential is defined with the usual attractive sign convention,
\begin{equation}
\PhiN(\x)\to 0 \quad \text{as } |\x|\to\infty,
\qquad
\PhiN(\x)=-\frac{GM}{r}
\quad \text{for an isolated point monopole}.
\label{eq:notation-potential-sign}
\end{equation}
For two bodies \(A\) and \(B\),
\begin{equation}
\mathbf r_{AB}=\mathbf X_A-\mathbf X_B,
\qquad
r_{AB}=|\mathbf r_{AB}|,
\qquad
\nvec=\frac{\mathbf r_{AB}}{r_{AB}},
\label{eq:notation-pair-separation}
\end{equation}
and the potential of body \(B\) evaluated at body \(A\) is
\begin{equation}
\Phi_B(\mathbf X_A)=-\frac{Gm_B}{r_{AB}}.
\label{eq:notation-pair-potential}
\end{equation}

The paper uses the standard Lagrangian sign convention
\begin{equation}
L=T-V.
\label{eq:notation-LTV}
\end{equation}
Since the Newtonian pair potential energy is
\begin{equation}
V_{AB}=-\frac{Gm_A m_B}{r_{AB}},
\end{equation}
the interaction term appears in the Lagrangian as
\begin{equation}
L_{\rm int}^{(AB)}=+\frac{Gm_A m_B}{r_{AB}}.
\label{eq:notation-newtonian-pair-term}
\end{equation}
This is the source of the positive Newtonian pair term throughout the paper.

In the test-mass orbit reduction, we write
\begin{equation}
\mu\equiv GM,
\qquad
u(\phi)\equiv \frac{1}{r(\phi)},
\qquad
\ell \ \text{for the specific angular momentum}.
\label{eq:notation-orbit}
\end{equation}
The standard Newtonian relation between semilatus rectum \(p\), semimajor axis \(a\), and
eccentricity \(e\) is
\begin{equation}
p=a(1-e^2),
\qquad
\ell^2=\mu p.
\label{eq:notation-kepler}
\end{equation}

\subsection{Equation-of-state and wake notation}
\label{app:notation-eos-wake}

The barotropic equation of state is written as
\begin{equation}
P(\rho)=\KEOS \rho^n.
\label{eq:notation-eos}
\end{equation}
A recurring source of confusion in earlier drafts was the use of the symbol \(K\) for two
unrelated quantities. In the present paper we reserve:
\begin{equation}
\KEOS
\quad \text{for the equation-of-state normalization},
\qquad
\Kvec
\quad \text{for the dimensionless wake/vector normalization}.
\label{eq:notation-two-Ks}
\end{equation}
These objects should never be identified with one another.

The wake-sector coefficients are written as
\begin{equation}
\alphaSq=\alpha^2,
\qquad
a_H,
\qquad
\Kvec,
\qquad
\Cpara,
\qquad
\CLong,
\label{eq:notation-wake-symbols}
\end{equation}
where:
\begin{itemize}[leftmargin=1.5em]
  \item \(\alphaSq\) is the longitudinal wake fraction,
  \item \(a_H\) is the parity-even helical amplitude,
  \item \(\Kvec\) is the overall dimensionless normalization of the cross tensor,
  \item \(\Cpara\) multiplies \(\vct_A\!\cdot\!\vct_B\),
  \item \(\CLong\) multiplies
  \((\vct_A\!\cdot\!\nvec)(\vct_B\!\cdot\!\nvec)\).
\end{itemize}
The final derived values are
\begin{equation}
\alphaSq=\frac34,
\qquad
a_H=0,
\qquad
\Kvec=\frac{2}{\pi^2},
\qquad
\Cpara=-\frac72,
\qquad
\CLong=-\frac12.
\label{eq:notation-wake-final-values}
\end{equation}

\subsection{1PN bookkeeping conventions}
\label{app:notation-pn}

We keep \(c\) and \(G\) explicit throughout; no geometric-unit convention
\((c=G=1)\) is adopted. The weak-field/slow-motion expansion parameters are
\begin{equation}
\epsilon_\Phi \equiv \frac{\PhiN}{c^2},
\qquad
\epsilon_v^2 \equiv \frac{v^2}{c^2},
\label{eq:notation-epsilons}
\end{equation}
and both are treated as order \(c^{-2}\) in the standard post-Newtonian counting.
Accordingly, the conservative \(1\)PN sector means:
\begin{equation}
\text{retain all terms through } O(c^{-2}) \text{ in the Lagrangian,}
\qquad
\text{drop } O(c^{-4}) \text{ and higher.}
\label{eq:notation-pn-counting}
\end{equation}

The coefficient
\begin{equation}
\betaPN \equiv \krho+\kadd+\kPV
\label{eq:notation-betaPN}
\end{equation}
is a paper-specific ledger quantity summarizing the self-sector/precession contribution.
It should not be confused with the standard PPN parameter \(\beta\). In the present paper,
\begin{equation}
\krho=1,
\qquad
\kadd=\frac12,
\qquad
\kPV=\frac32
\quad \text{(within the adiabatic closure)},
\qquad
\betaPN=3.
\label{eq:notation-betaPN-values}
\end{equation}

Likewise, the self-sector coefficient \(\Cself\) is the one-body coefficient multiplying
\(\PhiN v^2/c^2\) before conversion to the pair Lagrangian:
\begin{equation}
\Cself=-\frac32,
\label{eq:notation-Cself}
\end{equation}
so the corresponding pair coefficient becomes \(+\frac32\) after inserting the negative
Newtonian potential.

\subsection{Claim-class labels}
\label{app:notation-claimclasses}

The paper uses four claim-class labels to distinguish exact identities from statements
that depend on explicit approximations or closures:
\begin{description}[leftmargin=1.8em,style=nextline]
  \item[\Exact] identities that follow directly from definitions, the parent action,
  projection algebra, or exact pair counting.
  \item[\Controlled] statements that require an explicit regime assumption, such as
  quasi-staticity, compact-defect / small-body scaling, far-field wake closure, or
  weak-field truncation.
  \item[\Protocol] quantities that depend on a declared internal-response model, most
  importantly the adiabatic one-degree-of-freedom computation of \(\kPV\).
  \item[\FullMatch] end-to-end equalities that become exact once all previously derived
  ingredients are inserted, such as the final equality to the standard conservative
  two-body EIH Lagrangian.
\end{description}

These labels are part of the paper's scientific bookkeeping, not just its exposition.
They mark which statements are exact consequences of the parent framework and which depend
on the closure hierarchy adopted for the present \(1\)PN derivation.

\paragraph{Summary.}
The three conventions most worth keeping in mind while reading the paper are:
\begin{enumerate}[leftmargin=1.5em]
  \item projection and reduction are different operations,
  \item the Newtonian potential \(\PhiN\) is negative for attraction, so the Newtonian
  pair term in the Lagrangian is positive, and
  \item \(\KEOS\) and \(\Kvec\) are unrelated quantities despite the historical reuse of
  the letter \(K\) in earlier drafts.
\end{enumerate}
Those three points account for most of the notational ambiguities that can otherwise arise
in the derivation.

\section{Minimal variational formulas inherited from the 4D model}
\label{app:variational}

This appendix records only the action-side formulas from the parent \(4\)D model that are
actually used, directly or indirectly, in the present \(1\)PN derivation. The goal is not
to reproduce the full parent-paper development, but to keep the provenance of the main
text explicit: the continuity identity, the localized Maxwell equation, the barotropic
closure, and the geometry-response sector all descend from a single variational
specification.

The guiding principle is the same as in the main text. Exact identities are kept separate
from reductions. The formulas below are therefore organized into:
\begin{enumerate}[leftmargin=1.5em]
  \item the parent action and its three sectors,
  \item the exact Euler--Lagrange equations inherited from those sectors, and
  \item the minimal geometry-force formulas carried into the adiabatic-response ledger.
\end{enumerate}

\subsection{Parent action: matter, localized Maxwell, and geometry}
\label{app:variational-action}

The minimal parent model is written as
\begin{equation}
S = S_\psi + S_{\rm EM} + S_{\rm geom},
\label{eq:app-total-action}
\end{equation}
with
\begin{equation}
S_\psi = \int dt \int d^4X\, \mathcal L_\psi,
\qquad
S_{\rm EM} = \int dt \int d^4X\, \mathcal L_{\rm EM},
\qquad
S_{\rm geom} = \int dt\, \mathcal L_{\rm geom}.
\label{eq:app-sector-actions}
\end{equation}
Here \(d^4X\equiv dx\,dy\,dz\,dw\), and the bulk spatial coordinate is
\(\mathbf X=(x,y,z,w)\in\mathbb R^4\).

\paragraph{Matter sector.}
The bulk matter field is a complex order parameter \(\psi(\mathbf X,t)\) with density
\begin{equation}
\rho(\mathbf X,t)\equiv |\psi(\mathbf X,t)|^2.
\label{eq:app-density}
\end{equation}
We use minimal gauge coupling
\begin{equation}
D_t\psi = \partial_t\psi + \frac{i q}{\hbar}A_0\psi,
\qquad
D_i\psi = \partial_i\psi - \frac{i q}{\hbar}A_i\psi,
\qquad
i\in\{x,y,z,w\},
\label{eq:app-cov-derivs}
\end{equation}
and the matter Lagrangian density
\begin{equation}
\boxed{
\mathcal L_\psi
=
\frac{i\hbar}{2}\bigl(\psi^*D_t\psi-\psi\,D_t\psi^*\bigr)
-\frac{\hbar^2}{2m_\psi}(D_i\psi)^*(D_i\psi)
- V_{\rm conf}(\mathbf X;a,L)\,\rho
- U(\rho).
}
\label{eq:app-Lpsi}
\end{equation}
The confinement potential \(V_{\rm conf}\) encodes the localized throat/tube geometry and
depends on the collective geometry variables \(a(t)\) and \(L(t)\). The parameter
\(m_\psi\) is the bulk matter mass scale entering the parent GNLS sector; it is unrelated
to the point-particle masses \(m_A,m_B\) appearing later in the reduced two-body problem.

\paragraph{Barotropic closure.}
The matter self-interaction is encoded through a barotropic equation of state
\begin{equation}
P(\rho)=\KEOS \rho^n.
\label{eq:app-eos}
\end{equation}
The corresponding internal-energy density is chosen so that
\begin{equation}
P(\rho)=\rho\,h(\rho)-U(\rho),
\qquad
h(\rho)\equiv \frac{dU}{d\rho},
\label{eq:app-thermo-identity}
\end{equation}
which gives
\begin{equation}
U(\rho)=\frac{\KEOS}{n-1}\rho^n,
\qquad
h(\rho)=\frac{n\KEOS}{n-1}\rho^{n-1}.
\label{eq:app-U-h-general}
\end{equation}
In the present paper the optical sector fixes
\begin{equation}
n=5,
\end{equation}
so the carried-forward specialization is
\begin{equation}
U(\rho)=\frac{\KEOS}{4}\rho^5,
\qquad
h(\rho)=\frac{5\KEOS}{4}\rho^4.
\label{eq:app-U-h-n5}
\end{equation}

\paragraph{Localized Maxwell sector.}
The gauge sector is weighted by a localization profile \(\Zloc(w)\), so that
electromagnetic-type fields are brane-dominant while remaining variationally defined in
the bulk. With
\begin{equation}
F_{MN}\equiv \partial_M A_N-\partial_N A_M,
\qquad
\partial\!\cdot\!A\equiv \partial_M A^M,
\label{eq:app-FMN}
\end{equation}
the minimal gauge-field Lagrangian density is
\begin{equation}
\boxed{
\mathcal L_{\rm EM}
=
-\frac{\Zloc(w)}{4\mu_0}F_{MN}F^{MN}
-\frac{1}{2\xi\mu_0}\,(\partial\!\cdot\!A)^2
+ A_M J_{\rm ext}^M.
}
\label{eq:app-LEM}
\end{equation}
The explicit source \(J_{\rm ext}^M\) is reserved for prescribed external/background
contributions. The matter current is generated automatically by varying the covariantized
matter action and is therefore not duplicated inside \(\mathcal L_{\rm EM}\).

\paragraph{Geometry sector.}
The localized throat is represented by two collective coordinates,
\begin{equation}
a(t)\quad \text{(effective radius)},
\qquad
L(t)\quad \text{(effective axial extent)}.
\label{eq:app-a-L}
\end{equation}
A minimal conservative geometry Lagrangian is
\begin{equation}
\boxed{
\mathcal L_{\rm geom}
=
\frac12 M_a \dot a^2
+
\frac12 M_L \dot L^2
-
E_{\rm geom}(a,L)
-
E_{\rm ratio}(a,L),
}
\label{eq:app-Lgeom}
\end{equation}
with optional ratio penalty
\begin{equation}
E_{\rm ratio}(a,L)
=
\frac{\kappa}{2}\bigl(L-\Lambda a\bigr)^2.
\label{eq:app-Eratio}
\end{equation}
The ratio-penalty term is optional and may be omitted by setting \(\kappa=0\). It is
included here because it is the minimal way of encoding the possibility that the throat
prefers a quasi-fixed aspect ratio in slow breathing configurations.

\subsection{Exact Euler--Lagrange equations carried into the 1PN story}
\label{app:variational-eom}

Only a few exact Euler--Lagrange consequences of the parent action are used in the present
paper.

\paragraph{Matter equation.}
Varying \(S_\psi\) with respect to \(\psi^*\) gives the gauged \(4\)-spatial-dimensional
GNLS equation
\begin{equation}
\boxed{
i\hbar D_t\psi
=
\left[
-\frac{\hbar^2}{2m_\psi}D_iD_i
+
V_{\rm conf}(\mathbf X;a,L)
+
h(\rho)
\right]\psi.
}
\label{eq:app-gnls}
\end{equation}
This is the exact bulk matter evolution law inherited by all later reductions.

\paragraph{Matter current and continuity.}
Global \(U(1)\) phase invariance implies the exact number current
\begin{equation}
j^i
=
\frac{\hbar}{m_\psi}\,\Im\!\bigl(\psi^*D_i\psi\bigr),
\label{eq:app-current}
\end{equation}
and therefore the exact bulk continuity identity
\begin{equation}
\boxed{
\partial_t\rho+\partial_i j^i=0.
}
\label{eq:app-continuity}
\end{equation}
The gauge-coupled matter current entering Maxwell's equations is
\begin{equation}
J_\psi^0 = q\rho,
\qquad
J_\psi^i = q\,j^i.
\label{eq:app-gauge-current}
\end{equation}
Equation \eqref{eq:app-continuity} is the exact input that later produces the projected
brane continuity law and, after Helmholtz decomposition, the exact longitudinal identity.

\paragraph{Localized Maxwell equation.}
Varying \(S_{\rm EM}+S_\psi\) with respect to \(A_N\) gives
\begin{equation}
\boxed{
\partial_M\!\left(\Zloc(w)F^{MN}\right)
+
\frac{1}{\xi}\partial^N(\partial\!\cdot\!A)
=
\mu_0\bigl(J_\psi^N+J_{\rm ext}^N\bigr).
}
\label{eq:app-maxwell}
\end{equation}
This is the exact parent equation behind the controlled zero-mode reduction used in the
main text. Under the brane-compatible zero-mode ansatz
\begin{equation}
A_w=0,
\qquad
\partial_w A_\mu = 0,
\qquad
\mu\in\{0,1,2,3\},
\label{eq:app-zero-mode}
\end{equation}
integration over \(w\) yields the effective brane Maxwell equation quoted earlier,
\begin{equation}
\partial_\mu F^{\mu\nu}
=
\mu_0^{\rm eff}J_{\rm eff}^\nu,
\qquad
\mu_0^{\rm eff}
=
\frac{\mu_0}{Z_{\rm int}},
\qquad
Z_{\rm int}\equiv \int \Zloc(w)\,dw.
\label{eq:app-brane-maxwell}
\end{equation}
Equation \eqref{eq:app-brane-maxwell} is a controlled reduction; the variational formula
that underlies it is the exact bulk equation \eqref{eq:app-maxwell}.

\subsection{Minimal thermodynamic formulas used later}
\label{app:variational-thermo}

The main text repeatedly uses the same small collection of thermodynamic identities, so it
is useful to gather them in one place. From \eqref{eq:app-U-h-general} we obtain
\begin{equation}
P(\rho)=\KEOS \rho^n,
\qquad
U(\rho)=\frac{\KEOS}{n-1}\rho^n,
\qquad
h(\rho)=\frac{n\KEOS}{n-1}\rho^{n-1}.
\label{eq:app-thermo-triple}
\end{equation}
The characteristic barotropic propagation speed is then
\begin{equation}
c_s^2(\rho)=\frac{dP}{d\rho}=n\KEOS \rho^{n-1},
\label{eq:app-cs2}
\end{equation}
in the normalization used throughout the present paper. Linearizing about a background
density \(\rho_0\) with
\begin{equation}
\delta \equiv \frac{\rho-\rho_0}{\rho_0},
\label{eq:app-delta}
\end{equation}
gives
\begin{equation}
\frac{\Delta c_s}{c_s}
\simeq
\frac{n-1}{2}\,\delta,
\qquad
N(\rho)\equiv \frac{c}{c_s(\rho)}
\simeq
1-\frac{n-1}{2}\,\delta.
\label{eq:app-linearized-index}
\end{equation}
These are the formulas used in the main text to connect the optical weak-field coefficient
to the EOS exponent \(n\), and therefore to fix \(n=5\).

\subsection{Geometry forces and the adiabatic-response bridge}
\label{app:variational-geometry}

The parent action is also the source of the geometry-force formulas used later in the
adiabatic-response ledger. Define the conservative total energy
\begin{equation}
H_{\rm tot}
=
H_\psi[\psi;a,L]
+
H_{\rm EM}[A;a,L]
+
E_{\rm geom}(a,L)
+
E_{\rm ratio}(a,L),
\label{eq:app-Htot}
\end{equation}
where the \(a,L\) dependence enters both through the explicit geometry sector and through
the confinement and localization structure seen by the fields. The conservative geometry
forces are then defined by
\begin{equation}
F_a \equiv -\frac{\partial H_{\rm tot}}{\partial a},
\qquad
F_L \equiv -\frac{\partial H_{\rm tot}}{\partial L}.
\label{eq:app-FaFL}
\end{equation}
In particular, the matter-sector contributions take the Hellmann--Feynman form
\begin{equation}
-\frac{\partial H_\psi}{\partial a}
=
-\int d^4X\, \rho(\mathbf X,t)\,\partial_a V_{\rm conf}(\mathbf X;a,L),
\qquad
-\frac{\partial H_\psi}{\partial L}
=
-\int d^4X\, \rho(\mathbf X,t)\,\partial_L V_{\rm conf}(\mathbf X;a,L).
\label{eq:app-HF-matter}
\end{equation}

At the purely conservative variational level, the collective coordinates satisfy
\begin{equation}
M_a \ddot a = -\frac{\partial H_{\rm tot}}{\partial a},
\qquad
M_L \ddot L = -\frac{\partial H_{\rm tot}}{\partial L}.
\label{eq:app-geom-conservative}
\end{equation}
The main text later adopts the familiar damped extension
\begin{equation}
M_a \ddot a + \Gamma_a \dot a = -\frac{\partial H_{\rm tot}}{\partial a},
\qquad
M_L \ddot L + \Gamma_L \dot L = -\frac{\partial H_{\rm tot}}{\partial L},
\label{eq:app-geom-damped}
\end{equation}
as a baseline response closure. Equation \eqref{eq:app-geom-damped} is not itself an
exact Euler--Lagrange consequence of the conservative action; it is the phenomenological
relaxation/ringing extension used when the geometry is treated as a slowly responding open
subsystem.

This distinction is important for the later \(1\)PN bookkeeping. The adiabatic breathing
coefficient \(\kPV\) is not read directly off the conservative action alone. It is
computed only after one additional response protocol is chosen. The variational formulas
above provide the conservative energy landscape on which that protocol is imposed.

\subsection{What is imported into the present paper}
\label{app:variational-imported}

Only a narrow subset of the parent variational structure is actually used in the present
\(1\)PN derivation:
\begin{enumerate}[leftmargin=1.5em]
  \item the matter action \eqref{eq:app-Lpsi} and its exact continuity identity
  \eqref{eq:app-continuity},
  \item the localized Maxwell equation \eqref{eq:app-maxwell},
  \item the barotropic identities \eqref{eq:app-thermo-triple}--\eqref{eq:app-linearized-index},
  \item the existence of the collective geometry coordinates \(a(t)\) and \(L(t)\), and
  \item the conservative geometry-force definitions \eqref{eq:app-FaFL}.
\end{enumerate}
Everything else in the main text---projected continuity, the longitudinal Poisson hook,
the coherent-defect particle limit, the \(1\)PN self/cross/static sectors, and the final
EIH match---is obtained by combining these inherited variational formulas with the
explicit reductions and closures stated in the body of the paper.

\claimclass{\Exact for the action, the GNLS equation, the matter current, the bulk
continuity identity, the localized Maxwell equation, and the conservative geometry-force
definitions. \Controlled for the zero-mode brane Maxwell reduction
\eqref{eq:app-brane-maxwell}. \Protocol for any damped or adiabatic closure built on top
of the conservative geometry sector, including the later computation of \(\kPV\).}

\section{Moment integrals and coherent-defect reduction details}
\label{app:moments}

This appendix supplies the moment algebra underlying the coherent-defect / small-body
reduction used in \cref{sec:newtonian-limit}. The main text stated the reduction in its
most economical form: a compact defect moving in a smooth external field behaves, at
leading order, like a Newtonian point particle, with finite-size effects suppressed by
powers of \(a/r\). Here we record the explicit moment formulas that justify that claim and
make clear where the leading corrections enter.

The discussion has four parts:
\begin{enumerate}[leftmargin=1.5em]
  \item defect-centered coordinates and the low moments of a localized profile,
  \item the center-of-mass split of momentum and kinetic energy,
  \item the multipole expansion of the external force, and
  \item the size of the leading finite-size corrections.
\end{enumerate}
Nothing in this appendix changes the logic of the main text. Its purpose is only to make
the reduction auditable.

\subsection{Defect-centered coordinates and profile moments}
\label{app:moments-profile}

Let \(W_t\subset \mathbb R^3\) denote the compact spatial support of the brane-facing
defect at time \(t\), and let \(\mathbf X_{\rm cm}(t)\) be its center-of-mass position.
Introduce defect-centered coordinates
\begin{equation}
\boldsymbol{\xi}\equiv \x-\mathbf X_{\rm cm}(t),
\qquad
\x=\mathbf X_{\rm cm}(t)+\boldsymbol{\xi}.
\label{eq:app-xi-def}
\end{equation}
The defect density is written as
\begin{equation}
\rho_{\rm def}(\x,t)=\rho_*(\boldsymbol{\xi},t),
\label{eq:app-rho-star}
\end{equation}
with \(\rho_*\) strongly localized near \(\boldsymbol{\xi}=0\).

The total defect mass is
\begin{equation}
M\equiv \int_{W_t}\rho_{\rm def}(\boldsymbol{\xi},t)\,\dd^3\xi.
\label{eq:app-M-def}
\end{equation}
The center-of-mass convention is that the dipole vanishes:
\begin{equation}
\int_{W_t}\rho_{\rm def}(\boldsymbol{\xi},t)\,\xi_i\,\dd^3\xi=0.
\label{eq:app-dipole-zero}
\end{equation}
The second and fourth moments are
\begin{align}
Q_{ij}
&\equiv
\int_{W_t}\rho_{\rm def}(\boldsymbol{\xi},t)\,\xi_i\xi_j\,\dd^3\xi,
\label{eq:app-Qij}
\\[3pt]
Q_{ijkl}
&\equiv
\int_{W_t}\rho_{\rm def}(\boldsymbol{\xi},t)\,\xi_i\xi_j\xi_k\xi_l\,\dd^3\xi.
\label{eq:app-Qijkl}
\end{align}
For a general smooth localized profile, these tensors encode the leading finite-size
corrections to the monopole motion.

When the defect is isotropic to leading order, rotational symmetry implies
\begin{equation}
Q_{ij}=q_2\,\delta_{ij},
\label{eq:app-Qij-isotropic}
\end{equation}
and
\begin{equation}
Q_{ijkl}
=
q_4\,
\bigl(
\delta_{ij}\delta_{kl}
+
\delta_{ik}\delta_{jl}
+
\delta_{il}\delta_{jk}
\bigr),
\label{eq:app-Qijkl-isotropic}
\end{equation}
for scalar coefficients \(q_2,q_4\).

\subsection{Gaussian representative closure}
\label{app:moments-gaussian}

Although the main reduction does not require a Gaussian profile, it is useful to have one
explicit profile for which all moments are elementary. We therefore use the isotropic
Gaussian
\begin{equation}
\rho_{\rm def}(\boldsymbol{\xi})
=
\frac{M}{\pi^{3/2}a^3}
\exp\!\left(-\frac{\xi^2+\eta^2+\zeta^2}{a^2}\right),
\label{eq:app-gaussian-profile}
\end{equation}
where \(a\) sets the defect size.

The normalization follows from the standard one-dimensional Gaussian integrals
\begin{equation}
\int_{-\infty}^{\infty}e^{-x^2/a^2}\,\dd x=\sqrt{\pi}\,a,
\qquad
\int_{-\infty}^{\infty}x^2 e^{-x^2/a^2}\,\dd x=\frac{\sqrt{\pi}}{2}\,a^3,
\qquad
\int_{-\infty}^{\infty}x^4 e^{-x^2/a^2}\,\dd x=\frac{3\sqrt{\pi}}{4}\,a^5.
\label{eq:app-gaussian-1d}
\end{equation}
Using factorization in Cartesian coordinates,
\begin{equation}
\int_{\mathbb R^3}\rho_{\rm def}\,\dd^3\xi=M.
\label{eq:app-gaussian-normalized}
\end{equation}
Odd moments vanish identically by parity, so in particular
\begin{equation}
\int \rho_{\rm def}\,\xi_i\,\dd^3\xi=0.
\label{eq:app-gaussian-dipole}
\end{equation}

The second moments are
\begin{equation}
Q_{ij}
=
\int \rho_{\rm def}\,\xi_i\xi_j\,\dd^3\xi
=
\frac{Ma^2}{2}\,\delta_{ij}.
\label{eq:app-gaussian-Qij}
\end{equation}
Thus the scalar second-moment coefficient is
\begin{equation}
q_2=\frac{Ma^2}{2}.
\label{eq:app-q2}
\end{equation}
Similarly, the fourth moment is
\begin{equation}
Q_{ijkl}
=
\frac{Ma^4}{4}
\bigl(
\delta_{ij}\delta_{kl}
+
\delta_{ik}\delta_{jl}
+
\delta_{il}\delta_{jk}
\bigr),
\label{eq:app-gaussian-Qijkl}
\end{equation}
so
\begin{equation}
q_4=\frac{Ma^4}{4}.
\label{eq:app-q4}
\end{equation}

It is sometimes useful to record the rotationally invariant contractions:
\begin{equation}
\int \rho_{\rm def}\,\xi^2\,\dd^3\xi
=
Q_{ii}
=
\frac{3}{2}Ma^2,
\label{eq:app-xisq-mean}
\end{equation}
and
\begin{equation}
\int \rho_{\rm def}\,\xi^4\,\dd^3\xi
=
Q_{iijj}
=
\frac{15}{4}Ma^4.
\label{eq:app-xi4-mean}
\end{equation}
These formulas are convenient for estimating the magnitude of omitted multipole terms.

\subsection{Center-of-mass momentum and kinetic-energy split}
\label{app:moments-kinematics}

Decompose the local defect velocity as
\begin{equation}
\vct(\boldsymbol{\xi},t)
=
\dot{\mathbf X}_{\rm cm}(t)
+
\mathbf u_{\rm int}(\boldsymbol{\xi},t),
\label{eq:app-v-split}
\end{equation}
with the coherent-translation closure
\begin{equation}
\int_{W_t}\rho_{\rm def}\,\mathbf u_{\rm int}\,\dd^3\xi=0.
\label{eq:app-uint-zero}
\end{equation}
Then the total momentum is
\begin{align}
\mathbf P
&=
\int_{W_t}\rho_{\rm def}\,\vct\,\dd^3\xi
\nonumber\\
&=
\dot{\mathbf X}_{\rm cm}\int_{W_t}\rho_{\rm def}\,\dd^3\xi
+
\int_{W_t}\rho_{\rm def}\,\mathbf u_{\rm int}\,\dd^3\xi
\nonumber\\
&=
M\dot{\mathbf X}_{\rm cm}.
\label{eq:app-total-momentum}
\end{align}
Thus the center-of-mass momentum is exactly monopolar under the coherent-defect closure.

The kinetic energy is
\begin{align}
T
&=
\frac12\int_{W_t}\rho_{\rm def}\,|\vct|^2\,\dd^3\xi
\nonumber\\
&=
\frac12\int_{W_t}\rho_{\rm def}\,
\bigl(
|\dot{\mathbf X}_{\rm cm}|^2
+
2\dot{\mathbf X}_{\rm cm}\!\cdot\!\mathbf u_{\rm int}
+
|\mathbf u_{\rm int}|^2
\bigr)\dd^3\xi
\nonumber\\
&=
\frac12 M\dot{\mathbf X}_{\rm cm}^{\,2}
+
\frac12\int_{W_t}\rho_{\rm def}\,|\mathbf u_{\rm int}|^2\,\dd^3\xi,
\label{eq:app-kinetic-split}
\end{align}
since the mixed term vanishes by \eqref{eq:app-uint-zero}. The first term is the ordinary
center-of-mass kinetic energy, while the second is an internal contribution. At Newtonian
order, the internal term renormalizes only the defect rest-energy ledger and does not
alter the leading monopole equation of motion.

\subsection{External-force multipole expansion}
\label{app:moments-force}

Let the defect move in a smooth external Newtonian potential \(\PhiN(\x,t)\). The total
external force is
\begin{equation}
\mathbf F_{\rm ext}
=
-\int_{W_t}\rho_{\rm def}(\boldsymbol{\xi},t)\,
\gradThree \PhiN\!\bigl(\mathbf X_{\rm cm}(t)+\boldsymbol{\xi},t\bigr)\,\dd^3\xi
+
\mathbf F_{\partial W},
\label{eq:app-force-def}
\end{equation}
where \(\mathbf F_{\partial W}\) collects fluxes through the worldtube boundary. In the
small-body conservative reduction used in the main text, these boundary terms are assumed
subleading and are dropped at the retained order.

Taylor-expand the external field about the center:
\begin{equation}
\partial_i \PhiN(\mathbf X_{\rm cm}+\boldsymbol{\xi})
=
\partial_i\PhiN
+
\xi_j\,\partial_j\partial_i\PhiN
+
\frac12\xi_j\xi_k\,\partial_j\partial_k\partial_i\PhiN
+
\frac{1}{6}\xi_j\xi_k\xi_l\,\partial_j\partial_k\partial_l\partial_i\PhiN
+\cdots,
\label{eq:app-grad-taylor}
\end{equation}
with all derivatives on the right evaluated at \(\mathbf X_{\rm cm}(t)\).

Substituting \eqref{eq:app-grad-taylor} into \eqref{eq:app-force-def} gives
\begin{align}
F_i
={}&
-\partial_i\PhiN \int \rho_{\rm def}\,\dd^3\xi
-\partial_j\partial_i\PhiN \int \rho_{\rm def}\,\xi_j\,\dd^3\xi
\nonumber\\
&
-\frac12\partial_j\partial_k\partial_i\PhiN
\int \rho_{\rm def}\,\xi_j\xi_k\,\dd^3\xi
-\frac{1}{6}\partial_j\partial_k\partial_l\partial_i\PhiN
\int \rho_{\rm def}\,\xi_j\xi_k\xi_l\,\dd^3\xi
+\cdots
+
F_{\partial W,i}.
\label{eq:app-force-expanded-raw}
\end{align}
Using the mass and dipole definitions,
\begin{equation}
F_i
=
- M\,\partial_i\PhiN
-\frac12\,Q_{jk}\,\partial_j\partial_k\partial_i\PhiN
-\frac{1}{6}\,Q_{jkl}^{(3)}\,\partial_j\partial_k\partial_l\partial_i\PhiN
+\cdots
+
F_{\partial W,i},
\label{eq:app-force-multipole}
\end{equation}
where
\begin{equation}
Q_{jkl}^{(3)}
\equiv
\int \rho_{\rm def}\,\xi_j\xi_k\xi_l\,\dd^3\xi
\label{eq:app-third-moment}
\end{equation}
is the third moment. For profiles with inversion symmetry, such as the Gaussian closure,
all odd moments vanish, so the first correction after the monopole is quadrupolar:
\begin{equation}
F_i
=
- M\,\partial_i\PhiN
-\frac12\,Q_{jk}\,\partial_j\partial_k\partial_i\PhiN
+O(Q_{ijkl}\nabla^5\PhiN)
+
F_{\partial W,i}.
\label{eq:app-force-quadrupole}
\end{equation}

For an isotropic profile with \(Q_{jk}=q_2\delta_{jk}\), this simplifies to
\begin{equation}
F_i
=
- M\,\partial_i\PhiN
-\frac{q_2}{2}\,\partial_i\nabla^2\PhiN
+O(q_4\nabla^5\PhiN)
+
F_{\partial W,i}.
\label{eq:app-force-isotropic}
\end{equation}
For the Gaussian representative \eqref{eq:app-gaussian-Qij},
\begin{equation}
F_i
=
- M\,\partial_i\PhiN
-\frac{Ma^2}{4}\,\partial_i\nabla^2\PhiN
+O(Ma^4\nabla^5\PhiN)
+
F_{\partial W,i}.
\label{eq:app-force-gaussian}
\end{equation}

Equation \eqref{eq:app-force-isotropic} makes one subtle point transparent. In a vacuum
region where the external potential is harmonic,
\begin{equation}
\nabla^2\PhiN=0,
\label{eq:app-harmonic-vacuum}
\end{equation}
the isotropic quadrupole correction vanishes identically. In that case, the first
finite-size correction to the monopole force is pushed to still higher order. This is one
reason the point-particle reduction is especially clean in the exterior two-body near
zone.

\subsection{Suppression of finite-size corrections}
\label{app:moments-suppression}

Suppose the external field is sourced by a distant compact body of mass \(m_B\), so that
\begin{equation}
\PhiN\sim -\frac{Gm_B}{r},
\qquad
\nabla\PhiN\sim \frac{Gm_B}{r^2},
\qquad
\nabla^3\PhiN\sim \frac{Gm_B}{r^4},
\qquad
\nabla^5\PhiN\sim \frac{Gm_B}{r^6}.
\label{eq:app-potential-scaling}
\end{equation}
Then the leading monopole force scales as
\begin{equation}
F_{\rm mono}\sim \frac{MGm_B}{r^2}.
\label{eq:app-Fmono-scaling}
\end{equation}
The generic quadrupole correction scales as
\begin{equation}
F_{\rm quad}\sim Q\,\nabla^3\PhiN
\sim
(Ma^2)\frac{Gm_B}{r^4},
\label{eq:app-Fquad-scaling}
\end{equation}
so the relative suppression is
\begin{equation}
\frac{|F_{\rm quad}|}{|F_{\rm mono}|}
=
O\!\left(\frac{a^2}{r^2}\right).
\label{eq:app-quad-suppression}
\end{equation}
Likewise, the next isotropic correction from the fourth moment scales as
\begin{equation}
F_{(4)}
\sim
(Ma^4)\frac{Gm_B}{r^6},
\qquad
\frac{|F_{(4)}|}{|F_{\rm mono}|}
=
O\!\left(\frac{a^4}{r^4}\right).
\label{eq:app-fourth-suppression}
\end{equation}

Therefore, provided
\begin{equation}
a\ll r,
\label{eq:app-small-body}
\end{equation}
the center-of-mass motion is controlled by the monopole force to leading order. This is
the quantitative content of the coherent-defect / small-body limit used in the main text.

\subsection{Equation of motion and point-particle limit}
\label{app:moments-eom}

Combining the momentum identity \eqref{eq:app-total-momentum} with the force expansion
\eqref{eq:app-force-quadrupole}, the center-of-mass equation of motion is
\begin{equation}
M\ddot{\mathbf X}_{\rm cm}
=
- M\,\gradThree\PhiN(\mathbf X_{\rm cm},t)
-\frac12\,Q_{jk}\,\gradThree\partial_j\partial_k\PhiN(\mathbf X_{\rm cm},t)
+\cdots
+\mathbf F_{\partial W}.
\label{eq:app-cm-eom-expanded}
\end{equation}
Under the retained Newtonian-order assumptions
\begin{equation}
\frac{a}{r}\ll 1,
\qquad
\mathbf F_{\partial W}\ \text{subleading},
\qquad
\text{odd moments negligible or symmetry-suppressed},
\label{eq:app-retained-assumptions}
\end{equation}
this reduces to
\begin{equation}
M\ddot{\mathbf X}_{\rm cm}
=
- M\,\gradThree\PhiN(\mathbf X_{\rm cm},t)
+
O\!\left(M\,\frac{a^2}{r^2}\,\frac{\PhiN}{r}\right).
\label{eq:app-cm-eom-mono}
\end{equation}
Equivalently, at leading order the motion follows from the point-particle Lagrangian
\begin{equation}
L_{\rm N}
=
\frac12 M\dot{\mathbf X}_{\rm cm}^{\,2}
-
M\PhiN(\mathbf X_{\rm cm},t).
\label{eq:app-point-particle-lag}
\end{equation}

This is the precise sense in which the coherent defect reduces to a Newtonian point mass.
The monopole reduction is not an additional ansatz placed on top of the force law; it is
the leading term in an explicit moment expansion whose small parameter is \(a/r\), with
higher moments and boundary-flux effects organized systematically as corrections.

\paragraph{Special simplification in the exterior two-body problem.}
For the exterior field of an isolated source away from the source support,
\(\nabla^2\PhiN=0\), so the isotropic quadrupole correction in
\eqref{eq:app-force-isotropic} vanishes. In that common case,
\begin{equation}
M\ddot{\mathbf X}_{\rm cm}
=
- M\,\gradThree\PhiN
+
O(Ma^4\nabla^5\PhiN)
+
\mathbf F_{\partial W},
\label{eq:app-cm-eom-harmonic}
\end{equation}
and the first nonzero finite-size term is pushed one order higher for an isotropic,
parity-even defect. This is a useful technical reason why the Newtonian two-body monopole
limit is especially robust in the present construction.

\claimclass{\Controlled. The moment expansion itself is exact once the worldtube integrals
are defined, but the reduction to the monopole point-particle limit requires the small-body
hierarchy \(a/r\ll 1\), smooth external fields across the support, negligible boundary
fluxes at the retained order, and the coherent-defect centering conditions used in the
main text.}

\section{Added-mass computations and topology comparison}
\label{app:added-mass}

This appendix gives the explicit added-mass calculations underlying the discussion in
\cref{sec:added-mass}. The goal is to make two points precise.

First, in the low-Mach, irrotational exterior-flow regime, the added mass is a purely
kinematic quantity determined by the harmonic flow problem outside the excluded region.
Second, for the present toy model the relevant excluded region is not a compact \(4\)-ball
but a \(w\)-uniform throat. That distinction changes the effective dimension of the
harmonic problem and therefore changes the coefficient \(\kadd\).

The result is:
\begin{equation}
\kadd^{\rm throat}=\frac12,
\qquad
\kadd^{(B^4)}=\frac13.
\label{eq:app-kadd-results}
\end{equation}
The difference is not a fit. It is a topology/geometry discriminator.

\subsection{Exterior potential-flow setup}
\label{app:added-mass-setup}

We work in the standard added-mass regime:
\begin{enumerate}[leftmargin=1.5em]
  \item the ambient medium has approximately constant density \(\rho_0\),
  \item the flow is sufficiently slow that compressibility corrections are negligible at
  the retained order,
  \item the exterior flow is irrotational, so it may be written as
  \(\mathbf v_{\rm ext}=\nabla \phi\), and
  \item the disturbance decays at spatial infinity.
\end{enumerate}
The added mass is defined by the kinetic energy stored in the exterior rearrangement:
\begin{equation}
T_{\rm ext}
=
\frac{\rho_0}{2}\int_{\rm ext} |\nabla \phi|^2\,dV
=
\frac12 M_{\rm add} U^2,
\label{eq:app-added-mass-def}
\end{equation}
for a body translating at speed \(U=|\mathbf U|\).

In isotropic cases the added-mass tensor is proportional to the identity, so it is enough
to compute the scalar \(M_{\rm add}\). We denote the displaced mass by
\begin{equation}
M_{\rm disp}=\rho_0 V_{\rm exc},
\label{eq:app-displaced-mass}
\end{equation}
and define the dimensionless coefficient
\begin{equation}
M_{\rm add}=\kadd M_{\rm disp}.
\label{eq:app-kadd-def}
\end{equation}

\subsection{General \(d_{\rm eff}\)-dimensional ball}
\label{app:added-mass-general-d}

The cleanest derivation is for a rigid ball of radius \(a\) in an effective exterior-flow
dimension \(d_{\rm eff}\ge 2\). Let \(r=|\mathbf x|\), with \(\mathbf x\in\mathbb R^{d_{\rm eff}}\),
and let \(\mathbf U\) be the translation velocity. Outside the excluded ball \(r>a\), the
potential solves
\begin{equation}
\nabla_{d_{\rm eff}}^2 \phi = 0,
\qquad
\phi\to 0 \quad \text{as } r\to\infty,
\label{eq:app-laplace-d}
\end{equation}
with Neumann boundary condition
\begin{equation}
\partial_r \phi\big|_{r=a}=\mathbf U\!\cdot\!\mathbf n,
\qquad
\mathbf n=\frac{\mathbf x}{r}.
\label{eq:app-neumann-boundary}
\end{equation}
The decaying \(\ell=1\) harmonic with the required angular structure is
\begin{equation}
\phi(\mathbf x)
=
A\,\frac{\mathbf U\!\cdot\!\mathbf x}{r^{d_{\rm eff}}},
\label{eq:app-harmonic-ansatz}
\end{equation}
and since
\begin{equation}
\mathbf U\!\cdot\!\mathbf x
=
Ur\cos\theta,
\end{equation}
this is
\begin{equation}
\phi(r,\theta)=A\,U\cos\theta\,r^{1-d_{\rm eff}}.
\end{equation}
Applying \eqref{eq:app-neumann-boundary} gives
\begin{equation}
A(1-d_{\rm eff})a^{-d_{\rm eff}}=1,
\end{equation}
hence
\begin{equation}
\boxed{
\phi(\mathbf x)
=
-\frac{a^{d_{\rm eff}}}{d_{\rm eff}-1}\,
\frac{\mathbf U\!\cdot\!\mathbf x}{r^{d_{\rm eff}}}.
}
\label{eq:app-general-potential}
\end{equation}

The kinetic energy is most conveniently evaluated by Green's identity:
\begin{equation}
\int_{\rm ext}|\nabla\phi|^2\,dV
=
\int_{\partial {\rm ext}}\phi\,\partial_{n_{\rm ext}}\phi\,dS,
\label{eq:app-green-identity}
\end{equation}
where \(n_{\rm ext}\) is the outward normal of the exterior domain. The contribution at
infinity vanishes, and on the inner boundary \(r=a\) one has
\begin{equation}
n_{\rm ext}=-\mathbf n,
\qquad
\partial_{n_{\rm ext}}\phi=-\partial_r\phi=-\mathbf U\!\cdot\!\mathbf n.
\label{eq:app-normal-convention}
\end{equation}
From \eqref{eq:app-general-potential},
\begin{equation}
\phi(a,\theta)=-\frac{a}{d_{\rm eff}-1}\,\mathbf U\!\cdot\!\mathbf n,
\end{equation}
so
\begin{align}
\int_{\rm ext}|\nabla\phi|^2\,dV
&=
\frac{a}{d_{\rm eff}-1}
\int_{r=a}(\mathbf U\!\cdot\!\mathbf n)^2\,dS
\nonumber\\
&=
\frac{a}{d_{\rm eff}-1}\,
U_iU_j
\int_{r=a}n_i n_j\,dS.
\label{eq:app-surface-eval}
\end{align}
Using isotropy,
\begin{equation}
\int_{r=a}n_i n_j\,dS
=
\frac{S_{d_{\rm eff}-1}a^{d_{\rm eff}-1}}{d_{\rm eff}}\,
\delta_{ij},
\label{eq:app-sphere-isotropy}
\end{equation}
where \(S_{d_{\rm eff}-1}\) is the area of the unit \((d_{\rm eff}-1)\)-sphere, we obtain
\begin{equation}
\int_{\rm ext}|\nabla\phi|^2\,dV
=
\frac{S_{d_{\rm eff}-1}a^{d_{\rm eff}}}{d_{\rm eff}(d_{\rm eff}-1)}\,U^2.
\label{eq:app-integral-result}
\end{equation}
Therefore
\begin{equation}
T_{\rm ext}
=
\frac{\rho_0}{2}\,
\frac{S_{d_{\rm eff}-1}a^{d_{\rm eff}}}{d_{\rm eff}(d_{\rm eff}-1)}\,U^2.
\label{eq:app-Text-general}
\end{equation}

Now the displaced volume of the \(d_{\rm eff}\)-ball is
\begin{equation}
V_{d_{\rm eff}}(a)
=
\frac{\pi^{d_{\rm eff}/2}}{\Gamma\!\left(\frac{d_{\rm eff}}{2}+1\right)}\,a^{d_{\rm eff}},
\qquad
S_{d_{\rm eff}-1}=d_{\rm eff}V_{d_{\rm eff}}(1),
\label{eq:app-volume-area}
\end{equation}
so
\begin{equation}
\rho_0 \frac{S_{d_{\rm eff}-1}a^{d_{\rm eff}}}{d_{\rm eff}}
=
\rho_0 V_{d_{\rm eff}}(a)
=
M_{\rm disp}.
\label{eq:app-disp-identity}
\end{equation}
Hence
\begin{equation}
T_{\rm ext}
=
\frac12
\left(
\frac{1}{d_{\rm eff}-1}M_{\rm disp}
\right)U^2,
\label{eq:app-Text-final-general}
\end{equation}
which means
\begin{equation}
\boxed{
\kadd(d_{\rm eff})=\frac{1}{d_{\rm eff}-1}.
}
\label{eq:app-kadd-general}
\end{equation}

This is the master formula used in the main text.

\subsection{The \(w\)-uniform throat: effective \(d_{\rm eff}=3\)}
\label{app:added-mass-throat}

The actual object used in the toy model is not a compact \(4\)-ball. It is a
\(w\)-extended throat, and in the quasi-static regime relevant for the \(1\)PN reduction
the exterior rearrangement is taken to be uniform along \(w\). The excluded region may be
modeled schematically as
\begin{equation}
\mathcal T \simeq B^3(a)\times I_L,
\label{eq:app-throat-geometry}
\end{equation}
where \(B^3(a)\) is a 3-ball of radius \(a\) in the brane directions and \(I_L\) is an
interval of length \(L\) along \(w\).

For motion parallel to the brane, the exterior potential is assumed independent of \(w\):
\begin{equation}
\phi_{\rm throat}(x,y,z,w)=\phi_3(x,y,z),
\qquad
\partial_w\phi_{\rm throat}=0.
\label{eq:app-throat-factorization}
\end{equation}
Then the \(4\)-spatial Laplacian reduces slice-by-slice to the ordinary \(3\)-dimensional
Laplace equation,
\begin{equation}
\nabla_4^2\phi_{\rm throat}
=
\nabla_3^2\phi_3
=
0
\qquad
\text{outside } B^3(a).
\label{eq:app-throat-laplace}
\end{equation}
Thus the relevant harmonic problem is the \(d_{\rm eff}=3\) ball problem on each
\(w\)-slice, not the \(d_{\rm eff}=4\) problem of a compact hypersphere.

Applying \eqref{eq:app-general-potential} with \(d_{\rm eff}=3\) gives
\begin{equation}
\phi_3(\x)
=
-\frac{a^3}{2}\,
\frac{\mathbf U\!\cdot\!\x}{|\x|^3}
=
-\frac{a^3}{2}\,\frac{U\cos\theta}{r^2},
\qquad
r=|\x|.
\label{eq:app-throat-potential}
\end{equation}
The exterior kinetic energy per unit length in \(w\) is therefore the standard sphere
result
\begin{equation}
\frac{T_{\rm ext}^{\rm throat}}{L}
=
\frac{\pi\rho_0 a^3}{3}\,U^2.
\label{eq:app-throat-energy-density}
\end{equation}
Multiplying by \(L\),
\begin{equation}
T_{\rm ext}^{\rm throat}
=
\frac{\pi\rho_0 a^3 L}{3}\,U^2.
\label{eq:app-throat-energy}
\end{equation}

The displaced mass of the throat is the ambient density times the excluded volume
\begin{equation}
V_{\rm exc}^{\rm throat}=L\,V_3(a)=L\,\frac{4\pi a^3}{3},
\qquad
M_{\rm disp}^{\rm throat}
=
\rho_0\,\frac{4\pi a^3L}{3}.
\label{eq:app-throat-disp}
\end{equation}
Hence
\begin{equation}
T_{\rm ext}^{\rm throat}
=
\frac12
\left(
\frac12\,M_{\rm disp}^{\rm throat}
\right)U^2,
\label{eq:app-throat-added-mass}
\end{equation}
so
\begin{equation}
\boxed{
\kadd^{\rm throat}=\frac12.
}
\label{eq:app-kadd-throat}
\end{equation}

This is the value used in the main text. The essential point is that the uniform throat
reduces the added-mass problem to a \(3\)-dimensional slice problem.

\subsection{Counterfactual compact \(4\)-ball}
\label{app:added-mass-bubble}

For comparison, consider instead a compact excluded \(4\)-ball
\begin{equation}
B^4(a)=\{\mathbf X\in\mathbb R^4:\ |\mathbf X|<a\},
\label{eq:app-b4-geometry}
\end{equation}
moving through the medium. Now the exterior flow is genuinely four-dimensional, so
\(d_{\rm eff}=4\). Applying \eqref{eq:app-general-potential} gives
\begin{equation}
\phi_{B^4}(\mathbf X)
=
-\frac{a^4}{3}\,
\frac{\mathbf U\!\cdot\!\mathbf X}{R^4},
\qquad
R=|\mathbf X|.
\label{eq:app-b4-potential}
\end{equation}
The displaced volume and mass are
\begin{equation}
V_4(a)=\frac{\pi^2}{2}a^4,
\qquad
M_{\rm disp}^{(B^4)}=\rho_0\frac{\pi^2}{2}a^4.
\label{eq:app-b4-disp}
\end{equation}
Using \eqref{eq:app-kadd-general} with \(d_{\rm eff}=4\),
\begin{equation}
T_{\rm ext}^{(B^4)}
=
\frac12
\left(
\frac13 M_{\rm disp}^{(B^4)}
\right)U^2
=
\frac{\rho_0\pi^2 a^4}{12}\,U^2,
\label{eq:app-b4-energy}
\end{equation}
and therefore
\begin{equation}
\boxed{
\kadd^{(B^4)}=\frac13.
}
\label{eq:app-kadd-b4}
\end{equation}

This is the counterfactual value quoted in the main text. It is useful precisely because
it differs from the throat value:
\begin{equation}
\frac12 \neq \frac13.
\label{eq:app-half-not-third}
\end{equation}
The added-mass coefficient therefore distinguishes a \(w\)-uniform throat from a compact
\(4\)-ball.

\subsection{Interpretation as a topology/geometry discriminator}
\label{app:added-mass-interpretation}

The comparison above shows why \(\kadd\) should not be treated as a phenomenological knob.
In the incompressible potential-flow regime, once the exterior problem is specified, the
coefficient is fixed. What changes between the throat and the compact bubble is not the
equation of state and not an arbitrary coupling constant, but the effective dimension of
the harmonic rearrangement problem:
\begin{equation}
d_{\rm eff}=3
\quad\Longrightarrow\quad
\kadd=\frac12,
\qquad
d_{\rm eff}=4
\quad\Longrightarrow\quad
\kadd=\frac13.
\label{eq:app-deff-summary}
\end{equation}

For the present toy model, the relevant statement is therefore:
\begin{quote}
Because the defect is treated as a \(w\)-uniform throat whose exterior rearrangement is
slice-wise uniform in \(w\), the appropriate added-mass problem is effectively
three-dimensional, and the carried-forward value is \(\kadd=1/2\).
\end{quote}

This also explains the limited scope of the result. The derivation assumes a quasi-static,
low-Mach, irrotational exterior flow and uniformity along \(w\). If those assumptions are
relaxed---for example, if one allows nonuniform \(w\)-dependent exterior rearrangement or
strong compressibility---then the effective inertia could acquire additional corrections.
Within the controlled regime used in this paper, however, the coefficient is fixed.

\paragraph{Compact comparison table.}
For convenience, the two relevant cases are summarized in
\cref{tab:app-added-mass-comparison}.

\begin{table}[t]
  \centering
  \caption{Added-mass comparison for the two geometries discussed in the paper.}
  \label{tab:app-added-mass-comparison}
  \small
  \renewcommand{\arraystretch}{1.15}
  \begin{tabularx}{\linewidth}{@{}L c c c@{}}
    \toprule
    Geometry & Effective flow dimension \(d_{\rm eff}\) & \(M_{\rm add}/M_{\rm disp}\) & Value of \(\kadd\) \\
    \midrule
    \(w\)-uniform throat \(B^3(a)\times I_L\) & \(3\) & \(1/(3-1)\) & \(1/2\) \\
    compact \(4\)-ball \(B^4(a)\) & \(4\) & \(1/(4-1)\) & \(1/3\) \\
    \bottomrule
  \end{tabularx}
\end{table}

\claimclass{\Controlled. The added-mass formulas are exact within the quasi-static,
irrotational potential-flow regime, but the choice between \(\kadd=1/2\) and
\(\kadd=1/3\) depends on the geometric/topological reduction used for the defect:
\(w\)-uniform throat versus compact \(4\)-ball.}

\section{Wake-overlap algebra and EIH ratio matching}
\label{app:wake}

This appendix gives the algebra behind the constructive wake calculation summarized in
\cref{sec:wake}. The main text stated the result in its most compact form:
an isotropic Fourier-space wake basis produces the parity-even cross tensor
\begin{equation}
L_{\rm cross}^{(AB)}
=
\frac{Gm_A m_B}{c^2 r_{AB}}
\left[
\Cpara\,\vct_A\!\cdot\!\vct_B
+
\CLong\,(\vct_A\!\cdot\!\nvec)(\vct_B\!\cdot\!\nvec)
\right],
\end{equation}
with
\begin{equation}
\Cpara=-\frac72,
\qquad
\CLong=-\frac12.
\end{equation}
The purpose of this appendix is to show how those coefficients arise from the overlap
integral, how the EIH ratio fixes a one-parameter family, and how the already-imported
barotropic closure collapses that family to the unique parity-even point used in the main
text.

Two remarks should be kept in mind throughout.
\begin{enumerate}[leftmargin=1.5em]
  \item This appendix concerns only the \emph{parity-even conservative} cross sector.
  Any parity-odd or dissipative structures are outside the scope of the present paper.
  \item The self-sector coefficient \(+\frac32\) multiplying \(v_A^2+v_B^2\) is
  \emph{not} part of this derivation. By construction, the wake overlap is bilinear in the
  velocities of two distinct bodies and therefore generates only cross terms.
\end{enumerate}

\subsection{Isotropic Fourier-space wake basis}
\label{app:wake-basis}

Let \(\mathbf k\in\mathbb R^3\) be the Fourier wavevector and define
\begin{equation}
k\equiv |\mathbf k|,
\qquad
\hat{\mathbf k}\equiv \frac{\mathbf k}{k}.
\end{equation}
The parity-even wake basis used in the symbolic module is
\begin{equation}
\mathbf u(\mathbf k;\vct)
=
\mathbf u_T(\mathbf k;\vct)
+
\mathbf u_H(\mathbf k;\vct)
+
\mathbf u_L(\mathbf k;\vct),
\label{eq:app-wake-total}
\end{equation}
with
\begin{align}
\mathbf u_T(\mathbf k;\vct)
&=
i\,a_T\,\frac{\mathbf P_T(\hat{\mathbf k})\,\vct}{k},
\label{eq:app-uT}
\\[3pt]
\mathbf u_H(\mathbf k;\vct)
&=
i\,a_H\,\frac{\mathbf k\times \vct}{k^2},
\label{eq:app-uH}
\\[3pt]
\mathbf u_L(\mathbf k;\vct)
&=
i\,a_L\,\frac{\mathbf k(\mathbf k\!\cdot\!\vct)}{k^3}.
\label{eq:app-uL}
\end{align}
Here
\begin{equation}
\mathbf P_T(\hat{\mathbf k})=\mathbf I-\hat{\mathbf k}\hat{\mathbf k}
\label{eq:app-PT}
\end{equation}
is the usual transverse projector, and we fix the overall projector amplitude to
\begin{equation}
a_T=1
\label{eq:app-aT-one}
\end{equation}
as a normalization convention.

The three pieces have the intended interpretations:
\begin{itemize}[leftmargin=1.5em]
  \item \(\mathbf u_T\): translation/shear transverse response,
  \item \(\mathbf u_H\): helical transverse response,
  \item \(\mathbf u_L\): longitudinal/compressible response.
\end{itemize}
Each mode scales as \(v/k\), so the bilinear overlap behaves as \(u_A\!\cdot\!u_B\sim
1/k^2\). With the Fourier convention used in the symbolic scripts, this is the scaling
that produces a \(1/r\) pair interaction after transformation to position space. Any
conventional \(2\pi\)-normalization factors are absorbed into the overall coefficient
\(\Kvec\), so the only physically meaningful outputs are the final dimensionless ratios
and coefficients.

\subsection{Pair overlap kernel}
\label{app:wake-overlap}

For two bodies \(A\) and \(B\), separated by
\begin{equation}
\mathbf r\equiv \mathbf r_{AB}=\mathbf X_A-\mathbf X_B,
\qquad
r\equiv r_{AB}=|\mathbf r|,
\qquad
\nvec\equiv \frac{\mathbf r}{r},
\label{eq:app-wake-r}
\end{equation}
the parity-even overlap kernel is taken to be
\begin{equation}
I_{AB}(\mathbf r)
=
\Kvec
\int_{\mathbb R^3} d^3k\;
\mathbf u(\mathbf k;\vct_A)\!\cdot\!\mathbf u(-\mathbf k;\vct_B)\,
e^{i\mathbf k\cdot\mathbf r}.
\label{eq:app-overlap-kernel}
\end{equation}
The sign flip on the second wake is the natural one for the bilinear Fourier overlap of
two real-space disturbances.

Because the EIH target is parity-even, we keep only the parity-even part of the overlap.
The algebra is simplest when written in tensor language. Introduce
\begin{equation}
P_{T,ij}=\delta_{ij}-\hat k_i\hat k_j,
\qquad
P_{L,ij}=\hat k_i\hat k_j.
\label{eq:app-projector-tensors}
\end{equation}
Then
\begin{equation}
u_{T,i}(\mathbf k;\vct)= i\,\frac{P_{T,ij}v_j}{k},
\qquad
u_{L,i}(\mathbf k;\vct)= i\,a_L\,\frac{P_{L,ij}v_j}{k},
\label{eq:app-uT-uL-tensor}
\end{equation}
while the helical piece satisfies
\begin{equation}
(\mathbf k\times\vct_A)\!\cdot\!(\mathbf k\times\vct_B)
=
k^2\,\vct_A\!\cdot\!\vct_B
-
(\mathbf k\!\cdot\!\vct_A)(\mathbf k\!\cdot\!\vct_B).
\label{eq:app-cross-identity}
\end{equation}

Using \(P_TP_L=0\) and \(\mathbf k\!\cdot\!(\mathbf k\times\vct)=0\), the
\(T\)-\(L\) and \(L\)-\(H\) cross terms vanish identically. The \(T\)-\(H\) mixed term is
parity-odd and belongs to a pseudovector sector, not to the parity-even EIH tensor being
matched here. Therefore the retained even overlap comes entirely from the
\(T\)-\(T\), \(H\)-\(H\), and \(L\)-\(L\) pieces:
\begin{align}
\mathbf u(\mathbf k;\vct_A)\!\cdot\!\mathbf u(-\mathbf k;\vct_B)\Big|_{\rm even}
={}&
-\frac{\vct_A\!\cdot\!\vct_B}{k^2}
+\frac{(\mathbf k\!\cdot\!\vct_A)(\mathbf k\!\cdot\!\vct_B)}{k^4}
\nonumber\\
&
+a_H^2
\left[
\frac{\vct_A\!\cdot\!\vct_B}{k^2}
-\frac{(\mathbf k\!\cdot\!\vct_A)(\mathbf k\!\cdot\!\vct_B)}{k^4}
\right]
-a_L^2\frac{(\mathbf k\!\cdot\!\vct_A)(\mathbf k\!\cdot\!\vct_B)}{k^4}.
\label{eq:app-overlap-even-raw}
\end{align}
Collecting terms,
\begin{equation}
\boxed{
\mathbf u_A\!\cdot\!\mathbf u_B(-\mathbf k)\Big|_{\rm even}
=
\frac{-1+a_H^2}{k^2}\,\vct_A\!\cdot\!\vct_B
+
\frac{1-a_H^2-a_L^2}{k^4}\,
(\mathbf k\!\cdot\!\vct_A)(\mathbf k\!\cdot\!\vct_B).
}
\label{eq:app-overlap-even-collected}
\end{equation}

This is the key algebraic simplification. Once the wake basis is declared, the parity-even
tensor structure depends only on \(a_H^2\) and \(a_L^2\).

\subsection{Fourier tensor integrals}
\label{app:wake-fourier-integrals}

For \(r>0\), the two needed Fourier integrals are
\begin{equation}
I_0(r)
\equiv
\int d^3k\;\frac{e^{i\mathbf k\cdot\mathbf r}}{k^2}
=
\frac{2\pi^2}{r},
\label{eq:app-I0}
\end{equation}
and
\begin{equation}
I_{ij}(r)
\equiv
\int d^3k\;\frac{k_i k_j}{k^4}\,e^{i\mathbf k\cdot\mathbf r}
=
\frac{\pi^2}{r}\,\bigl(\delta_{ij}-n_i n_j\bigr).
\label{eq:app-Iij}
\end{equation}
A convenient derivation of \eqref{eq:app-Iij} is to introduce
\begin{equation}
J(r)\equiv \int d^3k\;\frac{e^{i\mathbf k\cdot\mathbf r}}{k^4},
\label{eq:app-J-def}
\end{equation}
for which the chosen Fourier convention gives
\begin{equation}
J(r)=-\pi^2 r
\qquad (r>0),
\label{eq:app-J-result}
\end{equation}
and then use
\begin{equation}
I_{ij}(r)= -\partial_i\partial_j J(r).
\label{eq:app-Iij-from-J}
\end{equation}
Since
\begin{equation}
\partial_i\partial_j r
=
\frac{1}{r}\bigl(\delta_{ij}-n_i n_j\bigr),
\qquad
n_i\equiv \frac{r_i}{r},
\label{eq:app-second-derivative-r}
\end{equation}
equation \eqref{eq:app-Iij} follows immediately.

Contracting \eqref{eq:app-Iij} with \(v_{A,i}v_{B,j}\) gives
\begin{equation}
v_{A,i}v_{B,j}I_{ij}(r)
=
\frac{\pi^2}{r}
\left[
\vct_A\!\cdot\!\vct_B
-
(\vct_A\!\cdot\!\nvec)(\vct_B\!\cdot\!\nvec)
\right].
\label{eq:app-Iij-contracted}
\end{equation}

\subsection{Extraction of the parity-even cross tensor}
\label{app:wake-extraction}

Substituting \eqref{eq:app-overlap-even-collected},
\eqref{eq:app-I0}, and \eqref{eq:app-Iij-contracted} into
\eqref{eq:app-overlap-kernel}, we obtain
\begin{align}
I_{AB}^{\rm even}(r)
={}&
\Kvec\,(-1+a_H^2)\,I_0(r)\,\vct_A\!\cdot\!\vct_B
\nonumber\\
&
+\Kvec\,(1-a_H^2-a_L^2)\,
v_{A,i}v_{B,j}I_{ij}(r).
\label{eq:app-IAB-even-step}
\end{align}
Hence
\begin{align}
I_{AB}^{\rm even}(r)
=
\frac{\Kvec\pi^2}{r}
\Big[
&\bigl(2(-1+a_H^2)+(1-a_H^2-a_L^2)\bigr)\,\vct_A\!\cdot\!\vct_B
\nonumber\\
&
-\bigl(1-a_H^2-a_L^2\bigr)
(\vct_A\!\cdot\!\nvec)(\vct_B\!\cdot\!\nvec)
\Big].
\label{eq:app-IAB-even-expanded}
\end{align}
Simplifying,
\begin{equation}
\boxed{
I_{AB}^{\rm even}(r)
=
\frac{\Kvec\pi^2}{r}
\left[
(-1+a_H^2-a_L^2)\,\vct_A\!\cdot\!\vct_B
+
(-1+a_H^2+a_L^2)\,
(\vct_A\!\cdot\!\nvec)(\vct_B\!\cdot\!\nvec)
\right].
}
\label{eq:app-IAB-even-final}
\end{equation}

Comparing with the standard cross-sector form
\begin{equation}
L_{\rm cross}^{(AB)}
=
\frac{Gm_A m_B}{c^2 r}
\left[
\Cpara\,\vct_A\!\cdot\!\vct_B
+
\CLong\,(\vct_A\!\cdot\!\nvec)(\vct_B\!\cdot\!\nvec)
\right],
\label{eq:app-cross-standard}
\end{equation}
we read off
\begin{equation}
\boxed{
\Cpara
=
\Kvec\pi^2\bigl(-1+a_H^2-a_L^2\bigr),
\qquad
\CLong
=
\Kvec\pi^2\bigl(-1+a_H^2+a_L^2\bigr).
}
\label{eq:app-Cpara-CLong}
\end{equation}

Defining the longitudinal wake fraction
\begin{equation}
\alphaSq \equiv a_L^2,
\label{eq:app-alphaSq-def}
\end{equation}
this becomes
\begin{equation}
\Cpara
=
\Kvec\pi^2\bigl(-1+a_H^2-\alphaSq\bigr),
\qquad
\CLong
=
\Kvec\pi^2\bigl(-1+a_H^2+\alphaSq\bigr).
\label{eq:app-Cpara-CLong-alpha}
\end{equation}
These are precisely the geometric shape formulas printed by the symbolic wake module.

\subsection{EIH ratio matching}
\label{app:wake-ratio}

The EIH target values are
\begin{equation}
\Cpara^{\rm(EIH)}=-\frac72,
\qquad
\CLong^{\rm(EIH)}=-\frac12.
\label{eq:app-EIH-targets}
\end{equation}
Their ratio is
\begin{equation}
\frac{\CLong^{\rm(EIH)}}{\Cpara^{\rm(EIH)}}=\frac{1}{7}.
\label{eq:app-EIH-ratio}
\end{equation}
Applying this to \eqref{eq:app-Cpara-CLong-alpha} yields
\begin{equation}
\frac{-1+a_H^2+\alphaSq}{-1+a_H^2-\alphaSq}=\frac{1}{7}.
\label{eq:app-ratio-equation}
\end{equation}
Solving,
\begin{equation}
7(-1+a_H^2+\alphaSq)= -1+a_H^2-\alphaSq,
\end{equation}
so
\begin{equation}
6a_H^2+8\alphaSq=6,
\qquad\Longleftrightarrow\qquad
3a_H^2+4\alphaSq=3.
\label{eq:app-family-constraint}
\end{equation}
Equivalently,
\begin{equation}
\boxed{
\alphaSq=\frac34(1-a_H^2).
}
\label{eq:app-alpha-family}
\end{equation}

Thus the EIH ratio alone does not yet select a unique point; it leaves a
one-parameter family of parity-even solutions. The overall normalization \(\Kvec\) is then
fixed by matching either target coefficient. Using \(\Cpara=-7/2\),
\begin{equation}
-\frac72
=
\Kvec\pi^2\bigl(-1+a_H^2-\alphaSq\bigr).
\label{eq:app-K-from-para}
\end{equation}
Substituting \eqref{eq:app-alpha-family},
\begin{equation}
-1+a_H^2-\alphaSq
=
-\frac74(1-a_H^2),
\end{equation}
so
\begin{equation}
\boxed{
\Kvec=\frac{2}{\pi^2(1-a_H^2)}.
}
\label{eq:app-K-family}
\end{equation}
Equations \eqref{eq:app-alpha-family} and \eqref{eq:app-K-family} are exactly the
one-parameter EIH-matching family reported in the main text.

\subsection{Collapse to the unique parity-even point}
\label{app:wake-unique-point}

To collapse the family, the main text imports the thermodynamic closure already fixed by
the barotropic sector:
\begin{equation}
\alphaSq = 1-\frac{1}{n-1}.
\label{eq:app-alpha-thermo}
\end{equation}
With the optical result
\begin{equation}
n=5,
\label{eq:app-n-five-wake}
\end{equation}
this gives
\begin{equation}
\alphaSq=\frac34.
\label{eq:app-alpha-three-quarters}
\end{equation}
Substituting into \eqref{eq:app-alpha-family},
\begin{equation}
\frac34=\frac34(1-a_H^2)
\qquad\Longrightarrow\qquad
a_H^2=0.
\label{eq:app-aHsq-zero}
\end{equation}
Hence the minimal real parity-even solution is
\begin{equation}
a_H=0,
\qquad
a_L=\pm \frac{\sqrt3}{2},
\label{eq:app-aH-aL-final}
\end{equation}
where only \(a_L^2=\alphaSq\) enters the conservative even tensor. Equation
\eqref{eq:app-K-family} then gives
\begin{equation}
\Kvec=\frac{2}{\pi^2}.
\label{eq:app-K-final}
\end{equation}

Back-substituting into \eqref{eq:app-Cpara-CLong-alpha},
\begin{align}
\Cpara
&=
\frac{2}{\pi^2}\pi^2\left(-1-\frac34\right)
=
-\frac72,
\label{eq:app-final-Cpara}
\\[3pt]
\CLong
&=
\frac{2}{\pi^2}\pi^2\left(-1+\frac34\right)
=
-\frac12.
\label{eq:app-final-CLong}
\end{align}
Therefore
\begin{equation}
\boxed{
\alphaSq=\frac34,
\qquad
a_H=0,
\qquad
\Kvec=\frac{2}{\pi^2},
\qquad
\Cpara=-\frac72,
\qquad
\CLong=-\frac12.
}
\label{eq:app-wake-final-box}
\end{equation}

\paragraph{Relation to the symbolic module.}
The WL script evaluates the same overlap by direct spherical integration after aligning the
separation vector with the \(z\)-axis. The output shapes printed by that module are
\begin{equation}
\texttt{VecParaShape}=\pi^2(-1+a_H^2-a_L^2),
\qquad
\texttt{VecLongShape}=\pi^2(-1+a_H^2+a_L^2),
\end{equation}
which are exactly the coefficients derived analytically in
\eqref{eq:app-Cpara-CLong}. The module then solves the EIH ratio constraint, finds the
minimal real solution \(a_H=0\), \(a_L=\pm \sqrt{3}/2\), and returns
\(\Kvec=2/\pi^2\), in agreement with \eqref{eq:app-wake-final-box}.

\claimclass{\Controlled for the isotropic linear-response wake basis, the restriction to
the parity-even conservative sector, and the thermodynamic closure used to select
\(\alphaSq\). Exact for the overlap algebra, the Fourier tensor integrals, the EIH-ratio
family \eqref{eq:app-alpha-family}--\eqref{eq:app-K-family}, and the final recovery of
\(\Cpara=-7/2\), \(\CLong=-1/2\) once the closure inputs are supplied.}

\section{Adiabatic breathing-response closure}
\label{app:adiabatic-closure}

This appendix records the algebra behind the explicit adiabatic internal-response protocol
used in \cref{sec:kappaPV-closure} to obtain
\begin{equation}
\kPV=\frac32.
\label{eq:app-kPV-target}
\end{equation}
The main text treated this as a declared \Protocol closure rather than as an exact theorem
of the parent \(4\)D dynamics. The purpose of the present appendix is to make that status
completely transparent: first by writing the generic adiabatic-elimination formula, and
then by specializing to the one-degree-of-freedom throat model actually used in the paper.

The essential idea is simple. A slowly varying external environment changes the
equilibrium configuration of the defect. In the adiabatic regime, the internal geometry
re-equilibrates quickly compared to orbital timescales, so one may eliminate the geometry
variables at fixed external control parameter and read off the resulting effective
rest-energy response. In general this response is protocol-dependent. Here we adopt one
specific protocol: a one-DOF throat with fixed aspect ratio, driven by the slowly varying
ambient density.

\subsection{Generic adiabatic-elimination formula}
\label{app:adiabatic-generic}

Let \(q^a\) denote the internal geometry variables of the defect. In the full parent
theory one may take, for example,
\begin{equation}
q^a=(a,L,\dots),
\end{equation}
but the derivation below is initially independent of how many such variables are retained.
Introduce a small external control parameter \(\epsilon\) encoding the weak environmental
change relevant for the orbital sector. In the present application one may think of
\begin{equation}
\epsilon \sim \frac{\delta\rho}{\rho_0} \sim \frac{\PhiN}{c^2},
\label{eq:app-epsilon-rho-phi}
\end{equation}
where \(\rho\) is the ambient medium density evaluated at the defect location and \(\PhiN\)
is the slowly varying Newtonian potential.

Expand the total conservative energy about an equilibrium point
\((q_0^a,\epsilon=0)\):
\begin{equation}
H_{\rm tot}(q,\epsilon)
=
H_0
+
H_\epsilon\,\epsilon
+
\frac12 H_{\epsilon\epsilon}\,\epsilon^2
+
\frac12 K_{ab}\,\delta q^a\delta q^b
+
f_a\,\delta q^a\,\epsilon
+\cdots,
\label{eq:app-H-expand}
\end{equation}
where
\begin{equation}
\delta q^a\equiv q^a-q_0^a,
\qquad
K_{ab}\equiv
\left.\frac{\partial^2 H_{\rm tot}}{\partial q^a\partial q^b}\right|_0,
\qquad
f_a\equiv
\left.\frac{\partial^2 H_{\rm tot}}{\partial q^a\partial \epsilon}\right|_0.
\label{eq:app-Kab-fa}
\end{equation}
Here \(K_{ab}\) is the stiffness matrix of the internal geometry and \(f_a\) is the
mixed environment--geometry coupling vector.

In the adiabatic regime, the geometry relaxes to the instantaneous minimum of
\(H_{\rm tot}\) at each \(\epsilon\), so
\begin{equation}
\frac{\partial H_{\rm tot}}{\partial q^a}=0.
\end{equation}
To linear order in \(\epsilon\), this gives
\begin{equation}
\boxed{
\delta q^a_{\rm ad}
=
-(K^{-1})^{ab}f_b\,\epsilon
+
O(\epsilon^2).
}
\label{eq:app-dq-adiabatic}
\end{equation}
Substituting back into \eqref{eq:app-H-expand} yields the adiabatically eliminated
effective energy
\begin{equation}
\boxed{
H_{\rm eff}(\epsilon)
=
H_0
+
H_\epsilon\,\epsilon
+
\frac12
\left(
H_{\epsilon\epsilon}
-
f_a(K^{-1})^{ab}f_b
\right)\epsilon^2
+
O(\epsilon^3).
}
\label{eq:app-Heff-adiabatic}
\end{equation}

Equation \eqref{eq:app-Heff-adiabatic} is the generic source of the breathing/pressure-volume
response in the orbital ledger. It also shows immediately why \(\kPV\) is not an exact
kinematic constant of the theory: the response depends on the chosen control parameter
\(\epsilon\), on which geometry variables are kept, and on the stiffness/coupling data
\((K_{ab},f_a)\). To make progress one must therefore declare a specific response
protocol. The remainder of the appendix does exactly that.

\subsection{Declared one-DOF adiabatic throat protocol}
\label{app:adiabatic-1dof}

The explicit closure adopted in the paper is the following.

\paragraph{External control variable.}
The environment is characterized by the ambient medium density \(\rho\) at the defect
location. In the weak-field regime,
\begin{equation}
\frac{\delta\rho}{\rho_0}\simeq \frac{\PhiN}{c^2},
\label{eq:app-rho-phi-relation}
\end{equation}
so varying \(\rho\) is equivalent, at the retained order, to varying the local weak
Newtonian environment.

\paragraph{Adiabatic regime.}
The internal geometry relaxes rapidly compared to orbital timescales, so for each slowly
varying \(\rho\) the defect sits at the instantaneous minimum of an effective rest-energy
functional.

\paragraph{Reduced geometry.}
To obtain an analytic closure, we retain only one geometric degree of freedom, the throat
radius \(a\), and impose a fixed aspect ratio
\begin{equation}
L=\Lambda a,
\label{eq:app-fixed-aspect}
\end{equation}
with \(\Lambda\) treated as a constant of the reduced closure.

\paragraph{Minimal reduced energy functional.}
Neglecting any \(\rho\)-independent ``dead-weight'' stabilizer terms in the weak-field
response extraction, we take
\begin{equation}
F(a,\rho)=E_w+E_f+E_{\rm PV},
\label{eq:app-F-closure}
\end{equation}
with
\begin{align}
E_w(a,\rho)
&=
\frac{\mathcal C_w\,c_s(\rho)}{a},
\label{eq:app-Ew}
\\[3pt]
E_f(a,\rho)
&=
\frac{\mathcal C_f}{\rho\,a^2},
\label{eq:app-Ef}
\\[3pt]
E_{\rm PV}(a,\rho)
&=
\mathcal C_{\rm PV}\,P(\rho)\,V(a)
=
\mathcal C_{\rm PV}\,\KEOS \rho^n\,(\pi\Lambda a^3).
\label{eq:app-EPV}
\end{align}
The three sectors have clear interpretations:
\begin{itemize}[leftmargin=1.5em]
  \item \(E_w\): wave-supported contribution,
  \item \(E_f\): flow/rearrangement contribution,
  \item \(E_{\rm PV}\): pressure-volume or breathing work against the ambient medium.
\end{itemize}

The sound speed follows from the barotropic closure,
\begin{equation}
P(\rho)=\KEOS \rho^n,
\qquad
c_s^2(\rho)=\frac{dP}{d\rho}=n\KEOS \rho^{n-1},
\label{eq:app-cs-polytrope}
\end{equation}
so
\begin{equation}
c_s(\rho)\propto \rho^{(n-1)/2}.
\label{eq:app-cs-scaling}
\end{equation}

A useful scope note belongs here. If an additional \(\rho\)-independent stabilizer makes a
comparable contribution to \(F\), it dilutes the logarithmic density response and can
shift, or even eliminate, the \(\kPV=3/2\) solution derived below. The present closure
therefore assumes that such terms are subdominant in the weak-field operating regime used
for the \(1\)PN extraction.

\subsection{Equilibrium virial identity}
\label{app:adiabatic-virial}

Adiabatic equilibrium means
\begin{equation}
\partial_a F=0.
\label{eq:app-equilibrium-condition}
\end{equation}
Using the \(a\)-scalings
\begin{equation}
E_w\propto a^{-1},
\qquad
E_f\propto a^{-2},
\qquad
E_{\rm PV}\propto a^{3},
\label{eq:app-a-scalings}
\end{equation}
one finds
\begin{equation}
-\frac{E_w}{a}-\frac{2E_f}{a}+\frac{3E_{\rm PV}}{a}=0,
\end{equation}
or equivalently
\begin{equation}
\boxed{
E_w+2E_f=3E_{\rm PV}.
}
\label{eq:app-virial}
\end{equation}
Define the internal partition parameter
\begin{equation}
x\equiv \frac{E_f}{E_w}.
\label{eq:app-x-def}
\end{equation}
Then \eqref{eq:app-virial} implies
\begin{equation}
\frac{E_{\rm PV}}{E_w}=\frac{1+2x}{3},
\label{eq:app-EPV-over-Ew}
\end{equation}
and therefore the total equilibrium energy becomes
\begin{equation}
F
=
E_w\left(1+x+\frac{1+2x}{3}\right)
=
E_w\,\frac{4+5x}{3}.
\label{eq:app-F-in-x}
\end{equation}

\subsection{Density-response slope and the \texorpdfstring{$\kPV$}{kappaPV} estimator}
\label{app:adiabatic-slope}

Because \(a=a_*(\rho)\) minimizes \(F\) at fixed \(\rho\), the equilibrium logarithmic
derivative obeys an envelope-theorem form:
\begin{equation}
\frac{dF_{\rm eq}}{d\rho}
=
\left.\frac{\partial F}{\partial \rho}\right|_{a=a_*(\rho)}.
\label{eq:app-envelope}
\end{equation}
Using the \(\rho\)-scalings
\begin{equation}
E_w\propto \rho^{(n-1)/2},
\qquad
E_f\propto \rho^{-1},
\qquad
E_{\rm PV}\propto \rho^n,
\label{eq:app-rho-scalings}
\end{equation}
we obtain
\begin{equation}
\boxed{
\frac{d\ln F}{d\ln \rho}
=
\frac{\left(\frac{n-1}{2}\right)E_w - E_f + nE_{\rm PV}}{F}.
}
\label{eq:app-dlnF-general}
\end{equation}
Now eliminate \(E_{\rm PV}\) using \eqref{eq:app-virial} and write everything in terms of
\(x=E_f/E_w\). This gives
\begin{equation}
\boxed{
\frac{d\ln F}{d\ln\rho}
=
\frac{\frac{5n-3}{2}+(2n-3)x}{4+5x}.
}
\label{eq:app-dlnF-in-x}
\end{equation}
For the optical value already fixed in the main text,
\begin{equation}
n=5,
\label{eq:app-n-five}
\end{equation}
this reduces to
\begin{equation}
\boxed{
\frac{d\ln F}{d\ln\rho}
=
\frac{11+7x}{4+5x}.
}
\label{eq:app-dlnF-n5}
\end{equation}

To connect with the orbital \(1\)PN ledger, we use the same estimator as in the main text:
the total density-driven rest-energy response corresponds to
\(\krho+\kPV\), so with \(\krho=1\),
\begin{equation}
\kPV
\equiv
\frac{d\ln F}{d\ln\rho}
-
\krho
=
\frac{d\ln F}{d\ln\rho}-1.
\label{eq:app-kPV-estimator}
\end{equation}
The target value \(\kPV=3/2\) is therefore equivalent to
\begin{equation}
\frac{d\ln F}{d\ln\rho}
=
\krho+\kPV
=
\frac52.
\label{eq:app-target-slope}
\end{equation}
Substituting \eqref{eq:app-dlnF-n5} gives
\begin{equation}
\frac{11+7x}{4+5x}=\frac52
\qquad\Longrightarrow\qquad
\boxed{
x=\frac{2}{11}.
}
\label{eq:app-x-solution}
\end{equation}
Thus the explicit adiabatic closure yields
\begin{equation}
\boxed{
\kPV=\frac32.
}
\label{eq:app-kPV-solution}
\end{equation}

\subsection{Internal partition implied by the closure}
\label{app:adiabatic-partition}

Once \(x=2/11\) is known, the virial relation fixes the internal energy partition
uniquely. Since
\begin{equation}
x=\frac{E_f}{E_w}=\frac{2}{11},
\end{equation}
we have
\begin{equation}
\frac{E_{\rm PV}}{E_w}
=
\frac{1+2x}{3}
=
\frac{1+4/11}{3}
=
\frac{5}{11}.
\end{equation}
Therefore
\begin{equation}
\boxed{
E_w:E_f:E_{\rm PV}=11:2:5.
}
\label{eq:app-11-2-5}
\end{equation}
The corresponding fractions of the total are
\begin{equation}
\boxed{
(f_w,f_f,f_{\rm PV})
=
\left(
\frac{11}{18},
\frac{2}{18},
\frac{5}{18}
\right).
}
\label{eq:app-partition-fractions}
\end{equation}

This point is worth emphasizing. Within the declared adiabatic protocol, the value
\(\kPV=3/2\) is not an isolated fit parameter. It comes with a definite predicted
partition of the internal energy among wave, flow, and pressure-volume sectors. That
partition is one of the useful outputs of the closure and a natural target for later
numerical checks.

\subsection{Predicted breathing slope}
\label{app:adiabatic-breathing-slope}

The same closure predicts how the throat radius responds to changes in the ambient
density. Differentiate the equilibrium condition
\eqref{eq:app-equilibrium-condition} with respect to \(\ln\rho\). Since
\begin{equation}
\partial_a F
=
-\frac{E_w}{a}
-\frac{2E_f}{a}
+\frac{3E_{\rm PV}}{a},
\end{equation}
the condition \(\partial_a F=0\) is equivalent to
\begin{equation}
-E_w-2E_f+3E_{\rm PV}=0.
\label{eq:app-equilibrium-rewritten}
\end{equation}
Now take the total derivative with respect to \(\ln\rho\). Using the scalings
\eqref{eq:app-a-scalings} and \eqref{eq:app-rho-scalings}, one finds
\begin{align}
0
={}&
\left[
-\frac{n-1}{2}E_w + 2E_f + 3nE_{\rm PV}
\right]\,d\ln\rho
\nonumber\\
&
+
\left[
E_w+4E_f+9E_{\rm PV}
\right]\,d\ln a.
\label{eq:app-differentiated-equilibrium}
\end{align}
Solving for the breathing slope gives
\begin{equation}
\boxed{
\frac{d\ln a}{d\ln\rho}
=
-
\frac{-\frac{n-1}{2} + 2x + n(1+2x)}{4+10x}.
}
\label{eq:app-dlna-general}
\end{equation}
For the present closure, \(n=5\) and \(x=2/11\), so
\begin{equation}
\boxed{
\frac{d\ln a}{d\ln\rho}
=
-\frac{57}{64}
\approx -0.890625.
}
\label{eq:app-dlna-57-64}
\end{equation}
This number is not needed to assemble the conservative two-body \(1\)PN Lagrangian, but
it is an explicit prediction of the adiabatic breathing model and therefore belongs in the
appendix-level closure ledger.

\subsection{Consequences for the 1PN ledger}
\label{app:adiabatic-ledger}

The role of the present appendix in the main paper is now immediate. With
\begin{equation}
\krho=1,
\qquad
\kadd=\frac12,
\qquad
\kPV=\frac32,
\end{equation}
the self-sector/precession ledger is
\begin{equation}
\boxed{
\betaPN
=
\krho+\kadd+\kPV
=
1+\frac12+\frac32
=
3.
}
\label{eq:app-beta-three}
\end{equation}
This value matches the independently derived self-sector coefficient and reproduces the
standard perihelion coefficient in the test-mass orbit reduction.

\subsection{Scale degeneracy and scope}
\label{app:adiabatic-scope}

The adiabatic closure fixes the dimensionless partition \(x\), and therefore fixes
\(\kPV\), but it does \emph{not} fix the absolute equilibrium scale \(a(\rho_0)\). The
overall constants
\begin{equation}
\mathcal C_w,\qquad \mathcal C_f,\qquad \mathcal C_{\rm PV},\qquad \Lambda
\label{eq:app-scale-constants}
\end{equation}
enter only multiplicatively in \(F\) and therefore determine the absolute normalization of
the equilibrium radius, not the dimensionless response slope derived above. Fixing those
constants requires an additional normalization convention or a more microscopic
identification of the wave, flow, and PV energies.

It is also important to keep the closure's scope explicit. The result
\(\kPV=3/2\) is the \emph{adiabatic} \(\omega\to 0\) limit of the response problem. If the
geometry does not re-equilibrate instantaneously, one should instead solve the dynamical
compliance problem and promote
\begin{equation}
\kPV
\;\to\;
\kPV(\omega,\text{protocol}),
\label{eq:app-kPV-omega}
\end{equation}
with the present result understood as the low-frequency limit of that more general
response function.

\claimclass{\Protocol. The generic adiabatic-elimination formulas
\eqref{eq:app-dq-adiabatic}--\eqref{eq:app-Heff-adiabatic} are exact consequences of the
quadratic response expansion once the control parameter and geometry coordinates are
declared, but the explicit results
\(\kPV=3/2\), \(x=2/11\), \(E_w:E_f:E_{\rm PV}=11:2:5\), and
\(d\ln a/d\ln\rho=-57/64\) belong specifically to the one-DOF adiabatic throat protocol
adopted in this paper.}

\end{document}
